% HAFF Paper D: observable-algebra selection, modular candidates and physical sources
\documentclass[12pt,a4paper]{article}
\usepackage{amsmath,amssymb,amsfonts,amsthm}
\usepackage{physics}
\usepackage{hyperref}
\hypersetup{pdftitle={Selection, Maintenance and Gravitational Response: Modular Constraints and a Physical Return},pdfauthor={Sidong Liu}}
\usepackage{geometry}
\usepackage{array}
\newcolumntype{P}[1]{>{\raggedright\arraybackslash}p{#1}}
\usepackage{booktabs}
\geometry{margin=1in}

\newtheorem{theorem}{Theorem}[section]
\newtheorem{proposition}[theorem]{Proposition}
\newtheorem{definition}[theorem]{Definition}
\newtheorem{remark}[theorem]{Remark}
\newtheorem{conjecture}[theorem]{Conjecture}

\title{Selection, Maintenance and Gravitational Response:\\
Modular Constraints and a Physical Return}

\author{
  Sidong Liu, PhD \\
  iBioStratix Ltd \\
  \texttt{sidongliu@hotmail.com}
}

\date{October 2026}

\begin{document}
\emergencystretch=2em
\raggedbottom
\hbadness=5000
\vbadness=5000

\maketitle


\begin{abstract}
Physical selection changes a gravitational source when it changes an actual
coupling, preparation or maintenance process. We develop this route alongside
a modular route that constrains which geometric laws can agree with selected
algebraic data. The two routes have different inputs: modular stability does
not select a unique algebra, and a supplied gravitational action does not
explain the origin of that law.

Calibrated correction, projection selection and learned protection produce
stress through the interactions that change access. We compute its
Jackiw--Teitelboim (JT) response and construct a finite semiclassical
mean-action return in which
acquisition changes a material source and later sensing changes the next
force. Quantum history, old references, support, synchronization, clock
work and recoil are retained. Joint perturbation estimates allow bounded controller coupling to
an unbounded common Hamiltonian on domain-preserving strong solutions. A separate full-mode two-field model
matches its stress to a work-corrected source at rate $O(I^{-1})$ as
clock inertia $I$ grows; the energy correction leaves a nonzero pressure
response. For a real-quartet carrier, an exact nonlinear normal-order
conversion couples its actual full-mode quantum state, physical work
clock and material director to constrained JT evolution. The resulting
bulk solution has a prepared positive causal domain.

Maintenance is priced by the entropy removed from the full retained state.
For a thermal chiral exhaust of capacity $c_{\rm ex}$, its long-window heat
rate obeys $\dot Q\ge k_BT\dot s+3\hbar\dot s^2/(\pi c_{\rm ex})$.
The complete output's relative entropy and entropy slack measure its
excess above this floor. Direct export of the current purifier retains
cross-species redundancy, so accurate renewal alone cannot make that
protocol sharp.
Matched horizon transport gives $\overline{\dot S_h}\ge
\dot s+6\dot s^2/c_{\rm ex}$ in the thermal rapidity clock.
The actual thermal-source construction includes the running clock's
nonextensive reference stress; physical realizations use finite operating
windows. Pointer-aligned dephasing gives different complete renewal debits
for the two selected axes. This identifies a candidate mechanism for a gravitational
maintenance reading to favor the stable pointer algebra. Sustainable
renewal, sharpness of the complete adaptive output and persistent causal
program contrast form a joint conjecture, with drive routing, moving-clock
conversion and support work charged. A common source action and the
origin of geometric selection remain further targets.
\end{abstract}

\section{The problem: what makes selection gravitational?}
\label{sec:intro}

An apparatus changes which distinctions it can preserve and which operations
it can execute. That change may reorganize a reconstructed geometry; it may
also change the state of matter. A gravitational claim must say which of
these changes it describes. Forgetting a record after an operation changes
what can be inferred, while the emitted stress persists. Changing the
interaction can change both access and the physical source.

The question of this paper is whether these two descriptions can meet:
can the structures selected by physical operations supply geometric data
whose consistency requires a gravitational law, and can the same operations
produce sources whose response changes subsequent selection? HAFF-A, B and E
provide the distinction between channels, stable content and actual records
\cite{Liu2026PaperA,Liu2026PaperB}; Foundation~I provides a hull relative to a
readable seed~\cite{FoundationI}. Here the object to explain is the relation
between these physical selections and gravitational dynamics.

One route starts with modular data. A supplied standard half-sided modular
family generates translations. With a holographic dictionary and the required
ball first laws, its variations match the linearized Einstein equation.
Consciousness~II supplies a different, same-spacetime conditional consistency
law for a boosted ball family~\cite{PaperCII}. Neither route derives its
geometric dictionary or consistency condition from modular stability alone.
The SLD metric of distinguishability is positive semidefinite; it is neither
the BKM Hessian of relative entropy nor a Lorentzian spacetime metric.

The other route starts with actual interactions. We derive sources from
calibrated correction, selected projections and learned protection, then
compute their response in a specified JT theory. The JT metric has fixed
curvature in the minimally coupled sector; its dilaton and invariant respond.
Direct dilaton-coupled support changes that sector explicitly. The resulting
finite return retains a quantum reader and its old reference while its acquired
record drives mechanics and new radiation. Its complete mean-action work and
recoil are included. This realizes a causal source mechanism without assuming
the proposed origin of the gravity law.

Sections~\ref{sec:source}--\ref{sec:gap} develop the derived sources,
finite physical return and maintenance price. The modular route follows
in Sections~\ref{sec:factorization}--\ref{sec:modular}; the final
conjectures ask how these routes could meet. The appendices retain
the instruments, profiles, clocks and apparatus proofs. In particular,
Table~\ref{tab:scope-ledger} states the objects and the interfaces between
them. It separates established results from the growing-window conjecture;
it is not a decomposition of one already completed apparatus.

\begin{table}[p]
\centering\small
\begin{tabular}{P{3.0cm}P{5.3cm}P{5.6cm}}
\toprule
Object and supplied inputs & Result & Remaining interface\\
\midrule
Modular family, dictionary and ball laws &
Conditional translations and linearized geometric consistency
(Sections~\ref{sec:factorization}--\ref{sec:modular}) &
Formation of the candidate family and its geometric dictionary.\\[.6ex]
Calibrated field operations and a JT action &
Derived outgoing source, retarded response and finite echo/contact
readout (Section~\ref{sec:source}; Appendix~\ref{app:boundary-return}) &
Total apparatus stress and physical references; the action is an input.\\[.6ex]
Quantum core with classical mechanics and mean-action source &
Finite retained-history return and strong-domain perturbation bounds
(Section~\ref{sec:physical-return}; Foundation~III) &
Fixed-window mean-action result; protected content, old references and
initial work remain charged. No quantum stress tensor for the scalar
fields or support is supplied.\\[.6ex]
Spatial quadratic fields and real-quartet bulk &
Work-corrected full-mode stress matching, a finite spatial return,
and a separately constrained quartet interior &
Joining their receivers, moving endpoints, preparation and thermal
ports in one material model.\\[.6ex]
Complete retained entropy and paid finite carrier &
Long-window heat floor and neutral-sector sorting of a resettable
source (Section~\ref{sec:gap}; Proposition~\ref{prop:pointer-sort}) &
Actual adaptive source concentration, protected-content compatibility
and local delayed readout.\\[.6ex]
One reused body over growing runs &
Conjecture~\ref{conj:gravitating} jointly requires complete-output
sharpness, normalized window contrast and causal transition gain &
Sustained inlet, support, sensing and resource account in that same
dynamics; none is supplied by sorting alone.\\
\bottomrule
\end{tabular}
\caption{Objects, results and interfaces. Finite constructions retain their
stated inputs; the joint growing-window target is a conjecture.}
\label{tab:scope-ledger}
\end{table}

\section{From a Selected Interface to a Physical Source}
\label{sec:source}

A calibrated interface acts through a field coupling whose outgoing stress
sources the supplied JT theory. This gives a physical selection--source map;
the origin of the gravitational law in Conjecture~\ref{conj:main} is a
separate question.

\subsection{What records can and cannot do to a source}
\label{sec:records}

Let a classical record $R\sim\mathrm{Bernoulli}(p)$ control whether a physical operation prepares the matter state $\psi$ or leaves the reference $\omega$. Without reading $R$, the matter state is $\rho_p=(1-p)\omega+p\psi$ and its stress is the mixture of the two stresses. Deleting or ignoring $R$ after the preparation does not act on the matter, so $\rho_p$ and its source are unchanged. Changing the preparation probability $p$ changes the source. The information side has a matching statement. For the joint state $\rho_{RM}=(1-p)|0\rangle\langle0|\otimes\omega+p|1\rangle\langle1|\otimes\psi$ and the reference $\pi_R\otimes\omega$, with $\pi_R$ the marginal of $R$, additivity of relative entropy on direct sums and data processing give
\begin{equation}
D(\rho_{RM}\Vert\pi_R\otimes\omega)=p\,D(\psi\Vert\omega),\qquad D(\rho_p\Vert\omega)\le p\,D(\psi\Vert\omega)
\end{equation}
\cite{OhyaPetz1993}. Losing the record lowers the distinguishability available to a reader; it does not remove the matter source. A gravitational consequence of selection must therefore pass through a preparation or interaction, and the conditional geometry of a branch, the average of branch geometries and the geometry of the averaged source are three different objects.

\subsection{The source of a calibrated feedback actuator}
\label{sec:actuator-source}

The companion paper~\cite{Liu2026QRAIFB} realizes a corrective flip of a degenerate bit as a field coupling. A right-moving free massless scalar with null coordinate $u$, current $j=\partial_u\phi$ and stress $T_{uu}=:j^2:$ is coupled through a real smooth switch $g$. We require a finite nonzero one-particle norm
$\lambda^2=\int_0^\infty\omega|\hat g(\omega)|^2d\omega/(4\pi)$ and
$\int(1+u^2)|g'(u)|^2du<\infty$, where
$\hat g(\omega)=\int g(u)e^{-i\omega u}du$ in the $L^2$ sense when needed.
For example, $g\in H^1(\mathbb R)$ together with the weighted derivative condition
is sufficient. Smooth compact switches and the full-line packet below both belong
to this class. The smeared current defines a Weyl unitary,
\begin{equation}
U=\exp\Bigl(-iX\!\int g(u)j(u)\,du\Bigr)=\exp[-i\lambda X(b+b^\dagger)],
\end{equation}
where $b$ is a normalized field mode with mean frequency $\bar\omega$ fixed by $g$, and $\lambda^2\bar\omega=\tfrac14\int g'^2du$. With the drive mode in the coherent state $|\sqrt N\rangle$ and $2\lambda\sqrt N=\pi/2$, the coupling flips the bit with success probability $s_N=(1+e^{-\pi^2/(8N)})/2$. From the current commutator $[j(u),j(v)]=\tfrac i2\delta'(u-v)$ one has $U^\dagger j(u)U=j(u)+\tfrac12Xg'(u)$ and
\begin{equation}
U^\dagger T_{uu}(u)U=T_{uu}(u)+g'(u)Xj(u)+\tfrac14g'(u)^2 .
\label{eq:stress-id}
\end{equation}
The cross term has zero expectation if the bit and the fresh field input are initially independent and $\operatorname{tr}(\rho X)=0$, or if their joint state has no coherence in the bit's $Z$ basis; a $Z$-diagonal bit marginal alone does not suffice~\cite{Liu2026QRAIFB}. Writing $j=\psi(u)b+\psi^*(u)b^\dagger+j_\perp$, the incoming coherent beam has stress $T_{\rm in}=4N[\operatorname{Re}\psi(u)]^2$ and each activation adds
\begin{equation}
F_g(u)=\tfrac14g'(u)^2=4\lambda^2[\operatorname{Im}\psi(u)]^2 .
\label{eq:Fg}
\end{equation}
This added stress is a property of the coupling. Its integral is the switching work $\lambda^2\bar\omega$ supplied by the external switching agent.

The feedback task is that of the Imprint Chain~\cite{Liu2026Chain} with the flip realized by this coupling. A hidden instrument label is estimated from calibration records with mismatch probability $q$; each task step reads the bit with reliability $\nu$, accepts feedback with probability $h$, switches $g$ on when accepted and the reading is negative, and lets a bath with contraction $c$ act. Each opportunity receives its own scheduled incoming beam, independent of the bit and earlier records. In the companion paper the stationary loss $L_\infty(q)$ and the stationary coupling rate satisfy~\cite{Liu2026QRAIFB}
\begin{equation}
p(q)=h\nu L_\infty(q)+\tfrac{h(1-\nu)}2 .
\label{eq:rate}
\end{equation}

\begin{proposition}[Source--loss relation]
\label{prop:source-loss}
At fixed drive, switch and task parameters, the stationary mean source per opportunity outside the apparatus is $\langle T_{uu}\rangle_q=T_{\rm in}+p(q)F_g$. For two calibrations,
\begin{equation}
\langle T_{uu}(u)\rangle_q-\langle T_{uu}(u)\rangle_{q'}=h\nu\,[L_\infty(q)-L_\infty(q')]\,F_g(u).
\label{eq:source-loss}
\end{equation}
\end{proposition}

\begin{proof}
The incoming beam is present at every opportunity. A coupling occurs with stationary probability $p(q)$ and, under the joint initial condition, adds $F_g$ in expectation by Eq.~\eqref{eq:stress-id}. Equation~\eqref{eq:rate} gives the difference.
\end{proof}

Better calibration lowers the task loss and, through it, the source. The coefficient comes from the same interaction that performs the correction; no separate rule maps a calibration error to an emitted pulse. The relation does not assign stress to knowledge. Deleting the calibration record after the commands have been executed leaves the outgoing source unchanged, in line with Section~\ref{sec:records}. The relation is a stationary average per opportunity; a single trajectory, with its fixed hidden match label, does not follow the population mean step by step.

\subsection{A source set by learned quantum protection}
\label{sec:learned-protection-source}

The activation-rate mechanism is not the only way that retained records can
change a source. Brain Anchor~\cite{BrainAnchor} constructs two degenerate
qubits with a programmable coupling
\begin{equation}
 A_a=Z_1+\alpha aZ_2,\qquad
 U_a=e^{-iA_aQ_f},\quad
 Q_f=\int f(u)j(u)du=\sqrt{k}(b+b^\dagger).
 \label{eq:protection-coupling}
\end{equation}
The real smooth compact switch $f$ and positive gain interval are supplied.
The gain $\alpha$ is unknown, and each encounter receives a fresh independent
vacuum field. On $\operatorname{span}\{|01\rangle,|10\rangle\}$ the
coupling is $(1-\alpha a)Z_L$. Calibration records from the commands $\pm h$
have means $r_\pm=e^{-2k(1\mp\alpha h)^2}$, so
$\log(r_+/r_-)=8k\alpha h$. The finite-record certificate in that paper
bounds $|1-\alpha a|$, including command rounding and known drift.
It does not promise an exact fixed algebra after finite training.

\begin{proposition}[One residual controls protection and source]
\label{prop:protection-source}
Fix the acquired command $a$, write $m=1-\alpha a$ and
$F_f=f'^2/4$, and keep the initial field independent of the logical
body and retained ancillas. For every logical input, one encounter has
added mean stress $m^2F_f$ and switching work $m^2E_f$, with
$E_f=\int F_fdu$. After $M$ independent identical encounters, with the
output ports discarded and no interleaved gates or active recovery,
its passive protection error is
\begin{equation}
 \epsilon_M=\tfrac12(1-e^{-2Mk m^2}).
 \label{eq:protection-source-risk}
\end{equation}
For a single pulse with stress moments $E_j=\int u^jF_fdu$ and the
initial data of Proposition~\ref{prop:jt}, the late invariant is
\begin{equation}
 \mathcal C(m)=H_0^2+2\kappa H_0m^2E_1
                 +\kappa^2m^4(E_1^2-E_0E_2).
 \label{eq:protection-source-invariant}
\end{equation}
Thus retained history can change the physical protection and the
gravitational response at fixed pulse shape and activation schedule.
\end{proposition}
\begin{proof}
The current shift is $j\mapsto j+A_af'/2$. The centered independent
field input makes the stress cross term vanish for every logical state,
and $A_a^2=m^2I_L$ on the code. Its characteristic function multiplies
the logical off-diagonal by $e^{-2km^2}$ per encounter; a dephasing
channel's half-diamond error is its phase-flip probability. Time ordering
may give a factor $e^{i\vartheta_f A_a^2}$, which is scalar on the code
but not on the whole two-qubit register. Finally substitute $m^2F_f$
in the fixed-initial-data source solution below.
\end{proof}

This link does not require a rule that emits a pulse only after a wrong
answer. The same interaction produces both the residual error and its
source, even when the logical body is entangled with retained memory.
At $m=0$ the body decouples in the code; a noiseless stored qubit need
not continuously export entropy. The finite-rate entropy-export bound
of Section~\ref{sec:gap} applies to actual renewal, not to every protected
state. Calibration itself adds field energy
$2N(1+\alpha^2h^2)E_f$ for $N$ independent tests at each sign.
Preparations, measurement, the controller and switching mechanism add
further charges. The source suppression during operation alone is not
a total-energy saving.

Brain Anchor also treats an unknown centered Gaussian variance $K_\omega$. Logical records at $0,\pm h$ give
$K_\omega=-\log r(0)/2$ and
$\alpha=[\log r(h)-\log r(-h)]/(8hK_\omega)$; two settings suffice when the signs of $1\mp\alpha h$ are known. Its finite-record mismatch certificate needs no supplied temperature label. The coherence is now $e^{-2K_\omega m^2}$, but the coupling-added source remains $m^2F_f$ on centered independent input.

Incoming stress must still be added. A $u$-clock chiral KMS input has background $\pi/(12\beta^2)$ in natural units. Within a supplied bounded one-temperature family, strictly decreasing $K_\beta$ permits a finite-error background estimate; a general Gaussian variance does not determine incoming energy. At common preparations and initial data, the linear dilaton equation compares coupling-source differences. Equation~\eqref{eq:protection-source-invariant} cannot be used for the full source after dropping its thermal background, nor is $u$ thermality the rapidity-clock reference of Section~\ref{sec:gap}.

The same command can also operate on one reused vacuum field. For nonnegative compact switches $f_j$ on disjoint intervals, with no noncommuting logical gates, the smeared currents commute but have negative correlations. Brain Anchor gives exact coherence $e^{-2m^2k_M}$ with
$k_M=\langle(\sum_jQ_{f_j})^2\rangle\le Mk$, and added source $m^2\sum_jf_j'^2/4$ of energy $Mm^2E_f$. Its JT response uses the actual pulse-train moments and initial coefficients in a valid operating domain. No independent-port approximation, thermal input or infinite-lifetime worldline is inferred.

There is a finite causal pulse in this construction. For example, set
$f(u)=A\exp[-1/(1-(u-3)^2)]$ on $2<u<4$, zero outside, and normalize
its one-particle norm to $k=.2$. Numerical quadrature gives
$(E_0,E_1,E_2)=(.4568522,1.3705565,4.3215400)$.
For $H_0=1$, $\kappa=.05$ and $\alpha=1.1$, the supplied command
$a=48/55$ has residual $m=.04$ and gives
$\mathcal C\simeq1.00021929$; the comparison $a=1$
has $m=-.1$ and $\mathcal C\simeq1.00137053$.
These are deterministic evaluations of the displayed pulse, with code
and outputs in \texttt{supplemental/pulse\_debit}, not statistics of
a calibration pilot. For a random learned command,
the invariant of the mean source uses $(\mathbb E m^2)^2$, whereas
the mean branch invariant uses $\mathbb E m^4$ in the quadratic term. The compact source ends before a later boundary echo is
scheduled, but the calibrated wall and timing instrument remain the
inputs of Section~\ref{sec:echo}.

Stress readback also has an exact blind direction. At fixed $a_0>0$,
the gains $\alpha_\pm=(1\pm d)/a_0$ give residuals $\mp d$ and therefore
the same protection channel and source. Field parity conjugates their
full code interactions while leaving the vacuum and stress invariant.
No parity-covariant stress instrument with invariant outcome labels can
distinguish them, nor can their mean JT responses at common initial data.
Their correcting commands $1/\alpha_\pm$ differ. Probes at $a_0\pm h$
break this ambiguity through the opposite signs of
$8k\alpha h(1-\alpha a_0)$ in their log contrast. A gravitational
response is therefore not a complete calibration reference for the
interface that produced it. Off-frame coupling probes, or a signed
field quadrature, are additional access, with their own sources and costs.
Full output-field parity and a conditional logical pulse can instead
restore the stored state without identifying the gain, as shown in
Ref.~\cite{BrainAnchor}. This recovery leaves the emitted source
unchanged. Active state correction and learning a match that suppresses
the actual source are therefore distinct operations.

The gain is a physical actuator parameter, not a freely changed
description. The selected finite code is not the geometric algebra of
Section~\ref{sec:modular}, and the supplied JT action still supplies the
gravitational response. A link to continuous modular selection requires
that further interface; it is not obtained by identifying a protected
$M_2$ sector with a half-sided modular inclusion.


\subsection{A physical seed changes the source}
\label{sec:projection-source}

The rate-controlled model changes how often a fixed interaction occurs. The coupled
observable can itself change, giving a direct relation between an exactly represented
algebra and a source. Keep the incoming field, switch and activation schedule fixed.
On a qubit put $P_\theta=(I+\cos\theta\,Z+\sin\theta\,X)/2$ and use
$U_\theta=\exp[-i\lambda P_\theta(b+b^\dagger)]$. The projection selects a pointer
algebra and its source population under this fixed schedule.

\begin{proposition}[Selected projection and outgoing stress]\label{prop:projection-source}
For $\lambda\ne0$, an initially independent qubit state $\rho$ and normal field state $\sigma$, the
exactly represented algebra of the induced qubit channel is $\{P_\theta\}'$. Let
$r_\theta=\operatorname{tr}(\rho P_\theta)$. The stress formula below requires its displayed current moments to exist. The outgoing field is
\[
 (1-r_\theta)\sigma+r_\theta e^{-i\lambda Q}\sigma e^{i\lambda Q},\qquad Q=b+b^\dagger,
\]
and its stress is
\begin{equation}
 T_\theta(u)=T_{\rm in}(u)+r_\theta
 \left[g'(u)\langle j(u)\rangle_\sigma+\tfrac14g'(u)^2\right].
\label{eq:projection-source}
\end{equation}
Thus two selected projections can change the source at fixed global input state and
fixed activation schedule.
\end{proposition}
\begin{proof}
The quadrature has non-atomic spectral distribution in a normal state, so the
multiplicative-domain theorem of Q-RAIF~C gives $\{P_\theta\}'$~\cite{QRC}.
Expanding $U_\theta$ in its two projection blocks and tracing the qubit gives the
field mixture; off-diagonal blocks have zero trace even when $\rho$ does not commute
with $P_\theta$. In the current identity replace the Pauli $X$ by $P_\theta$ and
use $P_\theta^2=P_\theta$. Initial independence gives~\eqref{eq:projection-source}.
\end{proof}

For $\rho=(I+0.8Z)/2$, the seeds $P_0$ and $P_{\pi/2}$ select different pointer
algebras, with $r_0=0.9$ and $r_{\pi/2}=0.5$. In the worked coherent packet,
the cross profile is $8\lambda\sqrt N\operatorname{Re}\psi\operatorname{Im}\psi$.
Its first three moments vanish, so the invariant increment is~\eqref{eq:dC} with
$p$ replaced by $r_\theta$. With the table's $N=16$ and background it is
$2.0807581\times10^{-5}$ versus $1.1559770\times10^{-5}$. The full source also
contains the cross profile, which enters the local dilaton even though its first
three moments do not enter the final invariant.

Here the source weight is derived from the state's population in the selected
projection; it is not assigned to an external mismatch label. Changing the projection
means changing a physical interaction, not renaming an unchanged algebra. The two
algebras are abstractly isomorphic but are differently embedded relative to the fixed
state and coupling. This supplies a seed--interface--source link. It does not identify
that finite-dimensional pointer algebra with a Type~III geometric candidate or explain
why the coupling selects this seed rather than another. Those are the remaining
selection and geometry bridges.


\subsection{Response of a given Jackiw--Teitelboim model}
\label{sec:jt}

Take the JT action~\cite{AlmheiriPolchinski2015,Moitra2020}
\begin{equation}
I=\frac1{16\pi G}\int\sqrt{-g}\,\Phi(R+2/L^2)+I_{\rm matter},\qquad
ds^2=-\frac{4L^2\,du\,dv}{(1+uv)^2},
\end{equation}
with the usual boundary terms, $\kappa=8\pi G$, $D=1+uv$ and the vacuum dilaton $\bar\Phi=H_0(1-uv)/D$, $H_0>0$. The field equations are $\nabla_\mu\nabla_\nu\Phi-g_{\mu\nu}\Box\Phi+g_{\mu\nu}\Phi/L^2=-\kappa T_{\mu\nu}$, and in these coordinates
\begin{equation}
\Phi_{uu}+\tfrac{2v}D\Phi_u=-\kappa T_{uu},\qquad
\Phi_{vv}+\tfrac{2u}D\Phi_v=-\kappa T_{vv},\qquad
\Phi_{uv}+\tfrac{2}{D^2}\Phi=\kappa T_{uv} .
\label{eq:jt}
\end{equation}
The matter is a conformal theory with matched left and right central charges, and the state-independent anomaly is removed by subtracting the same background. The left-moving input is unchanged, and the state-dependent source is right-moving, $T_{uu}=F(u)$ with $T_{uv}=T_{vv}=0$.

\begin{proposition}[Response with fixed initial data]
\label{prop:jt}
Let $F$ be integrable with $\int|u|^2|F|du<\infty$. On a connected coordinate domain with $D\ne0$, write the perturbation as $\Phi-\bar\Phi=(A(u)+vB(u))/D$ and impose the initial coefficient data $A(u),B(u)\to0$ as $u\to-\infty$. For a source bounded below in support this is equivalent to zero perturbation before the source. The unique solution with these data is
\begin{equation}
\Phi=\bar\Phi+\frac{A(u)+vB(u)}D,\qquad
A=-\kappa\int_{-\infty}^{u}(u-z)F(z)dz,\qquad
B=-\kappa\int_{-\infty}^{u}z(u-z)F(z)dz .
\label{eq:retarded}
\end{equation}
The perturbation is linear in $F$. With $E_j=\int u^jF\,du$, after the source (or in the asymptotic limit for sources with tails) the invariant $\mathcal C=\Phi^2-L^2(\nabla\Phi)^2$ takes the vacuum value
\begin{equation}
\mathcal C=H_0^2+2\kappa H_0E_1+\kappa^2(E_1^2-E_0E_2),
\label{eq:invariant}
\end{equation}
and inside the support $\partial_u\mathcal C=\kappa F\,(2H_0u+uA-B)$.
\end{proposition}

\begin{proof}
Write $\psi=D\,(\Phi-\bar\Phi)$. The $vv$ equation of~\eqref{eq:jt} reads $\psi_{vv}=0$, so $\psi=A(u)+vB(u)$, and the other two equations reduce to $A''=-\kappa F$, $B''=uA''$ and $B'=uA'-A$, which Eq.~\eqref{eq:retarded} satisfies. The absolute first and second moments make both retarded coefficients tend to zero at $-\infty$. Two solutions differ by $\alpha+\beta u+\gamma v-\alpha uv$, whose coefficients are $A_h=\alpha+\beta u$ and $B_h=\gamma-\alpha u$. Their two zero limits force $\alpha=\beta=\gamma=0$. Every vacuum solution has numerator $\alpha+\beta u+\gamma v-\alpha uv$ and invariant $\alpha^2+\beta\gamma$; after the source $\alpha=H_0+\kappa E_1$, $\beta=-\kappa E_0$, $\gamma=\kappa E_2$. The balance law follows by differentiation. 
\end{proof}

The coefficient condition cannot be replaced by $\Phi-\bar\Phi\to0$ at each fixed $v$. The homogeneous solution $\eta v/(1+uv)$ has that limit for every fixed $v$ along nonsingular points. It even tends to zero uniformly on the subdomain $u<0$, $v\le0$, since its absolute value is at most $|\eta|/|u|$; no uniform limit near $D=0$ is implied. Adding it leaves the source unchanged but changes the late invariant by $-\kappa E_0\eta$. It is excluded by $B(-\infty)=0$. Source comparisons below hold $H_0$, the coordinate matching and both coefficient boundary data fixed.

The JT curvature stays $-2/L^2$; the dilaton and the invariant respond. A dilaton difference is not by itself a higher-dimensional curvature or a mass, and comparing boundary proper times needs a boundary normalization.

The matching condition matters. Set $A_F,B_F\to0$ at $+\infty$, or exactly to zero after a compact source, instead of the initial data of Proposition~\ref{prop:jt}. The final-state-matched solution differs from the retarded solution by
\[
A_F=A_R+\kappa(uE_0-E_1),\qquad
B_F=B_R+\kappa(uE_1-E_2).
\]
For the chiral diffeomorphism states of Consciousness~I~\cite{PaperCI}, with a Cut at $a$ to the left of the support, the relative entropy is $D_a=2\pi\int(u-a)F(u)\,du$. The final-state-matched solution gives $\delta\Phi(a,0)/(4G)=-D_a$, while the initial-state-matched solution gives $\delta\Phi(a,0)=0$ at a Cut before the pulse. The first identity compares an earlier section with a matched final state; it does not say that a pulse removes area before it arrives. The dilaton--modular-energy relation itself is known~\cite{Callebaut2019}; what is used here is its type-correct pairing with a specified source and matching condition.

The same equations show that a null-only record of the source is incomplete. A genuine vacuum-density shift $T_{\mu\nu}=-\delta\rho\,g_{\mu\nu}$ has zero null components but gives the particular solution $\delta\Phi=\kappa L^2\delta\rho$. A probe that reads only null energy cannot certify that a source is absent, which is the source-completeness condition relevant to the vacuum curvature test of Ref.~\cite{QRNV}.


\section{Response that changes the next operation}
\label{sec:physical-return}

A source-to-geometry calculation becomes a return only when the geometry
changes an actual record and that record changes a later interaction.
The readout permission matters. In the minimally coupled JT sector,
instruments on the same matter metric with the same inputs have identical
adaptive record laws even when the dilaton differs
(Proposition~\ref{prop:blind-readback}). A boundary tied to a fixed dilaton
value evades that obstruction. Its echoes reconstruct the three vacuum
coefficients, and a source-compensated inverse includes finite probe
backreaction (Propositions~\ref{prop:echo} and
\ref{prop:compensated-echo}). An exact current contact instead obtains an
invariant-dependent bit, including the contact's own source
(Propositions~\ref{prop:driven-boundary-sensor} and
\ref{prop:weak-boundary-sensor}). The bit selects the next physical
projection. Its same-field conditional covariance remains part of the
history; detector replacement does not reset the field.

These are finite physical records, not permission to read an exact metric.
The worked current contact has total variation only
$2.3650\times10^{-4}$ between its two initial geometries.
Its fixed-cutoff energy budget also forbids indefinitely many contacts of
unchanged positive source charge. Appendix~\ref{app:boundary-return}
contains the timing, stress and accumulated-source bounds.

\subsection{A finite prefix is a sufficient response object}

The source need not have finite second moment over all time to determine a
finite episode. This matters for actual material radiation, whose tails
are not inferred from a finite smooth-source certificate.

\begin{proposition}[Finite-prefix response and its source tolerance]
\label{prop:finite-prefix-response}
Fix $u_*<U$ and matched data $(a_*,b_*,\gamma_*)$ at $u_*$ in the
traceless single-chiral JT sector on a common patch with $1+uv\ne0$.
For $F\in L^1([u_*,U])$, define
\[
 a(u)=a_*+\kappa\int_{u_*}^{u}zF(z)dz,\quad
 b(u)=b_*-\kappa\int_{u_*}^{u}F(z)dz,\quad
 \gamma(u)=\gamma_*+\kappa\int_{u_*}^{u}z^2F(z)dz.
\]
Then $\Phi=(a+bu+\gamma v-auv)/(1+uv)$ solves the JT equations
on this window with those data, and
$\mathcal C(u)=a(u)^2+b(u)\gamma(u)$. No future-source assumption
enters either reading. If two sources have
$\int_{u_*}^{U}(1+z^2)|\delta F|dz\le D$, then
\begin{align}
 |\delta a|,|\delta b|,|\delta\gamma|&\le\kappa D,\nonumber\\
 |\delta\Phi|&\le
 \frac{\kappa D(1+|uv|+|u|+|v|)}{|1+uv|},\label{eq:prefix-tolerance}\\
 |\delta\mathcal C|&\le
 \kappa D(2|a|+|b|+|\gamma|)+2\kappa^2D^2.\nonumber
\end{align}
\end{proposition}
\begin{proof}
The coefficient derivatives are $a'=\kappa uF$, $b'=-\kappa F$
and $\gamma'=\kappa u^2F$. They cancel from the first derivatives
of the numerator, giving the field equations in their weak sense.
Direct substitution in $\Phi^2-L^2(\nabla\Phi)^2$ gives
$a^2+b\gamma$. The three integrands are bounded in absolute value
by $1+z^2$. Expanding the numerator and invariant gives the estimates.
\end{proof}

For the complete quantum material, Q-RAIF~C bounds smooth interacting
stress tests on one positive Hamiltonian with all old references retained
\cite{QRC}. Its cross and creation contributions are essential.
The bound is a test-functional certificate, whereas
Eq.~\eqref{eq:prefix-tolerance} uses a weighted $L^1$ certificate.
They are not interchangeable. With a smooth preparation cutoff,
smearing a finite-resolution dilaton reading gives smooth retarded kernels.
Each such kernel can be used as a source test with its own Q-RAIF~C
form constant. Recovering a sharp prefix additionally requires a bound
on unresolved source structure. This gives a direct finite-window
interface without claiming an unproved all-time tail estimate or
discarding controller and support sources.

\subsection{The same quantum reader, support and source}

A reflecting boundary alone does not supply its support. The material
branch uses a positive dilaton-dependent mass and varies the worldline.
It contains the finite reader, history clock, blank and old purification
in its mass. Classical scalar fields couple to an angle $\theta$; storage
$y$ has positive kinetic and spring energies. Their emitted powers are
$P_\pm=g_\pm^2\dot\theta^2$, derived by characteristics, and their
reaction is included in the mechanical equation. The two exterior
invariants evolve with the proper worldline and obey the mass junction.
Appendix~\ref{app:material-return} gives this complete mean-action
construction and its admissible positive domain.

The quantum core is
\[
 H_c=H_{\rm hist}+E_{\rm off}I+
 \lambda(\Phi-\phi_0)Z_P+
 \frac{\chi_R}{2}P_R(y-d)^2.
\]
Its first active clock edge writes an actual $Y_P$ outcome; the next
swaps the probe with a paid blank. The record spring reacts while
acquisition is in progress. Variation gives
$F_y=\chi_R(y-d)\langle P_R\rangle$ and
$F_\Phi=\lambda\langle Z_P\rangle$, with
$\dot E_c=F_y\dot y+F_\Phi\dot\Phi$.
Consequently $M=m_{\rm rest}+U+e+E_c$ obeys
$\dot M=(U'+F_\Phi)\dot\Phi-P_+-P_-$.
These are quantum expectations in a retained pure-state calculation,
not a sampled record fed to an independent source generator.

Proposition~\ref{prop:record-mediated-source} proves a nonzero source
return through its shortest-path expansion. Either $\lambda=0$ or
$\chi_R=0$ removes the proper-power difference between its two initial
geometries. Both interactions present produce a nonzero coefficient at
order $t^{2r+7}$, where $r$ is the acquisition edge. The finite
33-state example retains its old probe reference and all work balances.
This is a mean-worldline realization; quantum support fluctuations
are not removed by taking the quantum expectation of its force.

The reader must also remain accurate when the response is allowed to
move, rather than being supplied beforehand. The following estimate
controls this coupled problem.

\begin{proposition}[Joint tolerance of an acquired-record source return]
\label{prop:coupled-return-tolerance}
Write the mean-action equations as
$\dot z=\mathcal F(z,\Psi)$ and $i\dot\Psi=H_c(z)\Psi$,
retaining the old reference in $\Psi$.
Allow $H_c(z)=H_*+V_c(z)$, where $H_*$ is a common semibounded self-adjoint
internal Hamiltonian and $V_c$ and its first two coordinate derivatives
are bounded on the stated domain. For an unbounded core, require strong
initial data in $D(H_*)$ and propagators preserving this common domain
along the stated coordinate paths. This includes the finite core above
and regular Gaussian internal data with bounded controller coupling;
it does not include an unbounded coordinate-force operator merely
because its free field is Gaussian. Include the complete mass $M$
among the classical variables,
with its actual initial value and work equation
\[
 \dot M=(U'+\langle\partial_\Phi V_c\rangle)\dot\Phi-P_+-P_-.
\]
Let exact and approximate solutions, with the same core and
reference space, lie in a compact subset of the positive domain of
Proposition~\ref{prop:mechanical-source-window}, with
$|y|\le Y$ and $|\dot\theta|\le V$.
Use the maximum norm for classical variables nondimensionalized in
fixed material and clock units. Suppose the
approximate equations have classical residual at most $r_z(t)$,
including their mass-work residual. Let $r_H(t)$ bound the Hermitian
Hamiltonian error applied to the approximate quantum vector; the
propagators and this residual must exist on their specified domains.
Choose constants on this domain such that
\[
 |\mathcal F(z,\Psi)-\mathcal F(\widehat z,\widehat\Psi)|
 \le L_z|z-\widehat z|+L_\Psi\|\Psi-\widehat\Psi\|.
\]
Here the same reference is compared, not just its reduced probe state.
For the displayed sensor and spring the Hamiltonian Lipschitz constant is
$L_H=|\lambda|+\chi_R(Y+|d|)$.
For any $\sigma>0$ put
$L=\max\{L_z+\sigma L_H,L_\Psi/\sigma\}$ and
$E=|z-\widehat z|+\sigma\|\Psi-\widehat\Psi\|$.
Then
\begin{equation}
 E(t)\le e^{Lt}E(0)+\int_0^t e^{L(t-s)}
                   [r_z(s)+\sigma r_H(s)]ds.
 \label{eq:coupled-return-error}
\end{equation}
Any terminal record probability differs by at most $2E/\sigma$,
and the total proper emitted power differs by at most $2\gamma VE$,
where $\gamma=g_+^2+g_-^2$.
These estimates cover the actual quantum-force reaction in the classical
response, including its contribution to mass and clock conversion.
\end{proposition}
\begin{proof}
For normalized states,
$|\langle A\rangle_\Psi-\langle A\rangle_{\widehat\Psi}|
\le2\|A\|\|\Psi-\widehat\Psi\|$.
In particular the $P_R$ and $Z_P$ expectations have that bound.
The force differences are bounded by
\[
 |\delta F_y|\le\chi_R|\delta y|
       +2\chi_R(Y+|d|)\|\delta\Psi\|,\quad
 |\delta F_\Phi|\le2|\lambda|\|\delta\Psi\|,
\]
The mechanical work $-F_y\dot y$ cancels its contribution
$F_y\dot y$ to the internal energy, leaving the stated mass equation.
Thus mass is compared through its initial value and integrated work,
not by bounding $\langle H_*\rangle$ with the state-vector error.
The mass, square roots $\beta_\pm$ and null-clock denominators
are smooth with bounded derivatives on the stated compact domain;
these formulas give finite $L_z,L_\Psi$ for the actual coupled
equations. All clock and material energy remains in the initial mass.
The quadratic spring and the probe term give the displayed $L_H$.
Duhamel comparison using the exact unitary propagator gives
\[
 \|\delta\Psi(t)\|\le\|\delta\Psi(0)\|
    +\int_0^t[L_H|\delta z(s)|+r_H(s)]ds.
\]
It contains no spurious growth from the norm of a common clock
Hamiltonian. Add the classical integral bound with weight $\sigma$
and apply Gronwall. Finally $P_++P_-=\gamma\dot\theta^2$ gives
the power estimate, and a record effect has norm at most one.
\end{proof}

For the strong-domain solutions just specified, the Schr\"odinger
equation gives
\[
 \frac{d}{dt}\langle H_*+V_c(z)\rangle
 =\sum_j\langle\partial_jV_c(z)\rangle\dot z_j.
\]
The common $H_*$ contributes no explicit coordinate work. In its
interaction picture only bounded perturbation differences enter
Duhamel's estimate, so the same proof applies without bounding
$\langle H_*\rangle$ by a vector-norm error. Extending to all
energy-form initial data would additionally require an energy-controlled
approximation preserving this work identity. Spatial field observables
still require the separate material form bounds of Q-RAIF~C.
This extension treats the core as internal material of the specified
thin worldline; it does not equate spatial field stress with a point
source. Smooth tests of the line stress
$M u^\mu u^\nu\delta_\Gamma$ are Lipschitz on the same compact tube,
so their errors are bounded by a constant times $\int E(s)ds$.
The constants and admissible tube may change when the initial clock
mass grows; no uniform-in-clock-size geometric limit follows.

The work equation is necessary even when the retained quantum states
are close. Q-RAIF~C gives a positive clock-driven oscillator whose
actual total energy is $mJ$, whereas its scheduled product-clock target
has energy $mJ+b^2/(2\omega)$ for every $m$, despite vanishing joint
state error~\cite{QRC}. A nearby rotated clock can match that scalar
energy, but does not thereby match local stress. For a scheduled
comparison with internal Hamiltonian jumps $h_{b-1}\to h_b$, subtract
the switching work
$W_b=\langle h_b-h_{b-1}\rangle_{t_b}$ from its clock ledger.
Together with the continuous coordinate work this gives the mass
equation used above. Initial mass mismatch and any missing work belong
in $E(0)$ and $r_z$; they cannot be dropped by declaring the clock
state error small. This ledger is not the energy expectation of the
unmodified product-clock target.

A source contrast remains nonzero whenever the two power-error bounds
sum to less than its ideal contrast at a chosen time. Thus finite
implementation errors can be admitted while preserving the causal
sensor-removal and feedback-removal tests. The domain condition is
verified by a strict ideal margin and an error tube smaller than that
margin; it cannot be certified from a reduced-state fidelity alone.
These are the fixed-window mean-action bounds in
Table~\ref{tab:scope-ledger}.

Foundation~III uses this action for a frame learned while the original
quantum content is held~\cite{FoundationIII}. In a finite positive
parameter family, the acquired record changes the apparatus's own
source; its geometric response enters an $X$-basis reading of the old
probe into the already spent cell. The first-order geometry signal is
generated after record acquisition, with no new probe or blank port.
The matched intervention keeps the complete
initial state, mass and geometric constraints fixed and removes only
the first record's mechanical use. A sensor-edge diagnostic separates
the geometric path from direct probe and cursor effects. This is one
finite semiclassical mean-action scalar/JT return. The completed-window
geometry-edge contrast, averaged over the stated plant-sign prior,
is about $-4.98\times10^{-9}$ in the example;
its small size is reported separately from the task advantage.
It is distinct from the spatial realization below
and from the sustainable sharp-exhaust conjecture.

\subsection{A complete spatial clock-work source}
\label{sec:spatial-clock-work}

The work correction can be a spatial tensor rather than a scalar debit.
Consider the local quadratic action
\begin{equation}
 S_{\chi u}=-\frac12\int\sqrt{-g}\left[
 I(\nabla\chi)^2+(\nabla u)^2+I\mu^2\chi^2+
                      (\omega u+\lambda\chi/\omega)^2\right]d^2x,
 \label{eq:spatial-clock-action}
\end{equation}
where $I,\mu,\omega>0$. In canonical variables $(\sqrt I\chi,u)$
its mass matrix is
\[
 \begin{pmatrix}\mu^2+\lambda^2/(I\omega^2)&\lambda/\sqrt I\\
                 \lambda/\sqrt I&\omega^2\end{pmatrix}.
\]
Its determinant is $\mu^2\omega^2>0$. Constant orthogonal
diagonalization gives two positive-mass Klein--Gordon fields with
their complete spatial mode towers. Choose the interacting static
vacuum and a finite-energy Weyl displacement with initial data
$\chi=u=\dot u=0$, $\dot\chi=\Omega f$,
$f\in C_c^\infty$. Its added clock energy in flat space is
$I\Omega^2\|f\|_2^2/2$. One fixed local renormalization prescription
is used for the interacting vacuum and all displaced states.

Put $L_a=\partial_t^2-\partial_x^2+a^2$. In flat space let
\begin{equation}
 L_\mu c=0,\qquad L_\omega u_0=-\lambda c,\qquad
 L_\mu\eta_0=-\lambda u_0-\lambda^2c/\omega^2,
 \label{eq:spatial-retarded-work}
\end{equation}
where $c$ has the prepared clock data and the other two fields have
zero data. Define
\begin{align}
 T^{(a)}_{\alpha\beta}[q]
  &=\partial_\alpha q\partial_\beta q-\tfrac12g_{\alpha\beta}
                         [(\nabla q)^2+a^2q^2],\\
 B_{\alpha\beta}[c,\eta]
  &=\partial_\alpha c\partial_\beta\eta+
    \partial_\alpha\eta\partial_\beta c
      -g_{\alpha\beta}(\nabla c\cdot\nabla\eta+\mu^2c\eta).
\end{align}
The work-corrected comparison source is
\begin{equation}
 T_{\rm cmp}=I T^{(\mu)}[c]+B[c,\eta_0]+T^{(\omega)}[u_0]
       -g(\lambda c u_0+\lambda^2c^2/(2\omega^2))+T_{{\rm vac},I}.
 \label{eq:spatial-comparison}
\end{equation}
The large positive clock source and the same interacting-vacuum
source are present on both sides of the comparison.

\begin{proposition}[Full-mode stress matching]
\label{prop:spatial-clock-matching}
For each fixed finite window and compact smooth tensor test $W$,
the actual coherent quantum state of \eqref{eq:spatial-clock-action}
satisfies
\begin{equation}
 \left|\int\sqrt{-g}\,W^{\alpha\beta}
       (\langle T_{\alpha\beta}\rangle-T_{{\rm cmp},\alpha\beta})\right|
       \le C_W/I.
 \label{eq:spatial-source-match}
\end{equation}
The constants are independent of $I$ above a finite threshold. In flat
space the comparison obeys the exact local Ward identity. The
source difference also tends to zero in every fixed $C^k$ norm on a
compact window. The same statements hold on a fixed static metric
using its finite-window Klein--Gordon energy constants.
\end{proposition}
\begin{proof}
The actual means obey
$L_\mu\chi_I=-\lambda u_I/I-\lambda^2\chi_I/(I\omega^2)$ and
$L_\omega u_I=-\lambda\chi_I$. For
$\eta_I=I(\chi_I-c)$ these imply the exact identities
\[
 L_\mu\eta_I=-\lambda u_I-\lambda^2\chi_I/\omega^2,
 \qquad u_I-u_0=-\lambda G_\omega\eta_I/I,
\]
where $G_\omega$ is the zero-data retarded operator. Use
$\|q\|_{s,a}^2=\|\dot q\|_{H^s}^2+\|q_x\|_{H^s}^2+a^2\|q\|_{H^s}^2$.
The energy inequality is
$\|G_a F(t)\|_{s,a}\le\int_0^t\|F(r)\|_{H^s}dr$.
Set $C=|\Omega|\|f\|_{H^s}$,
$a_T=|\lambda|T/\mu$, $b_T=|\lambda|T/\omega$,
$d_T=\lambda^2T/(\omega^2\mu)$ and $h_T=a_Tb_T+d_T$.
For $I\ge2h_T$ it gives
\begin{align}
 \sup\|\eta_I\|_{s,\mu}&\le2h_TC,&
 \sup\|u_I-u_0\|_{s,\omega}&\le2a_Th_TC/I,\\
 \sup\|\eta_I-\eta_0\|_{s,\mu}&\le2h_T^2C/I.
\end{align}
These are full spatial energy estimates, without a mode cutoff.
The exact quantum covariance equals the interacting-vacuum covariance;
only the mean is displaced. Metric variation therefore gives
\[
 \langle T\rangle=I T^{(\mu)}[c]+B[c,\eta_I]
   +I^{-1}T^{(\mu)}[\eta_I]+T^{(\omega)}[u_I]
   -g(\lambda\chi_Iu_I+\lambda^2\chi_I^2/(2\omega^2))+T_{{\rm vac},I}.
\]
Polarization of each quadratic tensor and the displayed energy
estimates prove \eqref{eq:spatial-source-match}. Smooth initial data,
higher Sobolev estimates and the wave equations give the $C^k$ result.
Finally
$\nabla^\alpha B_{\alpha\beta}
 =(\lambda u_0+\lambda^2c/\omega^2)\partial_\beta c$;
the material and contact terms have the opposite divergence.
\end{proof}

The correction has a measurable tensor component. Relative to the
common clock and vacuum terms in \eqref{eq:spatial-comparison}, the
flat-space preparation gives
\begin{equation}
 \delta T_{tt}=O(t^4),\qquad
 \delta T_{xx}=-\lambda^2\Omega^2 f^2t^2/\omega^2+O(t^4).
\end{equation}
The clock pays the order-$t^2$ energy transfer; pressure does not cancel.
This subtraction defines a comparison, not a second experiment with
the coupling removed and a different vacuum. For metric-only matter
on a fixed AdS patch, the matched sources drive the linear JT dilaton
equation with the same initial constraints and boundary data. Its
finite-window response inherits the source estimate. A positive
dilaton domain may require initial homogeneous data growing with $I$;
that resource and a detector's gain cannot be removed from the account.
This is the quadratic clock/material entry of
Table~\ref{tab:scope-ledger}; its receiver and thermal interfaces are separate.

At finite spatial bandwidth, these source and readback steps can be
placed in one mean-action evolution. Proposition~\ref{prop:finite-spatial-return}
includes an elastic bar, a physical synchronization field, two retained
quantum spins and the full oscillator Hilbert space of the chosen modes.
The two preparations have the same initial gravitational constraints.
Their later pressure contrast produces a dilaton contrast, which changes
the retained sensor and the next material force. The synchronization
field is varied, not supplied as a prescribed clock. This establishes the finite spatial return listed in
Table~\ref{tab:scope-ledger}.

The retained quartet admits a full-mode bulk step toward that construction.
Proposition~\ref{prop:covariant-quartet-bulk} evolves its physical clock,
material director and exact renormalized quantum source together with
constrained JT data. The normal-order conversion and causal-state
identification are nonlinear. The result supplies an interior patch for
specified global quantum data; moving endpoints, the other gate windows
and the finite-horizon renewal error must still be joined to it.

\section{Maintenance and its gravitational price}
\label{sec:gap}

\subsection{Entropy removal and the source it can constrain}
\label{sec:maintenance}

A fixed unitary correction can return its control port coherently. Removing entropy
injected by unobserved noise is a different task. To price it, the state whose entropy
is removed must include the controller's usable side information. The heat account
must also distinguish a thermal reservoir from a supply of pure memories.

HAFF~G supplies a record-driven program and executor in which selecting
a physical coupling changes its source under one fixed executor
\cite{Liu2026PaperG}. The dephased quartet below adds a specified full
renewal debit. It does not identify the two material realizations or
inherit a thermal controller from that program loop.

\begin{proposition}[Maintenance entropy account]\label{prop:entropyaccount}
Let $V=SC$ include the maintained system and every retained controller and
side-information register. A correction acts unitarily on finite-dimensional
$VAB$, initially in $\rho_V\otimes\alpha_A\otimes\tau_B$, where
$\tau_B=e^{-\beta H_B}/Z$, $\beta=(k_BT)^{-1}>0$. Put
\[
 s=S(V)-S(V'),\qquad p=S(A')-S(A),\qquad
 Q=\operatorname{tr}H_B(\rho'_B-\tau_B).
\]
Then, with $I_{V:A:B}=S(V)+S(A)+S(B)-S(VAB)$,
\begin{equation}
 \beta Q+p=s+I_{V':A':B'}+D(\rho'_B\Vert\tau_B)\ \ge s.
 \label{eq:entropyaccount}
\end{equation}
If the marginal entropy of $C$ is returned, $s$ equals the decrease of
$S(S|C)$. In a sequence of such corrections, a bounded total purity consumption
and a positive asymptotic entropy-removal rate $\dot s$ imply
$\liminf Q_{[0,t]}/t\ge k_BT\dot s$.
\end{proposition}
\begin{proof}
Unitary invariance of total entropy and the initial product give
$\Delta S(V)+\Delta S(A)+\Delta S(B)=I_{V':A':B'}$.
Also $\beta Q=\Delta S(B)+D(\rho'_B\Vert\tau_B)$.
Combine the identities and sum over corrections. The rate statement follows when
the accumulated $p$ divided by elapsed time tends to zero.
\end{proof}

This is the Landauer entropy balance, with the memory-purity term displayed
explicitly; the equality and the role of accessible side information are standard
\cite[Theorem~3 and Section~5.1]{ReebWolf2014}. For initially correlated or
nonthermal inputs the same calculation instead gives
\[
 \beta Q+p=s+(I_{\rm out}-I_{\rm in})+
 D(\rho'_B\Vert\tau_B)-D(\rho_B\Vert\tau_B).
\]
These consumed correlations and bath free energy can pay for correction. A classical
copy of the error permits a reversible conditional flip with $s=0$, although the
system's marginal entropy decreases. Negative quantum conditional entropy gives
another reason to include the side information. A pure-memory swap has $p=s$ and
can have $Q=0$. A retained oscillator's closed unitary loop also has no entropy
removal by itself. Thus a count of corrections or a marginal noise entropy is not
automatically $\dot s$. In a cyclic device $\dot s$ must be established from its
actual noise and controller dynamics. For nondegenerate systems, heat and supplied
work are further separated by the change in internal energy.

The memory can also be reused without being decorrelated after each correction.
Let correction $i$ act on $VA$ and a fresh thermal port $B_i$, initially independent
of the whole state of $VA$ and all earlier ports, and put
$s_i=S(V_i)-S(V_i')$ and $x_i=S(B_i')-S(B_i)$.
Suppose each intervening noise step is a CPTP map on $V$, tensored with the
identity on $A$ and the previously exported ports; updates involving $A$ belong
to the charged correction. Entropy conservation gives
\[
 x_i=s_i-[S(A_i')-S(A_i)]
       +I(V_i':A_i')-I(V_i:A_i)+I(V_i'A_i':B_i').
\]
Data processing between corrections and telescoping therefore imply
\begin{equation}\label{eq:reused-memory}
 \sum_{i=1}^n x_i\ge
 \sum_{i=1}^n s_i-[S(A_n')-S(A_1)]-I(V_1:A_1).
\end{equation}
For a fixed finite $A$, the last two terms are bounded. Reused memory purity and
initial correlations cannot absorb a positive entropy-removal rate indefinitely.
Here $x_i$ refers to each exported port's marginal entropy. Summing those marginals
can overcount correlated records. There is a stronger joint-port statement under
the same fresh-input and local-noise assumptions:
\begin{equation}\label{eq:joint-exhaust}
 S(B_{1:n}')-\sum_{i=1}^n S(\tau_{B_i})\ge
 \sum_{i=1}^n s_i-[S(A_n')-S(A_1)]-I(V_1:A_1).
\end{equation}
To see this, retain all previously exported ports in the mathematical state.
A noise step on $V$, with the identity on those ports, changes its
total entropy by $\Delta S(V)-\Delta I(V:A B_{<i})\ge\Delta S(V)$, by data
processing. The corrections preserve total entropy, and each new port adds
$S(\tau_{B_i})$. At the end, subadditivity between $VA$ and all exported ports
gives Eq.~\eqref{eq:joint-exhaust} after the $V$ entropies telescope. Thus finite
reused memory cannot hide the maintained rate in correlations between outputs.

The same account applies when the noise also acts on the retained controller.
Put $W=VA$ and let each correction be a unitary on $WB_i$ with a fresh,
independent input $\tau_{B_i}$, uncorrelated with $W$ and the earlier ports.
Noise between corrections acts on $W$ alone.
With $s_i^W=S(W_i)-S(W_i')$, the exact identity is
\begin{equation}\label{eq:joint-VA-exhaust}
 S(B_{1:n}')-\sum_iS(\tau_{B_i})
 =\sum_i s_i^W+I(W_n':B_{1:n}')+\sum_{i=1}^{n-1}\ell_i
 \ \ge\ \sum_i s_i^W,
\end{equation}
where $\ell_i=I(W_i':B_{1:i}')-I(W_{i+1}:B_{1:i}')\ge0$ is the information
lost to the intervening noise. Indeed, that noise changes the entropy of the
joint state by $\Delta S(W)+\ell_i$; the correction preserves it. Summing these
changes and expanding the final entropy as $S(W_n')+S(B_{1:n}')-I(W_n':B_{1:n}')$
gives the identity. Returning an old port to an operation would change these
hypotheses. The joint $VA$ rate, rather than a sum of system marginal entropies,
is what an apparatus with noisy internal memory must export.

\begin{proposition}[A finite-rate chiral exhaust floor]\label{prop:chiral-floor}
Consider a unitary, zero-chemical-potential chiral conformal exhaust with central charge
$c_{\rm ex}>0$, thermal input at $T\ge0$, and an extensive entropy-current account.
Energy and temperature are defined by the same translation generator and clock.
Assume that the outgoing modes export net dimensionless entropy current
$x\ge\dot s>0$, where $\dot s$ is the actual maintenance entropy-removal rate after
all retained side information and sublinear resource consumption are accounted for.
The entropy-current maximum at fixed energy is the thermal one. Then the excess
heat current over the incoming thermal stream obeys
\begin{equation}\label{eq:chiral-floor}
 \dot Q\ge k_BT\dot s+\frac{3\hbar}{\pi c_{\rm ex}}\dot s^2.
\end{equation}
\end{proposition}
\begin{proof}
Write $a=\pi c_{\rm ex}k_B^2/(12\hbar)$. The thermal energy and dimensionless
entropy currents are $J_E(T)=aT^2$ and $J_S(T)=2aT/k_B$.
The first is the standard chiral CFT current~\cite[Eq.~(1)]{BernardDoyon2012};
the second follows from its extensive Gibbs entropy. Entropy maximization at
fixed energy is the Gibbs variational principle, underlying the single-channel
information-flow bound~\cite{Pendry1983}.
Choose $T_*=T+k_Bx/(2a)$ so that the thermal output has the required entropy.
Any output with at least that entropy has energy current at least $aT_*^2$.
Subtracting $aT^2$ gives
\[
 \dot Q\ge k_BTx+\frac{k_B^2x^2}{4a}
            =k_BTx+\frac{3\hbar x^2}{\pi c_{\rm ex}}.
\]
The expression increases for $x\ge0$.
\end{proof}

The current account is a physical condition on the exported modes, not the entropy
of a Type~III local interval. Equation~\eqref{eq:joint-exhaust} establishes the
needed joint entropy rate in a fresh-port account; matching those ports to a conformal current
must preserve its capacity and entropy budget. At fixed capacity and positive
rate the quadratic term survives even at $T=0$. Thermal output saturates the
channel inequality, but this alone does not construct a correcting apparatus.
The currents are long-window, extensive Gibbs limits. For an observation window
$\tau$ with net exported entropy $X_\tau$, bounded rates and sublinear boundary
corrections give
\[
 Q_\tau\ge k_BT X_\tau+
 \frac{3\hbar X_\tau^2}{\pi c_{\rm ex}\tau}+o(\tau).
\]
The current formula is not an exact inequality for an arbitrarily short finite
packet. Applying the Gibbs bound to the entire output also includes correlations
between packets; it does not require each output to be a separate thermal state.

\begin{proposition}[Output deficit and direct renewal]
\label{prop:export-deficit}
Use the complete outgoing carrier $B_\tau$, its physical translation
generator $H_\tau$, and the thermal input of
Proposition~\ref{prop:chiral-floor}. Let the actual charged entropy
removal satisfy $s_\tau/\tau\to x>0$, set
$T_*=T+6\hbar x/(\pi c_{\rm ex}k_B)$, and write
$\tau_*=e^{-H_\tau/(k_BT_*)}/Z_\tau$.
Assume the extensive Gibbs limits used above and that input, charge,
interaction, resource and boundary remainders are $o(\tau)$.
If the full entropy account gives
$L_\tau=S(\rho_{B_\tau})-S_T-s_\tau\ge-o(\tau)$, then
\begin{equation}
 Q_\tau-\tau\left(k_BTx+\frac{3\hbar x^2}{\pi c_{\rm ex}}\right)
 =k_BT_*\left[D(\rho_{B_\tau}\Vert\tau_*)+L_\tau\right]+o(\tau).
 \label{eq:output-deficit}
\end{equation}
Thus asymptotic equality requires, and under these hypotheses is
equivalent to, both deficit and entropy-slack rates tending to zero.

For Q-RAIF~C's original quartet preparation followed by direct
per-species ring renewal~\cite{QRC}, let $d_{\rm ring}$ bound its
reset error against the actual old outside marginal, $\eta$ its
residual zero-mode amplitude, and $e_H$ the full-joint clock comparison
error over the specified horizon. Suppose each round has disjoint
emission modes, these bounds hold uniformly, and no later converter
changes the exported modes. With
$t=d_{\rm ring}+2\arcsin\eta+2e_H<1/2$, the output of $M$ rounds obeys
\begin{equation}
 D(\rho_{B_\tau}\Vert\tau_*)\ge
 M\left[h(\epsilon)-2h(t)\right],
 \qquad h(p)=-p\ln p-(1-p)\ln(1-p),
 \label{eq:direct-export-deficit}
\end{equation}
when the zero-chemical-potential free-rail reference factorizes across
the original species. If $M/\tau\to\nu$ and $h(\epsilon)>2h(t)$,
the asymptotic heat excess above the floor is at least
$k_BT_*\nu[h(\epsilon)-2h(t)]$ per unit time.
Temporal independence of thermal packets is not required.
\end{proposition}
\begin{proof}
The finite-carrier Gibbs identity is
$D(\rho\Vert\tau_*)=(E_\rho-E_*)/(k_BT_*)-S(\rho)+S_*$.
Here $S_*-S_T=x\tau+o(\tau)$ and
$E_*-E_T=\tau[k_BTx+3\hbar x^2/(\pi c_{\rm ex})]+o(\tau)$,
which gives \eqref{eq:output-deficit}.

The current purifier's occupation pair is $00$ or $11$, with
probabilities $1-\epsilon,\epsilon$. Operations on the receiver
leave its joint marginal with previous outputs unchanged. The reset
certificate compares that pair with the same actual past marginal;
it need not be independent of the current readout record.
The exact linear scattering column emitting an old zero mode $d_j$
has squared norm $1-\eta^2$. CAR therefore give the extracted mode
$b_j=\sqrt{1-\eta^2}\,d_j+\eta e_j$, where $e_j$ is a canonical
orthogonal mode in the actual body and incoming field. A two-mode
rotation of angle $\arcsin\eta$ compares $b_j$ with $d_j$, retaining
every old reference. The two purifier species cost $2\arcsin\eta$;
replacement of the comparison joint state and its outside marginal
costs at most $2e_H$.

Measure their even occupations in the disjoint emitted modes, giving
$X_{1:M},Y_{1:M}$. For each round the joint law of $(X_j,X_{<j})$
is within variation $t$ of
$\operatorname{Bern}(\epsilon)\otimes P_{X_{<j}}$, and
$\Pr(X_j\ne Y_j)\le t$. Conditional binary entropy continuity and
the binary Fano bound give
$H(X_j|X_{<j})\ge h(\epsilon)-h(t)$ and
$H(X_j|Y_j)\le h(t)$. Hence
$I(X_{1:M}:Y_{1:M})\ge M[h(\epsilon)-2h(t)]$.
The reference measurement law is $q_Xq_Y$, even though each sequence
may be thermally correlated. Data processing and
$D(P_{XY}\Vert q_Xq_Y)=I(X:Y)+D(P_X\Vert q_X)+D(P_Y\Vert q_Y)$
prove \eqref{eq:direct-export-deficit}.
\end{proof}

All entropies in this proposition use the finite-carrier account and
its extensive limit, not a trace entropy of a Type~III interval.
A fixed-charge canonical reference needs its sector-weight correction;
it has the same rate conclusion only when that correction is
sublinear. A conditional account must use the same charged removal.
For a complete usable system $W=ZM$ with a conserved classical record
$Z$, write $s^W=S(W)-S(W')$. A unitary interaction with a port $B$ gives
\[
 S(B'|Z)-S(B|Z)
 =s^W+I(M':B'|Z)-I(M:B|Z).
\]
When the port is initially independent of $W$, its thermal entropy
replaces $S(B|Z)$ and the last term vanishes. Old purifications remain
in the physical state; they are usable side information only if the
apparatus has that access. Fresh-port sequences with noise acting only
on $M$ also accumulate nonnegative information losses to earlier
outputs. Reading an old output or updating $Z$ needs a joint account
of that operation.

For a classical record $X$ retained at input and output, put
$L_c=S(B_{\rm out}|X)-S(B_{\rm in}|X)-s_\tau$. When the input marginal
entropy is $S_T$, the same Gibbs identity is
\begin{equation}
 \begin{aligned}
 Q_\tau-\tau f(x)
 &=k_BT_*\left[D(\rho_{B_{\rm out}}\Vert\tau_*)+L_c
       +I(X:B_{\rm out})-I(X:B_{\rm in})\right]+o(\tau),\\
 f(x)&=k_BTx+\frac{3\hbar x^2}{\pi c_{\rm ex}}.
 \end{aligned}
 \label{eq:conditional-export-ledger}
\end{equation}
Indeed $L_\tau=L_c+I_{\rm out}-I_{\rm in}$, while
$D(\rho_{XB}\Vert\rho_X\otimes\tau_*)
=D(\rho_B\Vert\tau_*)+I(X:B)$. These are two forms of one account;
thermal input marginals alone do not remove old correlations.
The direct protocol's deficit lower bound tends to
$\nu h(\epsilon)>0$ as its certified error vanishes. A physical
converter that redistributes those correlated records, with its own
work and retained memory charged, is therefore necessary for that
route to approach the sharp floor.

The incoming state and conserved charges matter. For a nonthermal stream with
energy current $e_{\rm in}$ and dimensionless entropy current $\sigma_{\rm in}$,
put $E_{\rm th}(\sigma)=3\hbar\sigma^2/(\pi c_{\rm ex})$. The same variational
argument gives
\begin{equation}\label{eq:nonthermal-exhaust}
 \dot Q+e_{\rm in}-E_{\rm th}(\sigma_{\rm in})
 \ge k_BT_\sigma x+\frac{3\hbar x^2}{\pi c_{\rm ex}},
 \qquad T_\sigma=\frac{6\hbar\sigma_{\rm in}}{\pi c_{\rm ex}k_B}.
\end{equation}
The nonequilibrium energy above the entropy-matched thermal stream can pay for
export. A pure coherent input therefore cannot be charged as though it were a
zero-energy thermal reservoir. At nonzero chemical potential the corresponding
Gibbs comparison must hold all retained charges fixed. For example, with a
level-one $U(1)$ current normalized as in Ref.~\cite[Eq.~(84)]{BernardDoyon2015},
the unitary CFT Gibbs currents in units $\hbar=k_B=1$ are
$e=\pi c_{\rm ex}T^2/12+\pi\mu^2/2$, $j=\pi\mu$ and
$\sigma=\pi c_{\rm ex}T/6$~\cite[Eq.~(93)]{BernardDoyon2015}.
Equal incoming and outgoing charge currents cancel the chemical energy in their
difference; an apparatus that absorbs charge can instead consume it.
Equation~\eqref{eq:chiral-floor} does not cover that consumption without a charge
ledger.

For independent exhausts at a common temperature, with central charges $c_j$
and nonnegative entropy currents $x_j$, the same argument and Cauchy--Schwarz give
\begin{equation}\label{eq:parallel-exhaust}
 \dot Q_{\rm tot}\ge k_BT\dot s+
 \frac{3\hbar\dot s^2}{\pi\sum_jc_j},\qquad \sum_jx_j\ge\dot s.
\end{equation}
The capacity term is sharp at the channel level: equality requires thermal outputs
and $x_j=c_j\dot s/\sum_kc_k$. Additional channels are an additional physical
resource. Intermittent export through the same capacity cannot lower the
long-window Gibbs floor: apply it to the complete output with its total entropy.
If a local-current description is valid, $\overline{x^2}\ge\bar x^2$ gives the
same conclusion.

Heat alone does not bound every geometric observable. Even a positive $E_0$ can
give zero or negative $\delta\mathcal C$ in Eq.~\eqref{eq:dC}; its response depends
on other moments and on the background. A positive weighted integral
$\int K(u)F(u)du\ge K_{\min}\int Fdu$ requires both a lower bound on $K$ over
the exhaust and $F\ge0$. Quantum excess stress need not have that sign. There is,
however, a simple JT observable which reads $E_0$ without either assumption.

\begin{proposition}[A heat account read through the dilaton]\label{prop:heatslope}
Compare two JT solutions with the same $H_0$, boundary coefficient data and
coordinate and clock normalization. Suppose their source difference $F$ is the
exhaust of the correction in Proposition~\ref{prop:entropyaccount}, with finite
absolute second moment and $\int Fdu=Q$. This last equality requires that the
specified outgoing channel carry all the heat in the same energy normalization.
Then its retarded coefficient difference obeys
\begin{equation}
 -\frac{A'(+\infty)}\kappa=Q\ \ge\ k_BT(s-p).
 \label{eq:heatslope}
\end{equation}
For a compact exhaust and any $u_2>u_1$ after it, the left side is exactly
$-[A(u_2)-A(u_1)]/[\kappa(u_2-u_1)]$. No pointwise sign of $F$ is required.
\end{proposition}
\begin{proof}
Differentiate Eq.~\eqref{eq:retarded}: $A'(u)=-\kappa\int_{-\infty}^uF(z)dz$.
Take the late limit and apply Eq.~\eqref{eq:entropyaccount}.
\end{proof}

The reference for $Q$ is the incoming thermal reservoir transported through the
same exhaust, before its heat uptake; $H_B$ is its $u$-translation generator.
A thermal background has a nonintegrable stationary flux, so the JT solution here
is the difference from that same transported thermal background.
The identification $E_0=Q$ is an energy-transport hypothesis:
other exhaust channels, energy stored in the device, work beams crossing the
observation surface and redshift must be accounted for separately. With a fixed
transport factor $E_0=\zeta Q$, Eq.~\eqref{eq:heatslope} acquires the same factor
$\zeta$. The coefficient slope is read from the dilaton in the declared matched
coordinate frame, not an invariant under arbitrary changes of boundary clock.
The relation links an actual entropy-removal process to a specified dilaton reading.
The JT dynamics and heat-routing map are its geometric inputs.


\subsection{A horizon reading of the maintenance price}

For a matched running single-chiral JT source, let
$\mathcal C=R^2>0$ and
$\xi=(bu^2+2au-\gamma)/(2R)>0$ on the regular future outer patch.
The marginal horizon has $\Phi_h=R$ and
$S_h=S_0+2\pi R/\kappa$. In its rapidity clock $d\eta=du/\xi$,
\[
 \frac{dS_h}{d\eta}=2\pi F\xi^2.
\]
When the charged heat is exactly $F\xi^2$ and the physical exhaust
input is thermal at $T_H=1/(2\pi)$ in this same clock, the
long-window channel bound gives
\[
 \overline{\dot S_h}\ge\dot s+\frac{6\dot s^2}{c_{\rm ex}}.
\]
Appendix~\ref{app:running-clocks} proves the horizon balance and its
domain conditions. The entropy here is classical Wald entropy; a
generalized-entropy statement needs its matter contribution.

Steady positive flux cannot be extrapolated in one fixed $u$ chart:
its boundary dilaton numerator eventually vanishes. Nor does a new
clock prepare a reservoir. For an actual source field thermal in
$\eta$, the vacuum-subtracted physical flux contains
$q_{\rm ref}=-c(\delta'+\delta^2/2)/(24\pi)$, where
$\delta=R'/R$. Its self-consistent equations are
\[
 R'=R\delta,\qquad
 \delta'=\frac{\kappa q-R\delta}{\kappa c/(24\pi)}
                         -\frac{\delta^2}{2}.
\]
The positive growing family of
Proposition~\ref{prop:thermal-source-clock} has
$R=R_0+\kappa q\eta+O(1)$ and finite accumulated reference
correction. Hence its asymptotic horizon rate is
$2\pi q=s+6s^2/c$. A partially transmitting local thermal contact
produces a nonthermal output with a strictly larger energy price at
the same entropy current; changes of angle add their own pump work.
The detailed state, source and clock calculations are in
Appendix~\ref{app:running-clocks}. Their global source-clock solution
does not grant infinite life to a finite material energy reserve.

\section{The Factorization Problem}
\label{sec:factorization}

\subsection{No canonical tensor factorization}

A Hilbert space of dimension $d=d_1d_2$ admits a tensor factorization,
but its Hilbert structure alone does not select one. Commuting observable
algebras can induce the subsystem structure~\cite{Zanardi2001,Zanardi2004}.
Which of those algebras an apparatus can use is a physical question about
its couplings and resources.

\subsection{Coarse-graining and effective descriptions}

Paper~A~\cite{Liu2026PaperA} gives the following conditional separation.

\begin{proposition}[Conditional coarse-graining separation]
\label{prop:inequiv}
Let $\Phi_1\ne\Phi_2$ be finite-dimensional coarse-graining channels with a common output space. Their outputs differ for generic pure inputs. Entropy profiles differ only under the additional nondegeneracy hypothesis stated in Paper~A, and metric inequivalence follows only when a supplied reconstruction dictionary separates those profiles modulo diffeomorphisms.
\end{proposition}

Channel inequivalence therefore establishes a difference of effective descriptions, not by itself a difference of spacetime metrics. Paper~A also decides, for qubit reductions, exactly when two channels give the same output entropy on every pure input, and it constructs one global state, one family of local noise channels and one fixed reconstruction rule that yield non-isometric finite metrics. Metric inequivalence of spacetime geometries remains a conditional model claim throughout this paper.

\subsection{Geometry from entanglement}

In a holographic model the Ryu--Takayanagi formula~\cite{RyuTakayanagi2006} relates an extremal area to boundary entropy,
\begin{equation}
S_A = \frac{\text{Area}(\gamma_A)}{4 G_N \hbar},
\end{equation}
and motivates an entanglement account of connectivity~\cite{VanRaamsdonk2010}.
Such a reconstruction requires its dictionary. Outside a specified model,
the accessible data determine operator relations rather than a spacetime metric.

\section{A Structural Distinction: Gauge Forces and Gravity}
\label{sec:distinction}

Gauge interactions evolve matter and connection fields within a supplied localization structure. The algebraic proposal asks whether gravity instead describes how physical correlations, accessible operations and their stability select that structure. With a geometric dictionary, changing those relations can change effective distances and causal structure. The preceding sources test a necessary interface: changed access reaches geometry through a specified matter process.

If the proposal holds, quantum gravity must describe selection of localization as well as fields within it. This is not an explanation of why quantization fails. Background independence~\cite{Rovelli2004,Smolin2006} asks for covariance under changes of selected geometric structure, not equality of the metrics assigned to operationally different algebras. Selection is constrained by interactions and stability, rather than subjective choice~\cite{Liu2026PaperB}; changed knowledge affects stress only through an operation.

A common reconstruction dictionary gives every matter representation in one selected geometric sector the same metric. This conditional universality is not a derivation of free fall or equality of inertial and gravitational mass. Degrees outside the sector cannot be declared geometrically irrelevant while changing its metric: a dark-sector application would need an enlarged source algebra and explicit backreaction map. No dark-matter or dark-energy conclusion follows here.

\section{Technical Formulation}
\label{sec:technical}

\subsection{Setup}

We separate the operator-algebraic and finite-dimensional regimes result by result. Modular statements use a Type~III$_1$ ambient factor with a faithful normal state and explicit cyclic and separating hypotheses. Density matrices, traces, the quantum Fisher information metric and the qubit chain are finite-dimensional statements. The Rindler example is an algebraic quantum field theory check. The linearized Einstein equation is imported only in the holographic regime of Theorem~\ref{thm:linearized}, and the JT model of Section~\ref{sec:source} is a supplied two-dimensional gravity theory.

\begin{definition}[Selected geometric algebra]
\label{def:accessible}
A \textbf{selected geometric algebra} $\mathcal{A}_{\mathrm{geo},\mathbf{c}}$ is a von Neumann subalgebra chosen as input to a specified algebra-to-geometry reconstruction model. It is not Paper~E's probe-reachable effect system, pointer or record algebra, or environmental fixed set. A coarse-graining model must declare the channel or inclusion relating any of those objects to $\mathcal{A}_{\mathrm{geo},\mathbf{c}}$.
\end{definition}

No canonical tensor factorization is assumed. State restrictions and relative-entropy data on $\mathcal{A}_{\mathrm{geo},\mathbf{c}}$ can be associated, as in holography and tensor networks, with an effective distance structure. That structure depends on the choice of $\mathcal{A}_{\mathrm{geo},\mathbf{c}}$, on the stability of correlations under coarse-grained dynamics and on the redundancy of the encoding; it is an effective description valid in the context the algebra defines.

\subsection{Central conjecture}

\begin{conjecture}[Gravity as adiabatic algebra evolution]
\label{conj:main}
Gravitational dynamics corresponds to an \textbf{adiabatic flow} of selected geometric algebras $\mathcal{A}_{\mathrm{geo},\mathbf{c}}(t)$:
\begin{enumerate}
    \item the global state evolves unitarily, $|\Psi_U(t)\rangle = U(t)|\Psi_U(0)\rangle$;
    \item a model-dependent rule selects a geometric candidate $\mathcal{A}_{\mathrm{geo},\mathbf{c}}(t)\subset\mathcal M_U$ at each time;
    \item as the state evolves, the selected candidate may shift, $\mathcal{A}_{\mathrm{geo},\mathbf{c}}(t) \to \mathcal{A}_{\mathrm{geo},\mathbf{c}}(t + dt)$;
    \item this shift manifests phenomenologically as the dynamical curvature of spacetime.
\end{enumerate}
Unitary evolution of observables within a fixed $\mathcal{A}_{\mathrm{geo},\mathbf{c}}$ manifests as gauge interaction. Gauge dynamics is evolution within the algebra (fixed stage, moving actors); gravitational dynamics is evolution of the algebra itself (moving stage).
\end{conjecture}

The conjecture answers one objection: if algebras were kinematical background, gravity could not be dynamical. Here the selection of which algebra is stable is state dependent, and state evolution induces algebra flow. The selection rule that produces the path is not derived in this paper; Section~\ref{sec:gap} returns to what a physical selection would have to supply. The adiabatic assumption takes algebra transitions to be slow compared with dynamics inside the algebra. Rapid transitions would correspond to strong-field regimes where the effective geometric description breaks down, and the dynamics of non-adiabatic transitions is left open.

\subsection{Comparison with quantum information geometry}

For comparison consider the quantum Fisher information metric (QFIM) on a smooth finite-dimensional family of faithful effective states $\rho(\lambda)$~\cite{Petz1996,Bengtsson2006},
\begin{equation}
g_{\mu\nu}(\lambda) = \frac{1}{2} \text{Tr}\left( \rho(\lambda) \{ L_\mu, L_\nu \} \right),
\qquad
\partial_\mu \rho = \frac{1}{2}(\rho L_\mu + L_\mu \rho).
\label{eq:qfim-d}
\end{equation}
It is coordinate independent, reduces to the classical Fisher metric in commuting limits, and measures how hard nearby states are to distinguish. It is also positive semidefinite, whereas a Lorentzian spacetime metric is indefinite and carries causal structure. Nothing in Eq.~\eqref{eq:qfim-d} supplies the negative direction, so the QFIM is a metric on model space and an analogy, not the spacetime metric. A change $\mathcal{A}_{\mathrm{geo},\mathbf{c}} \to \mathcal{A}_{\mathrm{geo},\mathbf{c}} + \delta\mathcal{A}_{\mathrm{geo},\mathbf{c}}$ may change this state-space metric; calling the change a spacetime perturbation requires an additional Lorentzian reconstruction map. In AdS/CFT, relative entropy and the entanglement first law constrain bulk perturbations~\cite{Lashkari2014,Faulkner2014}; those results use the holographic dictionary and do not make an arbitrary boundary QFIM the bulk metric.


\section{Algebraic Foundation: Modular-Stable Candidates}
\label{sec:modular}

This section places one candidate class in Tomita--Takesaki modular theory. It does not translate measurement accessibility or environmental redundancy into modular conditions, and it does not select a unique algebra.

Standard modular definitions and their two checks are in Appendix~\ref{app:modular-details}.

\subsection{Modular definition of candidate algebras}
\label{sec:modular-def}

Let $\mathcal M_U$ be a Type~III$_1$ factor in standard form on $\mathcal H_U$, let $|\Psi_U\rangle$ be cyclic and separating, and write $\omega(\cdot)=\langle\Psi_U|\cdot|\Psi_U\rangle$. Then $\omega$ is faithful and normal and its ambient modular group $\sigma_t^{\omega,U}$ exists. Let $\{\mathcal E_t\}_{t\ge0}$ be a normal unital completely positive (UCP) semigroup acting in the Heisenberg picture on $\mathcal M_U$.

\begin{definition}[Modular-stable geometric candidate]
\label{def:modular-accessible}
A von Neumann subalgebra $\mathcal{A}_{\mathrm{geo},\mathbf{c}}\subset\mathcal M_U$ is a \textbf{modular-stable geometric candidate} if:
\begin{enumerate}
\item[\textup{(U1)}] \textbf{Faithful restriction.} $\omega|_{\mathcal{A}_{\mathrm{geo},\mathbf{c}}}$ is faithful and normal. In the faithful ambient setting this holds for every unital subalgebra, so it is a regularity condition, not a selector.
\item[\textup{(U2)}] \textbf{Ambient modular stability.} $\sigma_t^{\omega,U}(\mathcal{A}_{\mathrm{geo},\mathbf{c}})=\mathcal{A}_{\mathrm{geo},\mathbf{c}}$ for all real $t$.
\item[\textup{(U3)}] \textbf{Semigroup stability.} $\mathcal E_t(\mathcal{A}_{\mathrm{geo},\mathbf{c}})\subseteq\mathcal{A}_{\mathrm{geo},\mathbf{c}}$ for all $t\ge0$.
\end{enumerate}
\end{definition}

\begin{remark}[Typed relation to the measurement objects]
\label{rem:translation}
Conditions (U1)--(U3) do not derive Paper~E's implementable effects $\mathrm{Eff}$, their operator system $\mathrm{Eff}_{\mathrm{span}}$ or the pointer algebra $\mathcal P_{\mathrm{env}}$, and setwise semigroup stability is not the exact fixed set $\operatorname{Fix}(\mathcal E)$. A downstream model may specify a channel or inclusion from those objects into $\mathcal{A}_{\mathrm{geo},\mathbf{c}}$; without it the relation is motivational. Environmental redundancy does not imply faithfulness of a state restriction. A causal restriction map and a Lindblad evolution are different types, and each model must supply its own $\mathcal E_t$ and verify invariance.
\end{remark}

\subsection{Selection is not settled by these conditions}
\label{sec:modular-uniqueness}

\begin{remark}[No uniqueness from (U1)--(U3)]
\label{thm:uniqueness}
Conditions (U1)--(U3), even with only scalars fixed by both dynamics, do not determine a unique candidate. There are four independent obstructions.
\begin{enumerate}
\item Uniqueness of the modular flow for a fixed pair $(\mathcal A,\omega)$ does not select the algebra $\mathcal A$.
\item A normal completely positive map may preserve two subalgebras separately without preserving their join, because it is not multiplicative. Join preservation holds for modular automorphisms but not for a general $\mathcal E_t$; it does hold for subalgebras of the fixed algebra of a state-preserving dynamics, on which $\mathcal E_t$ acts as the identity~\cite{FoundationI}.
\item Pointwise ergodicity constrains fixed elements, while (U2)--(U3) are setwise conditions; no implication to a unique maximal element of the lattice was proved.
\item With a faithful ambient state, $\mathcal M_U$ itself satisfies (U1)--(U3), so a bare maximality rule selects the ambient algebra unless an independent seed, properness, read-out or minimal-hull condition is supplied.
\end{enumerate}
Replacing $\mathcal B(\mathcal H_U)$ by a faithful Type~III factor repairs the existence of the ambient modular flow but none of these selection failures. The class of candidates may contain several members that are not unitarily equivalent, and classifying it requires additional structure. A nontrivial selection principle could start from a foreground seed and construct a minimal invariant hull, or impose an explicit read-out constraint. Foundation~I carries out the first route for a supplied seed: under state-preserving dynamics the hull of a seed in the fixed algebra lies between that seed and the fixed algebra. For a minimal compact symmetry, it equals the seed-stabilizer fixed algebra when that algebra is preserved by the hull generators and the hull contains the group-fixed algebra~\cite{FoundationI}. These hypotheses establish uniqueness relative to the seed; the seed's physical origin is a further question.
\end{remark}

\subsection{Connection to linearized gravity}
\label{sec:linearized}

\begin{theorem}[Entanglement first law in a density-matrix setting]
\label{thm:first-law}
Let $\rho_0$ be a faithful finite-dimensional density matrix and
$\rho_\varepsilon=\rho_0+\varepsilon\,\delta\rho+O(\varepsilon^2)$ a normalized
perturbation, with $K_0=-\log\rho_0$. Then $\delta S=\delta\langle K_0\rangle$.
An infinite-dimensional extension requires entropy differentiability and control
of the modular-energy variation; faithfulness alone is insufficient. For a Type~III
local algebra the statement must instead use Araki relative entropy.
\end{theorem}

\begin{proof}[Proof sketch]
$S(\rho_\varepsilon\|\rho_0)=-S(\rho_\varepsilon)+\langle K_0\rangle_{\rho_\varepsilon}$ plus a constant is nonnegative and stationary at $\varepsilon=0$.
\end{proof}

Conjecture~\ref{conj:main} posits a family $\mathcal{A}_{\mathrm{geo},\mathbf{c}}(\lambda)$. Modular theory treats two kinds of perturbation differently. For a state perturbation on a fixed algebra the Connes cocycle $u_t=(D\omega':D\omega)_t$ relates the automorphism groups by
$\sigma_t^{\omega'}(A)=u_t\sigma_t^\omega(A)u_t^*$~\cite{Connes1994,Takesaki2003}. This is not a left-multiplication formula for the two full modular operators. If the algebra family is generated by the positive translation of an HSMI,
$\mathcal A_{\mathrm{geo},\mathbf c}(a)=e^{iaP}\mathcal A_{\mathrm{geo},\mathbf c}(0)e^{-iaP}$,
and the vacuum is fixed by this translation, its full modular generator obeys
$K(a)=e^{iaP}K(0)e^{-iaP}$. On a common smooth domain,
\begin{equation}
\label{eq:deltaK-translation}
\frac{dK(a)}{da}\Big|_{a=0}=i[P,K(0)] .
\end{equation}
Equation~\eqref{eq:deltaK-translation} concerns an HSMI-generated family and a derivative
of matrix elements on that domain, not arbitrary subalgebra changes or an operator-norm
Taylor expansion. The one-sided CHM charge used for a ball below is distinct from the
full modular generator; their relevant local-unitary expectation differences agree,
but the operators themselves do not.

\begin{theorem}[Linearized gravity from entanglement, {\cite{Faulkner2014}}]
\label{thm:linearized-Faulkner}
In a holographic CFT dual to Einstein gravity in asymptotically AdS spacetime, the entanglement first law $\delta S = \delta\langle K_{\mathrm{mod}}\rangle$ for all boundary balls is equivalent to the linearized bulk Einstein equations $G_{\mu\nu}^{(1)} + \Lambda g_{\mu\nu}^{(1)} = 8\pi G_N\, T_{\mu\nu}^{(1)}$.
\end{theorem}

The derivation uses the ball modular Hamiltonian~\cite{CasiniHuertaMyers2011}
\begin{equation}
K_{\mathrm{mod}}^B = 2\pi \int_B d^{d-1}x\, \frac{R^2 - |\mathbf{x} - \mathbf{x}_0|^2}{2R}\, T_{00}(x),
\end{equation}
together with the Ryu--Takayanagi formula and the relation between bulk and boundary modular Hamiltonians~\cite{JLMS2016}; complementary arguments are due to Jacobson~\cite{Jacobson2016}.

\begin{theorem}[Conditional holographic matching for an HSMI-generated family]
\label{thm:linearized}
Let $\mathcal{A}_{\mathrm{geo},\mathbf{c}}$ be a modular-stable candidate in a holographic CFT vacuum representation, and assume:
\begin{enumerate}
\item[(H1)] the AdS/CFT correspondence holds;
\item[(H2)] the Ryu--Takayanagi formula and its quantum corrections hold;
\item[(H3)] the JLMS relation between bulk and boundary modular Hamiltonians holds;
\item[(H4)] the compared algebra family is an explicitly supplied HSMI/ball-region family whose generator $P$ implements the boundary-region variation used in the holographic first-law argument;
\item[(H5)] the regulated entropy variation, or its relative-entropy replacement, satisfies the hypotheses of the Faulkner et al.\ theorem for all required balls.
\end{enumerate}
Then Eq.~\eqref{eq:deltaK-translation} describes the modular-generator variation within that family, and the holographic first law matches the linearized Einstein equations by the imported result:
\begin{align}
 \text{HSMI/ball family (H4)}&\Longrightarrow\delta K_{\mathrm{mod}}
 \xrightarrow{\text{(H5)}}\delta S=\delta\langle K_{\mathrm{mod}}\rangle,
 \nonumber\\
 \text{(H1)--(H3)}&\Longrightarrow
 \delta(G_{\mu\nu}+\Lambda g_{\mu\nu})=8\pi G_NT_{\mu\nu}^{(1)}.
\end{align}
The theorem does not show that a generic stability-selected algebra variation is an HSMI translation, nor that HAFF algebra flow is gravitational dynamics.
\end{theorem}

\subsection{The circularity issue}

Jacobson's thermodynamic derivation~\cite{Jacobson1995} presupposes a background geometry, while HAFF wants geometry to emerge from algebra. Three strategies are discussed in the literature. (A) In the algebraic programme, wedges are defined through modular inclusions without a background metric, and Theorem~\ref{thm:Wiesbrock} recovers a positive translation from an HSMI; constructing a sufficient net and showing that the candidates of this paper form such inclusions remain open tasks. (B) A bootstrap starts from a seed geometry, derives linearized equations from the first law, updates the geometry and iterates. (C) In holography the boundary theory supplies the algebra and the bulk geometry is derived from boundary data. Theorem~\ref{thm:linearized} uses strategy (C) and assumes the HSMI/ball interface. A geometry-free reconstruction from the present candidate class remains open.


\section{Selection and response over growing windows}

\subsection{A conjecture on a sustainable selection--response loop}

The unresolved question is whether one reused local body can combine
selection, renewal, sharp exhaust and informative geometric feedback.
Fix two physical preparations on a logical receiver $A$ and its current
purifier $R$:
\[
 |\Phi_Z\rangle=\sqrt{1-\epsilon}|00\rangle+
                   \sqrt\epsilon|11\rangle,\qquad
 |\Phi_X\rangle=(H_A\otimes I_R)|\Phi_Z\rangle,
 \quad 0<\epsilon<1/2.
\]
The logical labels denote Q-RAIF~C's neutral pair codes, with their
spectator sea retained. Here $R$ is the paid, resettable current purifier
in the third and fourth body species, available to joint neutral gates.
Old references outside this source quartet are denoted $R_{\rm old}$
and are not acted upon. The quartet is a resettable working source
alongside the series' protected unknown content $H$, not a replacement
for it; access to the current $R$ comes from those paid body species
and joint neutral gates, not from an old-purifier permission. Each cycle's inaccessible environment applies
the same pointer-aligned channel
\[
 \mathcal D_q(\rho)=(1-q)\rho+qZ_A\rho Z_A,
 \qquad 0<q<1/2,\qquad h(x)=-x\ln x-(1-x)\ln(1-x).
\]
No noise error label is available to the apparatus. Renewal returns the
complete quartet, including the current purifier.

\begin{proposition}[Pointer alignment and the complete renewal debit]
\label{prop:pointer-debit}
For $\rho_p=\mathcal D_q(|\Phi_p\rangle\langle\Phi_p|)$, return to its
specified pure quartet removes entropy $h_p=S(\rho_p)$, where
\begin{equation}\label{eq:pointer-debit}
 h_Z=h\left(\frac{1-\sqrt{1-16q(1-q)\epsilon(1-\epsilon)}}2\right)>0,
 \qquad h_X=h(q)>h_Z.
\end{equation}
The receiver-only debits are $0$ for $Z$ and
$h(\epsilon+q(1-2\epsilon))-h(\epsilon)$ for $X$.
Approximate quartet return subtracts its residual joint entropy from
$h_p$, with finite-dimensional entropy continuity controlling the error.
Program updates, sensors and usable controller registers enter the full
maintenance debit in Proposition~\ref{prop:entropyaccount} as well.
\end{proposition}
\begin{proof}
The mixture of $|\Phi_p\rangle$ and $Z_A|\Phi_p\rangle$ has nonzero
eigenvalues $(1\pm d_p)/2$, with
$d_p^2=1-4q(1-q)(1-|\langle\Phi_p|Z_A|\Phi_p\rangle|^2)$.
The overlaps are $1-2\epsilon$ and zero. Strict monotonicity of $h$ on
$(0,1/2)$ proves the contrast. Tracing $R$ gives the receiver debits.
The reference marginal is unchanged by the noise, so the difference
between each complete and marginal debit equals the receiver--purifier
mutual information lost to the noise.
\end{proof}

Pointer alignment lowers the complete renewal price without making it
zero. For $\epsilon=.1,q=.05$,
$h_Z=.0877516288$ and $h_X=.1985152433$ nats per cycle.
The marginal $Z$ state is unchanged, but its purifier correlations are not.
The loop must establish that its additional charged operations preserve
an actual complete-debit contrast $\bar s_X>\bar s_Z$.

\begin{conjecture}[Pointer-aligned selection with a sharp complete exhaust]
\label{conj:gravitating}
Fix $\nu>0$, $c_{\rm ex}>0$ and thermal input at $T_H=1/(2\pi)$ in the
running rapidity clock, in natural units. For this dephasing family there
are local models with one reused body during each $N$-cycle run, for
arbitrarily large $N$, restoring both preparations to the declared error
and realizing the following requirements jointly.

Renewal acts on the current resettable $A/R$ source, with joint control
of both pairs explicitly permitted. Independently held unknown content
$H$ and its inaccessible old reference $R_{\rm old}$ remain preserved
by the specified content-holding task throughout the loop. No reset of
$H$ or operation on $R_{\rm old}$ is granted by source renewal.
Compatibility of that task with sustained renewal and the response
controller is part of the conjecture.

The selected coupling supplies a complete chiral output carrying the
charged heat and entropy. Every retained usable register, sensor and
program update enters $s_{1:N}$, with $s_{1:N}/N\to\bar s>0$ for the
actual adaptive run. Initial usable purity and correlation capacity are
$o(N)$; there is no positive rate of uncharged purity or nonequilibrium
resource consumption. The excess above the complete-output floor
\begin{equation}
 \limsup_{N\to\infty}\left[
 \frac{2\pi Q_{1:N}-s_{1:N}}{N}
 -\frac{6\nu\bar s^2}{c_{\rm ex}}\right]
 \label{eq:sharp-exhaust}
\end{equation}
can be made arbitrarily small along the family. This is sharpness for
the whole adaptive output, including correlations and every atypical label.

The physical comparison has a declared block and decision clock.
Let $L_N$ be the sorting block length, $W_N\ge L_N$ the number of
cycles used for one reading, and $d_N\ge0$ the causal transport,
sorting and sensing delay, all in cycle units. The program is held
through each frame of $B_N=W_N+d_N$ cycles. Its outgoing observation
interval $I_j=[\eta_j^{\rm out},\eta_j^{\rm out}+W_N/\nu]$ is physically
assigned to that program, and its record is available only by the
frame's end, before the next decision. Require
$K_N=\lfloor N/B_N\rfloor-1\to\infty$, so each tested decision and
its following program are inside the run; when block sorting is used,
$L_N\to\infty$. For the block converter, start at a sorting boundary
and choose $W_N=m_NL_N$, $m_N\in\mathbb N$. Pad the output latency to $r_{\rm out}L_N$
and the subsequent readout latency to $r_{\rm read}L_N$, choosing
$r_i=\max\{1,\lceil\delta_i/L_N\rceil\}$ for the physical latencies
$\delta_i$ in cycle units, $i\in\{\mathrm{out},\mathrm{read}\}$; these may depend on $N$.
The window $I_j$ contains $m_N$ complete output blocks after the
$r_{\rm out}$ blocks; the final $r_{\rm read}$ blocks allow its record
to arrive before the next decision. Thus
$d_N=(r_{\rm out}+r_{\rm read})L_N$ aligns all frames and windows.
Buffers, padded waiting, operation during delay and the final flush
enter the complete entropy and energy account.

For fixed $Q_{\rm ref}>0$ the normalized reading is
\begin{equation}\label{eq:window-horizon-reading}
 \overline Y_j=
 \frac{S_h(\eta_j^{\rm out}+W_N/\nu)-S_h(\eta_j^{\rm out})}
      {2\pi Q_{\rm ref}W_N}.
\end{equation}
Whole-frame $X$ and $Z$ program interventions, made at the same
pre-program state, history and geometric data and averaged over the
declared actual-run pre-frame histories, give
$\mathbb E(\overline Y_X-\overline Y_Z)\ge\delta_Y>0$
uniformly as implementation and sharpness errors vanish.
This is a per-cycle-normalized window contrast, not a one-cycle switch
inside a pooled block.

The response changes the next program toward $Z$ when the measured
maintenance response is larger. To isolate this edge, let
$P_j^{\rm on/off}(Z\mid\mathcal F_j)$ be the next-program probabilities with
the physical response input present or blocked, respectively.
The paired intervention keeps the upstream history and executed commands
fixed up to that decision, including all non-response controller inputs.
With $\mathcal F_j$ drawn from the actual run before this tested input, require
\begin{equation}\label{eq:window-response-gain}
 G_N=\frac1{K_N}\sum_{j=1}^{K_N}
       \mathbb E\!\left[P_j^{\rm on}(Z\mid\mathcal F_j)
                       -P_j^{\rm off}(Z\mid\mathcal F_j)\right]\ge g_0>0.
\end{equation}
The denominator includes every declared decision opportunity, not only
selected informative or exploration windows. The response-channel
replacement used in the test and its fixed reference are specified as
part of the sensor. Its finite resolution, range, stress and entropy
are charged. Neither a counterfactual fixed-program contrast nor a
conditional gain on rare events substitutes for this operational test.
For binary $X/Z$, the same opportunity/history average obeys
$\langle P^{\rm off}(X)\rangle_N
=\langle P^{\rm on}(X)\rangle_N+G_N
\ge\langle P^{\rm on}(X)\rangle_N+g_0$.
These are blocked-input probabilities on matched actual-run histories,
not frequencies of a separately evolved off-history; if the on average
tends to zero, the blocked average remains at least $g_0$.
The actual adaptive output still satisfies Eq.~\eqref{eq:sharp-exhaust}.

Each model supplies a conserved complete stress tensor, the heat transport
and a timelike worldline in a positive dilaton domain. Its actual clock
has $0<z_{\min}\le d\tau/d\eta\le z_{\max}<\infty$ uniformly along the
family, with physical inlet temperature $T_\tau=1/(2\pi z)$.
Thermal preparation, reference stress and matching errors give a
sublinear accumulated remainder in the horizon account. Drives are
recovered, stored or routed away from the heat reading, or enter as
explicit nonthermal sources in a matched background comparison. Their
program-dependent residual remains in the contrast. Support work,
clock preparation, thermal banks, buffer stages, ground offsets and
final stored energy belong to the same ledger. Size, installed energy,
control capacity and the dilaton domain may grow with $N$; each run
keeps one body and a controller whose initial entropy and usable purity
and correlation capacity are $o(N)$.
Temperature in the rapidity frame, external exhaust capacity and
cycle frequency stay fixed.
\end{conjecture}

The conjecture keeps sustained renewal, sharp complete export and an
informative causal return in one dynamics. At fixed
$\nu,c_{\rm ex},T_H,\bar s$, the quadratic excess over bare Landauer
remains. Proposition~\ref{prop:export-deficit} excludes direct purifier
dumping from the sharp target. Arbitrary-unitary reservoir constructions
approach Landauer's bound~\cite[Section~6]{ReebWolf2014}, but do not
keep this exhaust capacity fixed.

\begin{proposition}[Finite-carrier sorting of the dephased source]
\label{prop:pointer-sort}
Q-RAIF~C's charge-sector sorter~\cite{QRC} applies to $L$ independently
realized copies of either $\rho_p$ under its finite-carrier hypotheses:
the source marginal is $\rho_p^{\otimes L}$, initially a product with
the complete thermal carrier; their external purifications are retained;
joint source operations act on both the receiver and its current $R$
pair, while identity on old references refers to systems outside that
quartet;
the source is neutral in the original zero-mode quartet sector; the
original neutral gates and paid loading, unloading and pipeline stages
are available; and weighted energy tails control the implementation
limit. It returns the source near its blank with entropy slack $o(L)$
and complete output Gibbs deficit $o(V)$, $V=L/\nu$, at
\[
 T_{*,p}=T_H+\frac{6\nu h_p}{\pi c_{\rm ex}},\qquad c_{\rm ex}=4.
\]
The same spectral argument applies to a supplied deterministic word
with product marginal $\bigotimes_{j=1}^L\rho_{p_j}$ and actual entropy
$s_L=\sum_jh_{p_j}$. The finite-bank target uses $s_L$ plus the same
subextensive margin. If $s_L/L\to s$, its temperature tends to
$T_H+6\nu s/(\pi c_{\rm ex})$, with no convergence rate required;
otherwise there need not be one limiting temperature. This is a block
converter, not an adaptive moving apparatus or a per-cycle heat-reading law.
\end{proposition}
The spectral proof is in Appendix~\ref{app:source-clock-support}.
It uses a supplied product marginal and controlled energy tails, not
a local reset certificate for an adaptive history.

This converter supports specified joint-source renewal, not correction
of unknown content against an inaccessible purifier. For example, if
$R$ were instead forbidden to all operations, the noisy $Z$ pair would
have negativity $(1-2q)\sqrt{\epsilon(1-\epsilon)}$, below its target
$\sqrt{\epsilon(1-\epsilon)}$. A channel on $A$ and independent thermal
ports cannot increase that negativity: partial transpose on $R$
commutes with the channel, and trace norm contracts on Hermitian
operators. It therefore cannot restore the entangled target.
Preserving the distinct held content $H/R_{\rm old}$ alongside this
converter requires the stated separate content-holding interface.

In a long fixed-program calibration window, suppose its full debit is
$\bar s_p$ and converter and horizon remainders are sublinear.
For $F(s)=s+6\nu s^2/c_{\rm ex}$ the matched mean reading at the floor is
\begin{equation}\label{eq:pointer-gain}
 \overline Y_p=\frac{F(\bar s_p)}{2\pi Q_{\rm ref}}+o(1),\qquad
 \Delta\overline Y=\frac{F(\bar s_X)-F(\bar s_Z)}{2\pi Q_{\rm ref}}+o(1).
\end{equation}
For the quartet alone, $\nu=1,c_{\rm ex}=4$ give
$F(h_Z)=.0993021514$, $F(h_X)=.2576276961$ and
$\Delta\overline Y=.02519829/Q_{\rm ref}$.
Here $\overline Y_p$ is the normalized window reading
\eqref{eq:window-horizon-reading}, under the stated fixed-program
calibration assumptions, not a consequence of the adaptive sharpness
clause. A single changed preparation within a pooled $L$-cycle block
changes its mean debit by $O(L^{-1})$ and its normalized heat by
$O(L^{-1})$; it does not produce this whole-window contrast.

Separate fixed-program sharpness does not prove the whole adaptive
target. If time fractions $a_p$ are occupied by long individually sharp
phases with debits $\bar s_p$, their combined heat has the excess
\begin{equation}\label{eq:adaptive-jensen-gap}
 \sum_pa_pF(\bar s_p)-F\left(\sum_pa_p\bar s_p\right)
 =\frac{6\nu}{c_{\rm ex}}\operatorname{Var}_a(\bar s_p).
\end{equation}
This is positive for nonzero fractions and unequal debits. Joint block
sorting can remove that energy gap at the channel level, but the present
proof supplies no local program contrast or causal return after that
sorting. Correlations with retained program history also enter the
slack in Eq.~\eqref{eq:conditional-export-ledger}. These are joint
obstacles to Eq.~\eqref{eq:sharp-exhaust}. Sparse exploration with mean $X$-window share $a_N\to0$ removes this
phase-variance term but settles neither window gain. For independent
noisy window tests with a fixed nonzero information per trial, vanishing
estimation error requires the expected test count $K_Na_N\to\infty$. If window accuracy improves
with $N$, the requirement concerns the accumulated physical detection
information instead; counting windows alone is insufficient.
Let $b_N$ be the frequency of response-sensitive opportunities under
the same opportunity/history measure as $G_N$. If on/off kernels
coincide elsewhere, $|G_N|\le b_N$. When updates occur only in $X$
frames, $b_N\le a_N$ with $a_N$ evaluated over the same tested
opportunities and actual-run histories. For equal-length windows,
$a_N$ is the expected $X$ share of observed cycles, not the fraction
of the full run covered by $X$ observations. Hence $a_N\to0$ violates
Eq.~\eqref{eq:window-response-gain} despite fixed conditional contrast.
Exploration and update duty are not
identical in general: continuous monitoring during $Z$ windows could
keep $b_N$ nonvanishing. Its compatibility with charged sensing,
whole-output sharpness and a causal return is unproved, not excluded.

Under isotropic depolarization the two noisy quartet spectra and full
debits instead coincide. Their separately entropy-optimal matched heat
outputs have zero asymptotic program contrast; unsorted source contrast
can become recovered work. This excludes that particular sharp-heat
mechanism, not different full control debits or every causal realization
using depolarization. A fixed initial record also obeys
$I(X:Y_{1:N}|Z)\le H(X|Z)$ for charged baseline $Z$; sublinear initial
record entropy cannot supply a positive new-information rate indefinitely.

The exact marginal-horizon balance gives
$\overline Y_j=Q_{I_j}/(W_NQ_{\rm ref})$ when the entire matched
source is charged heat. Drive, other and reference-source integrals
add $[E_{{\rm drive},I_j}+E_{{\rm other},I_j}
+E_{{\rm ref},I_j}]/(W_NQ_{\rm ref})$ in the same normalization. A common background cancels in
program comparisons only with a justified common clock and transport;
program-dependent residuals remain. Table~\ref{tab:scope-ledger} separates the available finite readouts
from this joint inlet, renewal, sensor and outlet task. In particular,
Q-RAIF~C's massless sorter and massive-outside full-angle gate are
different layouts~\cite{QRC}.

The sorting-based route fails without source entropy concentration or
controlled energy tails. Any joint realization must also avoid positive
asymptotic output deficit or slack, loss of program contrast after drive
and sensing charges, and failure of causal routing, temperature matching
or the operating worldline.

Fresh, independently discarded quadrature drives have a separate
accuracy cost~\cite{Liu2026QRAIFB}: at error $\epsilon_{\rm drv}$ their
input energy is at least $\omega_h[\pi/(4\sqrt{\epsilon_{\rm drv}})+O(1)]$,
where $\omega_h^{-1}=\int|f|^2\omega^{-1}d\omega$.
Reused drives retain correlations and correlated failures; recovery is
not established by fresh-drive accuracy. At $\epsilon_{\rm drv}=.01$
the finite bound is $6.74238\omega_h$. Choose the source parameter
$\epsilon=.01$ independently of $\epsilon_{\rm drv}$.
With $q=.05$, $\nu=1$ and $c_{\rm ex}=4$, the complete quartet debits are
$(h_Z,h_X)=(.01370657,.19851524)$. Their heat floors
$F(h_p)/(2\pi)$ are $(.00222632,.04100272)$.
With the additional benchmark choice $\omega_h=1$, the drive/heat
ratios are about $(3028,164)$; unrecovered fractions
$(.0003302,.0060813)$ already equal one corresponding heat unit.
At the main example's $\epsilon=.1$, the $Z$ floor is $.01580443$
and that ratio is about $427$, while the $X$ values are unchanged.
These quartet-only benchmarks exclude additional charged controller
debits and do not identify the drive with the quartet's full renewal
device. Gate spectrum, duration and debit set the ratio; no universal
$\omega_h\ge\nu$ or universal drive/heat multiplier follows. Outgoing coherent
drive energy uses the nonthermal ledger in Eq.~\eqref{eq:nonthermal-exhaust}.
Recoverable energy, stored energy and exported heat are separate entries;
the same outgoing energy is never charged twice.

The exact moving-clock conversion is
\begin{equation}\label{eq:loop-moving-clock}
 z=\frac{Ly}{R},\qquad
 y=\sqrt{R_b^2+L^2\dot\Phi_b^2}-L\dot\Phi_b,
 \quad R_b^2=\Phi_b^2-R^2,
\end{equation}
with proper-time dot. The tracking-clock conversion is derived in
Appendix~\ref{app:source-clock-support}, Eq.~\eqref{eq:tracking-clock-exact}.
Thus $z=L\rho$ is a stationary limit. The physical input needs
$zT_\tau=T_H$, and support power is
$(L^{-1}-K)\dot\Phi_b/\kappa$ in
Eq.~\eqref{eq:boundary-work-balance}, including its actual sign and source.
Prescribing the trajectory or transforming temperature constructs
neither support nor inlet. When the monitored source carries heat plus
a sublinear reference remainder, its source--clock sector has
$R=R_0+\kappa Q_{1:N}+o(N)$ at a positive bounded mean heat; an exterior
operating domain must therefore reach $\Phi_b>R=\Omega(N)$.
With bounded $z$, an isolated reserve paying all per-cycle losses also
requires initial energy $\Omega(N)$; replenished charged work is an
alternative. The family must report its upper as well as lower size, installed-energy,
clock-domain and compiler scalings. A stronger resource target would be
$O(N)$ installed energy and dilaton range, with polynomial control
capacity; this is an optional benchmark, not a hypothesis of
Conjecture~\ref{conj:gravitating}. Faster growth would miss that benchmark
without refuting the stated conjecture. Fixed-gap and fixed-cutoff
examples alone give no uniform growing-run certificate.

This loop would connect selection to its source and backreaction.
Its finite pointer algebra still needs a map to the modular-stable
geometric algebra, and Conjecture~\ref{conj:main} still asks for the
origin of the gravitational action.


\subsection{The two routes and a common source principle}
\label{sec:source-principle}

The finite return demonstrates that a law supplied to a model can feed
back into the physical operations that select access. It does not explain
that law. Modular reconstruction supplies a constraint on the law, but
does not manufacture the apparatus. Their union requires a further
mechanism: the same physical history must determine both the accessible
structure and its complete source response.

\begin{conjecture}[Common-action selection and geometric consistency]
\label{conj:common-action-selection}
There is a local matter--controller theory in which readable histories
form and maintain a nontrivial selected geometric structure, and whose
causal effective action uses the same prepared fields, support and clock
in its source and response. In a low-energy sector with a common cone,
its selected ball data satisfy the geometric consistency condition of
Consciousness~II. Its first-order response is the corresponding Einstein
response with one source normalization; its higher-order response includes
the interaction, contact and clock work that produced the selection.
Selection changes this response through physical couplings and sources,
while a change of description at fixed matter history leaves the source
unchanged. The sustainable family of Conjecture~\ref{conj:gravitating}
is a finite operating sector of this theory, with a response-to-operation
gain that survives the growing-window limit.
\end{conjecture}

The conjecture has tests beyond assigning gravity to a massive apparatus.
Removing the history's access to a command must remove the
record-mediated component of the source, while holding the commanded
interaction fixed and merely discarding a log must preserve it.
Source identification must survive switched physical references;
logical noise records alone fail this test in
Proposition~\ref{prop:weighted-source-ambiguity}.
The same calibrated source must then pass the all-ball consistency
test with fixed initial or boundary data. A law that reproduces only
the JT dilaton or one binary return has not passed the higher-dimensional
test. At second order the source, geometry and changed preparation must
be varied together; entropy Hessians do not become SLD or spacetime
metrics by name.

Q-RAIF~C supplies a local common action with dynamical support, memory
and clock fields and a specified positive finite canonical ordering
\cite{QRC}. This permits a concrete attack on the missing source:
freeze the texture only as a controlled limit, quantize the retained
material and controls together, and compare the resulting local gates
with the actual acquired-history Hamiltonian. Direct substitution fails
already because dynamical support and clock introduce kinetic cross
terms and forces absent from a fixed-profile gate. Moreover a spatial
canonical cutoff preserves its positive form but does not by itself
provide a covariant renormalized Ward identity. Classical on-shell
conservation therefore cannot certify the complete quantum source.
The coupled tolerance above closes the trajectory/reader comparison in
the retained mean-action branch; it cannot remove these two operator
and continuum obstacles.

For the real-quartet carrier, Proposition~\ref{prop:covariant-quartet-bulk}
supplies that Ward identity and the actual full-mode bulk state. Its
projective and Weyl conversion is varied with the clock and director;
adding a complete vacuum effective action to the retained canonical
stress would double count the quantum response. This bulk construction
does not yet transport the entire reset protocol through dynamical walls.

The finite receiver-renewal family of Q-RAIF~C holds clock mean energy
and norm fixed while its space, time and paid Gibbs bank grow. It is a
real finite update result, with old references retained. It renews a
receiver rather than the entire interacting material. In the 31-site
full-program diagnostic, the receiver occupation $.157712$ exceeds
the target $.13$ by $.027712$; this is not a full-program state error.
The independently reported completion probability is $.857524$.
An additional complete-ring construction renews the same body to fixed
per-cycle tolerance through four thermal outlet species, retaining old
correlations and restoring the original mixed receiver and its current
purifier by local neutral pumps~\cite{QRC}. At the same incoming
history, changing the maintained logical axis changes the stress
emitted by that reset, with a strictly positive finite-error bound;
no separate program-tagged emitter is added. Its prescribed drives are
paid, and its energy-constrained storage stays bounded at fixed
parameters. Its finite smooth schedule can also be driven by a work-clock
field whose source includes the old incoming correlations~\cite[Section
on a spatial work clock]{QRC}. That result has static endpoints and a
supplied frame. The bulk covariant construction and this fixed-frame
schedule have not yet been joined with those endpoints and the sharp
excess-heat limit required here.
The finite-program diagnostic remains evidence against combining successful
components without a common state and dynamics. The conjectured common theory
would be refuted by an unavoidable extra extensive source or renewal
price, by a persistent mismatch of its selected modular data, or by a
response gain that vanishes under the stated fixed resources. These
are different failure modes, not a general impossibility claim.

\section{Relation to Existing Approaches}
\label{sec:relation}

AdS/CFT~\cite{Maldacena1999} and its area--entropy relation~\cite{RyuTakayanagi2006} supply the dictionary used in Theorem~\ref{thm:linearized}; it is not derived here or granted to arbitrary selected algebras. MERA~\cite{Vidal2008,Swingle2012} implements coarse-graining through disentanglers and isometries, with network depth resembling a radial AdS direction. Different networks can describe the same boundary data. Relating either construction to physical selection requires the explicit algebra-to-geometry map.

Jacobson's horizon Clausius argument~\cite{Jacobson1995} shares the entropy-based approach to gravitational dynamics. Section~\ref{sec:source} addresses its source-side input: which selecting processes produce the stress whose response is computed. Loop quantum gravity~\cite{Rovelli2004,Thiemann2007} instead quantizes the metric with background independence. Its relation to a path of selected algebras would require a representation map from spin-network states. These connections motivate the proposal without supplying its missing selection law.

\section{Conclusion}
\label{sec:conclusion}

Physical selection reaches gravity through the selecting interaction's source. Calibration changes activation rates; selected projections and learned protection change stress at fixed schedule. Logical noise records alone can miss that source. Prepared references or signed field readings remove specific ambiguities, with their resources charged.

Within matched JT data, a finite return retains the quantum reader, old reference, record spring, clock energy, mechanics, radiation and recoil in one mean action. Acquisition and feedback are separately necessary for its source contrast, and joint error bounds preserve the test under finite perturbations. The spatial construction includes support and synchronization, with a nonzero source--dilaton--sensor path and a subsequent complete material-force response. The latter also contains ordinary material and clock effects; it is not entirely gravitational. The full-mode quartet bulk supplies a covariant, actual-state source with clock and director work, rather than replacing the retained state by a vacuum.

Continued maintenance is a different result. The entropy of the complete retained system, consumed purity, correlations and conserved charges fixes the export debit. A thermal chiral channel of fixed capacity imposes its quadratic rate price. Matched horizon transport records that price; the actual thermal-source calculation also retains the moving clock's reference stress. Renewal accuracy alone does not establish sharp exhaust or indefinite operation.

The modular route constrains the law through a supplied geometric family and
dictionary; the source route computes within an assigned law. The common-action
and sustainable-loop conjectures ask for formed access, geometric consistency,
complete causal sources and persistent return gain in one dynamics. Pointer
alignment gives a complete-debit contrast; finite-carrier sorting can approach
its channel floor. Individually sharp program phases still leave a positive
variance price at nonzero duty fractions, so joint sharpness with an informative
return remains open. Endpoint and gate matching, moving support, sustained
thermal renewal and the origin of geometric selection remain construction tasks.

\appendix

\section{Modular definitions and candidate checks}
\label{app:modular-details}

\subsection{Preliminaries}
\label{sec:modular-prelim}

Standard references are Takesaki~\cite{Takesaki2003}, Bratteli--Robinson~\cite{BratteliRobinson1997} and Haag~\cite{Haag1996}.

\begin{definition}
A \textbf{von Neumann algebra} $\mathcal{M}$ on $\mathcal{H}$ is a $*$-subalgebra of $\mathcal{B}(\mathcal{H})$ containing the identity and closed in the weak operator topology; equivalently $\mathcal{M} = \mathcal{M}''$, where $\mathcal{M}'$ is the commutant. It is a \textbf{Type~III$_1$ factor} if $\mathcal{M} \cap \mathcal{M}' = \mathbb{C} \mathbf{1}$, it admits no semifinite normal trace, and its Connes spectrum is $\mathbb{R}_{\geq 0}$.
\end{definition}

Under standard conditions, local algebras of quantum field theories are Type~III$_1$ factors~\cite{Haag1996, Yngvason2005}.

A vector $|\Omega\rangle$ is \textbf{cyclic} for $\mathcal{M}$ if $\mathcal{M}|\Omega\rangle$ is dense, and \textbf{separating} if $M|\Omega\rangle = 0$ implies $M = 0$. For cyclic and separating $|\Omega\rangle$ the antilinear map $S_\Omega : M|\Omega\rangle \mapsto M^*|\Omega\rangle$ is closable, and its closure has the polar decomposition $S_\Omega = J_\Omega \Delta_\Omega^{1/2}$ with modular conjugation $J_\Omega$ and modular operator $\Delta_\Omega$.

\begin{theorem}[Tomita--Takesaki {\cite{Takesaki2003}}]
\label{thm:TT}
With $\mathcal M$, $|\Omega\rangle$, $J_\Omega$, $\Delta_\Omega$ as above, $J_\Omega \mathcal{M} J_\Omega = \mathcal{M}'$, and the modular group $\sigma_t^\Omega(M) := \Delta_\Omega^{it} M \Delta_\Omega^{-it}$ satisfies $\sigma_t^\Omega(\mathcal{M}) = \mathcal{M}$ for all real $t$.
\end{theorem}

A state $\omega$ is KMS at parameter $\beta$ in the modular-time convention used here
if for all $A,B\in\mathcal M$ there is a function $F_{A,B}$ analytic in
$0<\operatorname{Im}z<\beta$ and continuous on the closure, with
$F_{A,B}(t)=\omega(\alpha_t(A)B)$ and
$F_{A,B}(t+i\beta)=\omega(B\alpha_t(A))$. This ordering gives $\beta=1$ for
$\sigma_t^\omega(A)=\rho^{it}A\rho^{-it}$ in a density-matrix representation;
the usual future physical Hamiltonian time has the opposite modular orientation.

\begin{theorem}[Takesaki {\cite{Takesaki2003}}]
\label{thm:KMS-unique}
A faithful normal state $\omega$ on $\mathcal{M}$ is KMS at $\beta = 1$ for $\sigma_t^\omega$, and $\sigma_t^\omega$ is the unique one-parameter automorphism group with this property.
\end{theorem}

The modular flow is thus the unique time evolution for which the given state looks thermal. In the Rindler wedge it is the Lorentz boost, and the KMS condition at $\beta = 2\pi$ gives the Unruh temperature $T_U = a/(2\pi)$.

\begin{definition}[{\cite{Wiesbrock1993}}]
An inclusion $\mathcal{N} \subset \mathcal{M}$ with common cyclic and separating vector $|\Omega\rangle$ is a \textbf{half-sided modular inclusion} (HSMI) if $\sigma_t^{\mathcal{M}}(\mathcal{N}) \subset \mathcal{N}$ for all $t \leq 0$.
\end{definition}

\begin{theorem}[Wiesbrock {\cite{Wiesbrock1993}}]
\label{thm:Wiesbrock}
For an HSMI in the stated $t\le0$ convention, the closure of
$P=(\log\Delta_{\mathcal N}-\log\Delta_{\mathcal M})/(2\pi)$ is positive. The canonical
unitary group is $U(a)=e^{iaP}$, fixes $|\Omega\rangle$, and obeys
$\mathcal N=U(1)\mathcal M U(1)^*$ and
$\Delta_{\mathcal M}^{it}U(a)\Delta_{\mathcal M}^{-it}=U(e^{-2\pi t}a)$~\cite[Thm.~2.1]{ArakiZsido2005}.
\end{theorem}

The translation from $\mathcal M$ to $\mathcal N$ is therefore encoded algebraically, without background geometry. This is the standard tool for extracting translation structure from modular data.


\subsection{Example 1: Rindler wedge}
\label{sec:rindler}

Consider a free massless scalar field in $(1+3)$-dimensional Minkowski spacetime with Fock space $\mathcal{H}_U = \mathcal{F}(L^2(\mathbb{R}^3))$ and vacuum $|\Omega\rangle$. Let $\mathcal{W}_R = \{x > |t|\}$ be the right Rindler wedge and $\mathcal{R}(\mathcal{W}_R) = \{e^{i\phi(f)} : \operatorname{supp} f \subset \mathcal{W}_R\}''$. For this check we take $\mathcal M_U=\mathcal R(\mathcal W_R)$, which makes the ambient modular action well defined without performing a nontrivial selection.

The Reeh--Schlieder theorem~\cite{Haag1996} makes $|\Omega\rangle$ cyclic and separating for $\mathcal{R}(\mathcal{W}_R)$, so (U1) holds. By the Bisognano--Wichmann theorem~\cite{BisognanoWichmann1975, BisognanoWichmann1976}, $\Delta_\Omega = e^{-2\pi K}$ with $K$ the global boost generator in the $x$-direction. The modular action is $\sigma_t^\Omega(X)=e^{-2\pi itK}Xe^{2\pi itK}$; its spacetime action preserves the wedge, giving $\sigma_t^\Omega(\mathcal{R}(\mathcal{W}_R)) = \mathcal{R}(\mathcal{W}_R)$, which is (U2). No decoherence semigroup is specified, so (U3) is not verified. Haag duality, $\mathcal R(\mathcal W_R)'=\mathcal R(\mathcal W_L)$, does not imply that every proper extension of the wedge algebra contains the opposite wedge algebra, nor any maximality statement.

On the $t=0$ Cauchy slice the full modular generator is
\begin{equation}
K_{\mathrm{mod}} = -\ln \Delta_\Omega = 2\pi K,
\qquad K=\int_{\mathbb R^3}x\,T_{00}(0,x,y,z)\,d^3x .
\end{equation}
Formally this is $2\pi(K_R-K_L)$, with $K_R=\int_{x>0}xT_{00}\,d^3x$ and $K_L=\int_{x<0}(-x)T_{00}\,d^3x$. The right-wedge expression $2\pi K_R$ alone is not $-\log\Delta_\Omega$: in a factorizing regulator it is the right reduced-state modular Hamiltonian, up to its normalization constant. In the Type III theory the separate one-sided charges require an appropriate quadratic-form or regulator interpretation. For a local unitary excitation in the right wedge, the opposite-wedge expectation is unchanged, so the change of the full boost charge equals the right-sided change when those differences exist. This does not identify the two operators. The example reproduces the Bisognano--Wichmann geometric modular action; it does not show that (U1)--(U3) select the wedge algebra.

\subsection{Example 2: Qubit chain with dephasing}
\label{sec:qubit}

Consider $n$ qubits with Hamiltonian $\hat{H} = -J \sum_{i} \sigma_z^{(i)} \sigma_z^{(i+1)} - h \sum_{i} \sigma_x^{(i)}$, $J > 0$, coupled to an environment through
\begin{equation}
\label{eq:lindblad}
\mathcal{T}_t(\rho) = e^{t\mathcal{L}}(\rho), \qquad \mathcal{L}(\rho) = -i[\hat{H}, \rho] + \gamma \sum_i \left(\sigma_z^{(i)} \rho\, \sigma_z^{(i)} - \rho\right),
\end{equation}
with Heisenberg dual $\mathcal E_t:=\mathcal T_t^*$ and global state $\omega = Z^{-1} e^{-\beta \hat{H}}$. For an exact check we take $h=0$; for $h\ne0$ the diagonal algebra is stable only approximately in a strong-dephasing regime.

For $h=0$ the diagonal observables are fixed by the Heisenberg semigroup, and the candidate is the maximal abelian algebra
\begin{equation}
\mathcal{A}_{\mathrm{geo},\mathbf{c}} = \{f(\sigma_z^{(1)}, \ldots, \sigma_z^{(n)})\}'' = \mathcal{A}_{\mathrm{diag}}.
\end{equation}
The thermal state is positive definite at finite $\beta$, so $\omega(|s\rangle\langle s|) = \langle s | \rho_{\mathrm{th}} | s \rangle > 0$ for every computational basis state, which gives (U1); diagonal elements of a positive definite matrix are positive in any basis. The ambient modular flow is $\sigma_t^{\omega,U}(X) = \rho_{\mathrm{th}}^{it} X \rho_{\mathrm{th}}^{-it}$; for $h=0$, $\rho_{\mathrm{th}}$ is diagonal and commutes with $\mathcal A_{\mathrm{diag}}$, giving (U2). For $h\ne0$ the conjugation generally leaves $\mathcal A_{\mathrm{diag}}$. Under the dephasing dual, off-diagonal matrix units decay as $e^{-2\gamma\, d(s,s')\, t}\, |s\rangle\langle s'|$ with Hamming distance $d$; for $h=0$ every diagonal observable is fixed, so (U3) holds and $\mathcal A_{\mathrm{diag}}\subset\operatorname{Fix}(\mathcal E)$. This does not prove maximal setwise invariance, since the full matrix algebra is always mapped into itself.

The modular flow acts trivially on the abelian pointer algebra, and the entropy of a block $A$ within $\mathcal{A}_{\mathrm{diag}}$ is the Shannon entropy of the marginal distribution $p(s_A)$. The example checks one modular- and semigroup-stable candidate, which under exact fixation is the pointer algebra; it neither verifies the Type~III regime nor selects the candidate uniquely by (U1)--(U3). Selection of another type is available here: for $h=0$ the fixed algebra is the commutant of the $\sigma_z^{(i)}$, and the window theorem of Paper~B shows that every sufficiently long observation window selects it as the unique maximal stable algebra~\cite{Liu2026PaperB}.


\section{Source identification, carriers and profiles}
\label{app:source-details}

\subsection{Continuing calibration and identification of the source}
\label{sec:source-identification}

Source control need not interrupt stored quantum content. With an independently
prepared vacuum mode and full output collection, Q-RAIF~B's corrected parity
instrument has branches
\[
 \mathcal I_\pm(\rho)=\tfrac12(1\pm e^{-2k(1-\alpha a)^2})\rho.
\]
It yields the same calibration statistic without preparing a new body state.
The coherent controlled-parity identity
$R e^{-imZ_LQ_f}R=I_L\otimes e^{-imQ_f}$ goes further: it protects
the body's marginal for an arbitrary mode state by transferring the interaction
to the field alone~\cite{Liu2026QRAIFB}. For a centered independent input the
coupling-added stress is still $m^2F_f$. The meter can be returned coherently;
retaining a calibration outcome is a separate read and memory operation.
Partial collection loses coherence during the inquiry itself. Brain Anchor
prices that loss, bounds continuing retuning under drift and gives an exact
posterior criterion for when an extra inquiry pays for its source and apparatus
price~\cite{BrainAnchor}. No geometric response or entropy export is needed
for these preservation identities.

The inquiry and maintenance choices can also share one running body.
Brain Anchor tracks both the interface posterior $b_x$ and the
coherence-weighted mass $c_x$ in a finite hidden-state model. Equal $b_x$
and even equal reduced body states can require different next commands,
because past loss has different correlations with the interface.
Q-RAIF~B realizes collection by finite resonant contacts, including successive
contacts with a finite packet. In the matched field realization, the source
at the next encounter is still $m_j^2F_f$, irrespective of the accessible
correction, at the interaction output before collection. A collector can
redistribute that energy, so this is not its full downstream stress.
For a history-dependent policy the coefficient is computed
from that policy's hidden-path weights, not from a new independent gain draw
at every step. The four-encounter joint example gives
$\sum_j\mathbb E m_j^2=.588965$, against $.640000$ with permitted
maintenance but no retained calibration reads. These are source accounts;
the supplied objective's prices are not an entropy-export measurement.
Using these sources in the JT response also requires their profiles and
emission times and any collector contribution. The continuous collection
route below specifies a free-field remainder, but does not supply the wall,
thermal exhaust or entire sensing apparatus of the response model.

\subsubsection{Continuous collection changes the outgoing profile}
\label{sec:band-source}

Q-RAIF~B also gives an excitation-preserving, positive-frequency collection
model~\cite{Liu2026QRAIFB}. A retained oscillator of frequency $\Omega$
couples to a continuum $\omega\in[\Omega-B,\Omega+B]$, with real even
form factor $g(\omega-\Omega)$. A prescribed finite source drive prepares
the matched packet. The one-particle propagator $S$ then determines both
the collected fraction and the outgoing wave. Its finite waveform certificate
allows clock, arrival and switching errors; it is not a white-noise limit.
The field interaction is band-smeared, so its spatial tails are part of the
model. Parameters are fixed within each encounter.

Let $z_{\rm out}(\omega)$ denote the asymptotic free outgoing amplitude for a
unit source drive, at the same null-time origin as the incoming packet,
and let the body-conditioned
amplitudes be $\pm m\beta z_{\rm out}$. With the current normalization of
Eq.~\eqref{eq:stress-id}, put
\begin{equation}
 \chi(u)=-i\!\int_{\Omega-B}^{\Omega+B}
 \sqrt{\frac{\omega}{4\pi}}z_{\rm out}(\omega)e^{-i\omega u}d\omega,
 \quad F_{\rm out}(u)=4m^2\beta^2[\operatorname{Re}\chi(u)]^2.
 \label{eq:band-out-stress}
\end{equation}
For independent vacuum preparation, this is the outgoing free-field stress
for every logical input. Tracing the body gives a mixture of opposite
coherent amplitudes, whose normal-ordered squared currents agree.
Accessible parity correction does not change this marginal stress.
For a capture pulse ending at $T$, the rotating-frame amplitude is
$z_{\rm out}(\omega)=e^{i(\omega-\Omega)T}
 [P_{\rm field}S(0,w)](\omega)$, with the prescribed source carrier phase
included. Removing the final free propagation phase fixes the common time
origin; it must not be counted again when placing the source in the JT model.
The input occupation is $k_{\rm eff}m^2=\beta^2m^2\|w\|^2$.
For real envelopes and an even band, the source, receiver and free remainder
energies in units $\hbar=1$ are, respectively,
\begin{equation}
 \Omega k_{\rm eff}m^2,\qquad
 \Omega\xi k_{\rm eff}m^2,\qquad
 E_0^{\rm out}=\Omega(1-\xi)k_{\rm eff}m^2.
 \label{eq:band-source-energy}
\end{equation}
Capture moves energy into the receiver. A later vacuum swap exports it to
a separate exhaust. The net prescribed capture-switch work is zero for
this symmetric packet family, while source preparation supplies the first
energy in Eq.~\eqref{eq:band-source-energy}. The receiver, exhaust, controls
and supports do not disappear from a gravitational account.

Frequency regularity matters for using this source in the JT moment formulas.
A hard flat band generally gives $F_{\rm out}(u)$ tails of order $u^{-2}$,
so finite energy alone does not ensure a finite second moment. Instead choose
$g(\delta)=\sqrt{J_0}[1-(\delta/B)^2]$ inside the band and zero outside.
For a finite source envelope and finite collection time,
$z_{\rm out}(\omega)$ is $g$ times a smooth spectral function. Thus
$\sqrt\omega z_{\rm out}\in H^1_0$ and Plancherel gives
\begin{equation}
 E_2^{\rm out}=\int u^2F_{\rm out}(u)du
 =m^2\beta^2\!\int
 \left|\partial_\omega(\sqrt\omega z_{\rm out})\right|^2d\omega<\infty.
 \label{eq:band-second-moment}
\end{equation}
$E_1^{\rm out}$ is finite by Cauchy--Schwarz. This specifies a legitimate
free-field source for the response dictionary, rather than reusing the
uncaptured packet unchanged. It does not determine the total apparatus stress
or assert compact temporal support. The JT background, boundary access and
matter normalization remain the declared inputs of the response model.

In the worked collection setting $\Omega=20$, $B=5$, $T=16$ and
$J_0=1/(2\pi)$, a fixed 8721-bit waveform table and stated clock tolerances
certify $1-\xi<4.43\times10^{-7}$ and a relative preparation-strength
uncertainty below $.001$. This is enough to keep the collection-induced loss
below $.001$ in the preceding 4096-per-setting calibration example. The bound
does not include parity-gate or detector errors. Repeated encounters require
independent vacuum fields and exhausts, or a separately proved return/reset
bound on one reused field. A fixed-capacity thermal cycle and the source--sensor
return of Section~\ref{sec:gap} are not supplied by this finite-bank construction.

With several Gaussian field species, however, logical calibration alone may
not identify the actual source. Let the species have the same switch profile
and independent centered inputs, with quadrature covariance $kW$ in their
species coordinates. Put $G=(\mathbf g_1,\mathbf g_2)$ and
$\mathbf v(a)=\mathbf g_1-a\mathbf g_2$. The logical coherence and
coupling-added stress are respectively
\begin{equation}
 r_W(a)=e^{-2k\mathbf v(a)^TW\mathbf v(a)},\qquad
 T_{\rm add}(u)=\|\mathbf v(a)\|^2F_f(u).
 \label{eq:weighted-noise-source}
\end{equation}
The first measures $D_W=G^TWG$; the second needs $C=G^TG$.
The incoming stress of the reference preparation is an additional source.

\begin{proposition}[Protection records do not always identify stress]
\label{prop:weighted-source-ambiguity}
Take $W=\operatorname{diag}(1,4)$ and the two coupling families
\[
 \mathbf g_1=(1,0),\quad\mathbf g_2=(1,1),\qquad
 \widetilde{\mathbf g}_1=(0,1/2),\quad
 \widetilde{\mathbf g}_2=(2,1/2).
\]
Every command gives the same logical channel, with weighted residual
$1-2a+5a^2$. Their stress coefficients are
$1-2a+2a^2$ and $1/4-a/2+17a^2/4$.
Thus no finite adaptive experiment using only these fresh-port logical
channels distinguishes them. The noise-minimizing command is $a=1/5$
for both, but their source-minimizing commands are $1/2$ and $1/17$.
\end{proposition}
\begin{proof}
Insert each vector pair in Eq.~\eqref{eq:weighted-noise-source}.
Equality of conditional channels at every command gives equality of every
adaptive logical record law by induction, also with retained ancillas.
Minimizing each positive quadratic gives the stated commands.
\end{proof}
At the noise optimum their added stress coefficients are $17/25$ and
$8/25$. This is an ambiguity of the specified logical access, not an
inability to observe stress through another instrument. It persists with
the same incoming preparation and does not come from a changed background.

\begin{proposition}[Switched references identify the physical source]
\label{prop:reference-source-identification}
In the same two-species family, suppose the apparatus can independently
prepare known centered Gaussian references with covariances $kW_1,kW_2$,
where $W_1=\operatorname{diag}(1,4)$ and
$W_2=\operatorname{diag}(4,1)$. At each reference measure the logical
coherence means at $0,+h,-h$, $h>0$. They determine the physical source
Gram matrix exactly:
\begin{align}
 D_{j,11}&=-\frac{\log r_j(0)}{2k},&
 D_{j,12}&=\frac{\log r_j(h)-\log r_j(-h)}{8kh},\nonumber\\
 D_{j,22}&=-\frac{\log r_j(h)+\log r_j(-h)-2\log r_j(0)}{4kh^2},&
 C&=\frac{D_1+D_2}{5}.
 \label{eq:reference-source-reconstruction}
\end{align}
More generally, known reference weights with
$I=\sum_jt_jW_j$ suffice, with $C=\sum_jt_jD_j$.

For independent $\pm1$ trials, $N$ at each of the six settings, suppose
all true means are at least $r_{\min}>0$. Let
$\eta=\sqrt{2\log(12/\zeta)/N}\le r_{\min}/2$, clip sample means to
$[r_{\min}/2,1]$ and use the same formulas. With probability at least
$1-\zeta$, uniformly on $|a|\le A$,
\begin{equation}
 |\widehat V(a)-\|\mathbf v(a)\|^2|
 \le\Delta_C:=\frac{2\eta}{5kr_{\min}}
            \left(1+\frac A h+\frac{2A^2}{h^2}\right).
 \label{eq:reference-source-error}
\end{equation}
Minimizing $\widehat V$ on any supplied finite command grid incurs
actual excess source coefficient at most $2\Delta_C$ over its best point.
\end{proposition}
\begin{proof}
The log characteristic function is the quadratic form $-2kD_j$.
Its three coefficients give Eq.~\eqref{eq:reference-source-reconstruction};
$W_1+W_2=5I$ gives $C$ and proves the general sufficient condition.
The union Hoeffding bound covers all six means. Each log error is at most
$2\eta/r_{\min}$, so the errors of $C_{11},C_{12},C_{22}$ are at most
$2\eta/(5kr_{\min})$, $\eta/(5khr_{\min})$ and
$4\eta/(5kh^2r_{\min})$. Insert them in
$V=C_{11}-2aC_{12}+a^2C_{22}$ to prove the uniform bound.
Two applications of that bound compare the chosen grid point with the best one.
\end{proof}

These are $6N$ new preparations, couplings and reads, not six free exact
means. The reference covariances, their switches and their incoming energies
have to be calibrated and charged. Separate logical test preparations are
permitted here; applying these references to held content would also require
the protection operations of Q-RAIF~B and their collection account.
Equation~\eqref{eq:reference-source-reconstruction} identifies the
coupling-added source under the declared profile and species architecture,
not the full stress of an unknown apparatus. There is another route when
the apparatus has an oriented field readout rather than only logical records.

\begin{proposition}[Signed field calibration with held content]
\label{prop:signed-source-calibration}
For $d$ independent input modes put $Q_s=b_s+b_s^\dagger$ and
$P_s=i(b_s^\dagger-b_s)$. Assume centered independent Gaussian modes,
initially independent of the body and its retained reference,
a known common $k$, and $\operatorname{Var}(P_s)/(4k)\le\sigma_{\max}^2$.
Controlled total parity conjugates
$Z_L\sqrt{k}\sum_s v_s(a)Q_s$ to
$I_L\sqrt{k}\sum_s v_s(a)Q_s$. After this coherent echo the signed
readings $Y_s=-P_s/(2\sqrt{k})$ have mean $v_s(a)$ while the logical
state remains unchanged. Two commands, $a=0,h$ with $h\ne0$, identify
\[
 \mathbf g_1=\mathbb E\mathbf Y(0),\qquad
 \mathbf g_2=\frac{\mathbb E\mathbf Y(0)-\mathbb E\mathbf Y(h)}h.
\]
This does not require the Gaussian variances as labels.

Suppose $|v_s(0)|,|v_s(h)|\le B$ and use a finite readout that clips $Y_s$
to $[-Y,Y]$, $Y>B$, then rounds to the nearest point of a prescribed
uniform grid of spacing $q>0$ in that interval, adding the two endpoints
if needed. This alphabet has at most $\lceil2Y/q\rceil+2$ symbols and
absolute rounding error at most $q/2$.
For $N$ fresh queries at each command, define
\[
 \eta=Y\sqrt{\frac{2\log(4d/\zeta)}N}+\frac q2+
       2\sigma_{\max}\varphi\!\left(\frac{Y-B}{\sigma_{\max}}\right),
 \qquad \varphi(z)=\frac{e^{-z^2/2}}{\sqrt{2\pi}}.
\]
With probability at least $1-\zeta$, the estimated residual on $|a|\le A$
has vector error at most $e_v=\sqrt d\,\eta(1+2A/|h|)$.
If its true norm is bounded by $B_v$, then
\begin{equation}
 |\widehat V(a)-\|\mathbf v(a)\|^2|\le2B_ve_v+e_v^2.
 \label{eq:signed-source-error}
\end{equation}
The readout uses a finite alphabet, with at most
$\lceil\log_2(\lceil2Y/q\rceil+2)\rceil$ bits per species per query.
\end{proposition}
\begin{proof}
Total parity reverses every $Q_s$, so the coherent parity calculation of
Q-RAIF~B~\cite{Liu2026QRAIFB} gives a field-only displacement.
The canonical relation $[Q_s,P_s]=2i$ shifts $P_s$ by
$-2\sqrt{k}v_s(a)$ and leaves its variance unchanged.
For a Gaussian with $|\mu|\le B$ and standard deviation at most
$\sigma_{\max}$, the clipping bias is at most
$2\sigma_{\max}\varphi((Y-B)/\sigma_{\max})$, by integrating the two
tails. Hoeffding for the bounded clipped readings and a union over the
$2d$ means give the first term of $\eta$; rounding adds at most $q/2$.
Thus each coordinate error of $\mathbf g_1$ is at most $\eta$ and of
$\mathbf g_2$ at most $2\eta/|h|$. The Euclidean bound follows.
Expanding $\|\mathbf v+\delta\mathbf v\|^2$ proves
Eq.~\eqref{eq:signed-source-error}.
\end{proof}

The sign reference, species-resolved quadrature detector and complete mode
control are its additional permissions. Two commands here mean $2N$ new field
queries, each with $d$ digitized readings; mode preparation, coupling, capture
and memory remain charged. Unknown incoming means would confound the shift.
Together with the matched JT equations either identified source predicts a
response from actual records rather than from a supplied coupling label.
Boundary sensing still checks the physical response itself. Neither route
replaces the wall, clock and thermal-export construction in a closed loop.


\subsection{The command carrier matters}

How strongly the gravitational field records the command statistics depends on how the command is physically carried.

\begin{proposition}[Phase-carried commands]
\label{prop:carrier}
Keep the coupling switched on at every opportunity and carry the command in the drive phase: a correction uses the input $|\sqrt N\rangle$ and no correction uses $|i\sqrt N\rangle$. Then (i) the no-correction branch flips the bit with probability $1-s_N$ and the correction branch with probability $s_N$; (ii) both inputs carry the same field energy $N\bar\omega$, and the branch sources are $4N[\operatorname{Re}\psi]^2+F_g$ and $4N[\operatorname{Im}\psi]^2+F_g$; (iii) for the worked packet both branch sources have the same $E_0$, $E_1$, $E_2$, so $\mathcal C_\infty$ does not depend on the command rate, while the dilaton coefficients $A(u)$, $B(u)$ at finite $u$ do.
\end{proposition}

\begin{proof}
With $Q=b+b^\dagger$, the bit channel multiplies coherences in the $X$ basis by $\langle\beta|e^{-2i\lambda Q}|\beta\rangle=e^{-2\lambda^2}e^{-4i\lambda\operatorname{Re}\beta}$. For $\beta=i\sqrt N$ this factor is $e^{-2\lambda^2}$, real and positive, which gives (i). The stress of a coherent input $\beta$ is $(\psi\beta+\psi^*\beta^*)^2$, which is $4N[\operatorname{Re}\psi]^2$ for $\beta=\sqrt N$ and $4N[\operatorname{Im}\psi]^2$ for $\beta=i\sqrt N$, and the coupling adds $F_g$ in both branches. The equal moments of the two quadrature profiles give (iii) through Eq.~\eqref{eq:invariant}; the difference of $A$ between branches is proportional to $\operatorname{Re}[a+i(u-u_0)]^{-2}$, which is nonzero for generic $u$.
\end{proof}

The hiding has a price. With phase-carried commands the uncorrected steps also err, so the world mean obeys $m\mapsto c(2s_N-1)[(1-h)m+h\nu]$ on the calibrated branch. At $N=16$ and a known interface the stationary loss rises from $0.1896$ with gated couplings to $0.2334$, and the switching work per opportunity rises from $0.0039\bar\omega$ to $0.0386\bar\omega$. The phase must also be set on the incoming packet before the encounter, which for full-line packets has the same causal limitation as on-demand emission. At $N=16$ and $\kappa=10^{-4}$ the coefficient $A$ at the packet centre differs between the two branch sources by $3.4\times10^{-4}$, and the difference decays as $u^{-2}$.

So the gravitational trace of a selection is not universal. With gated couplings it sits in the rate of an energy-carrying event; with phase-carried commands it sits in the shape of the source and disappears from the asymptotic invariant of this example. A general law connecting selection to gravitational response must therefore name the carrier of the selection.


\subsection{A worked packet}
\label{sec:packet}

Take the normalized mode $f(\omega)=2a\sqrt\omega\,e^{-a\omega}e^{i\omega u_0}$, for which $\bar\omega=1/a$ and
\begin{equation}
\psi(u)=-\frac{ia}{\sqrt\pi\,[a+i(u-u_0)]^2}.
\end{equation}
The two quadrature profiles $4[\operatorname{Re}\psi]^2$ and $4[\operatorname{Im}\psi]^2$ have different shapes but the same normalized moments $(E_0,E_1,E_2)=(1,u_0,u_0^2+a^2)/a$. They have power-law tails: $[\operatorname{Im}\psi]^2\sim a^2/(\pi u^4)$ and $[\operatorname{Re}\psi]^2=O(u^{-6})$. The absolute third moment of $F_g$ diverges logarithmically; the first three moments are finite and are exactly those used in Eq.~\eqref{eq:invariant}. The switch determined by $g'=4\lambda\operatorname{Im}\psi$ and $g(\pm\infty)=0$ is
\begin{equation}
 g(u)=-\frac{4\lambda a(u-u_0)}{\sqrt\pi\,[a^2+(u-u_0)^2]}.
 \label{eq:exact-switch}
\end{equation}
Both $g$ and $F_g$ have full-line support. The formulas and table below are exact
for this scattering idealization with asymptotic coefficient data. A command decided
at a finite time cannot gate its entire switch; finite-time feedback uses the
following compact approximation.

\begin{proposition}[Causal compact approximation]\label{prop:compact}
Choose a fixed even $\eta\in C_c^\infty(\mathbb R)$ with $0\le\eta\le1$,
$\eta=1$ on $[-1,1]$ and $\eta=0$ outside $[-2,2]$, and set
$g_T(u)=\eta((u-u_0)/T)g(u)$. For $T\ge a$, write $F_T=(g_T')^2/4$.
The unrescaled switching stresses obey
\begin{align}
 \left|\int(F_T-F_g)du\right|&=O(\lambda^2a^2T^{-3}),\nonumber\\
 \int(u-u_0)(F_T-F_g)du&=0,\nonumber\\
 \left|\int(u-u_0)^2(F_T-F_g)du\right|&=O(\lambda^2a^2T^{-1}).
 \label{eq:compact-moments}
\end{align}
The one-particle smearing vectors converge in norm, with difference
$O(|\lambda|a/T)$. Hence the corresponding coherent-drive gate channels converge.
A gate chosen before $u_0-2T$ can be applied causally. Its duration and any delay
between the read and the gate are part of its task schedule.
\end{proposition}
\begin{proof}
With $w=u-u_0$ and $C=4|\lambda|a/\sqrt\pi$, the tails satisfy
$|g|\le C/|w|$ and $|g'|\le C/w^2$ for $|w|\ge T\ge a$.
The derivative of the cutoff obeys the same $w^{-2}$ bound on its annulus,
up to a constant depending only on $\eta$. Integrating $w^{-4}$ and $w^{-2}$
gives the first and third bounds; parity gives the second.
For $h=g_T-g$, $\|h\|_2=O(|\lambda|aT^{-1/2})$ and
$\|h'\|_2=O(|\lambda|aT^{-3/2})$. Plancherel and Cauchy--Schwarz bound
$\int_0^\infty\omega|\hat h|^2d\omega/(4\pi)$ by
$\|h\|_2\|h'\|_2/2$, giving the one-particle estimate.
The coherent-state characteristic function is continuous in this norm.
\end{proof}

At finite $T$, the drive mode, its overlap with the incoming beam and the gate error
must be recomputed from $g_T$ before assigning values to a finite-cutoff table.
If the drive is matched to the new mode, its
quadrature profiles and first three stress moments converge as well. The real and
imaginary profiles are Hilbert-transform pairs. For $r=(g_T-g)'$, the identity
$u\mathcal Hr=\mathcal H(ur)+\pi^{-1}\int r$ in the convention
$\mathcal Hr(u)=\pi^{-1}\operatorname{pv}\int r(v)/(u-v)dv$ has zero last term.
The $L^2$ isometry of $\mathcal H$ therefore gives
$\|u\mathcal Hr\|_2=\|ur\|_2\to0$. Together with convergence of the mode
normalization, this proves the weighted profile convergence.
Equation~\eqref{eq:source-loss} holds exactly for each compact
implementation with its actual gate parameters. The same discrete feedback law
applies when the target is held between its read and gate; noise during a new delay
must instead be included in the transition law.

With $A_{\rm tot}=N+p\lambda^2$ the asymptotic invariant is $\mathcal C_\infty=H_0^2+2\kappa H_0u_0A_{\rm tot}/a-\kappa^2A_{\rm tot}^2$. Relative to the same beam without couplings,
\begin{equation}
\delta\mathcal C(p)=p\lambda^2\Bigl[\frac{2\kappa H_0u_0}a-\kappa^2\bigl(2N+p\lambda^2\bigr)\Bigr],
\label{eq:dC}
\end{equation}
which has an explicit quadratic term in $p$. Table~\ref{tab:response} uses $a=1$, $u_0=3$, $H_0=1$, $\kappa=10^{-4}$, the task parameters $\nu=0.8$, $c=0.9$, $h=0.4$ of Ref.~\cite{Liu2026QRAIFB}, and the uncalibrated ($q=\tfrac12$) and calibrated ($32$ passive blocks, $q=0.0404888$) coupling rates.

\begin{table}[ht]
\centering
\footnotesize
\setlength{\tabcolsep}{4pt}
\begin{tabular}{@{}rrrrrrrr@{}}
\toprule
$N$ & gate error & $p$, $q=\tfrac12$ & $p$, $q_{32}$ & $\delta\mathcal C$, $q=\tfrac12$ & $\delta\mathcal C$, $q_{32}$ & beam share & quad./lin.\\
\midrule
1 & 0.3544 & 0.15525 & 0.11413 & $5.746\times10^{-5}$ & $4.224\times10^{-5}$ & $7.0\times10^{-2}$ & $-1.2\times10^{-6}$\\
4 & 0.1327 & 0.15153 & 0.10698 & $1.402\times10^{-5}$ & $9.897\times10^{-6}$ & $4.1\times10^{-3}$ & $-2.7\times10^{-7}$\\
16 & 0.0371 & 0.15033 & 0.10468 & $3.476\times10^{-6}$ & $2.420\times10^{-6}$ & $2.5\times10^{-4}$ & $-6.7\times10^{-8}$\\
64 & 0.0095 & 0.15002 & 0.10408 & $8.657\times10^{-7}$ & $6.006\times10^{-7}$ & $1.6\times10^{-5}$ & $-1.7\times10^{-8}$\\
\bottomrule
\end{tabular}
\caption{JT invariant response to the coupling-added source of the worked packet. ``Beam share'' is $\delta\mathcal C(p)$ at $q_{32}$ divided by the beam's own change $\mathcal C_\infty(N)-\mathcal C_\infty(0)$; ``quad./lin.'' is the ratio of the $p^2$ term of Eq.~\eqref{eq:dC} to the term linear in $p$.}
\label{tab:response}
\end{table}

At $N=16$ the ratio of the two responses is $0.696342074$, against the coupling-rate ratio $0.696342054$. For $\kappa=0.01$ and $0.05$ ($\kappa N=0.16$ and $0.8$) the response ratios are $0.6963442$ and $0.6963560$; the quadratic term stays small because $p\lambda^2$ is small, not because $\kappa N$ is. The calibration-dependent part is a small fraction of the beam's own response, and that fraction varies roughly as $N^{-2}$: it is $7\%$ at $N=1$, where the gate fails $35\%$ of the time, and $1.6\times10^{-5}$ at $N=64$. Its size is therefore a property of the chosen drive and background, not a general feature of selection, and a field-energy ratio such as $p\lambda^2/N$ is not a general invariant response ratio. In this example the bracket of Eq.~\eqref{eq:dC} is positive, so better calibration lowers $\delta\mathcal C$; for other backgrounds the bracket can vanish or change sign.

The mean-source solution is not the mean of branch solutions. For a single command $I\sim\mathrm{Bernoulli}(p)$, $\mathbb E\,\mathcal C_\infty(N+\lambda^2I)-\mathcal C_\infty(N+p\lambda^2)=-\kappa^2\lambda^4p(1-p)$, which counts only the classical command randomness and not field fluctuations within a branch. With $M$ independent copies of the whole experiment, each with its own hidden instrument, calibration records and drives, and with $G=G_1/M$, the classical source-fraction variance of smooth stress readings falls as $1/M$. A common hidden label shared by all copies does not give this independence, and a coherent drive still has nonzero energy variance. The exact mean-source solution does not show that a single quantum world has a definite classical geometry.

\subsection{What remains input}
\label{sec:inputs}

The construction takes as given the JT action and its anti-de Sitter background, the
matter field and its central-charge matching, the boundary data and the identification
of the downstream null coordinate with the preparation coordinate, the semiclassical
use of mean stress, and the actuator's field, switch, clock and fresh ports. The
two-dimensional sector theorem permits either one or several cones; it neither
derives this matter--gravity identification nor forbids a dynamical route to it~\cite{Liu2026QRAIF_A}.
The source is computed outside the apparatus, for the free outgoing field; the stress
of the switch and its support is not included, and a covariant detector stress tensor
shows that it cannot simply be dropped in the coupling region~\cite{Perche2025}.
Within these inputs the response is exact.


\section{Finite boundary sensing and returned records}
\label{app:boundary-return}

\subsection{A conditional readback route}
\label{sec:echo}

The matter action in Section~\ref{sec:inputs} does not directly couple to
$\Phi$. A bulk detector of the matter metric alone therefore does not acquire a
local dilaton readout. Boundary motion offers a different route if a physical
reflecting boundary with a known dilaton condition is available.

\begin{proposition}[A dilaton-blind record law]\label{prop:blind-readback}
Compare two dilaton solutions with the same matter metric, the same accessible
matter input states and the same initial controller state. Suppose every
instrument and conditional evolution used by the controller depends on this
metric and its earlier records, but not on the dilaton or a boundary that senses
it. Then every finite adaptive record sequence has the same distribution in the
two solutions.
\end{proposition}
\begin{proof}
The first instrument has the same input and conditional outcome law. Given
any earlier record sequence, the next instrument and input evolution are again
the same. Induction gives equality of the joint record laws. Private random bits
can be included in the controller state and the induction.
\end{proof}
This statement does not exclude inference from previously recorded sources.
It excludes a new dilaton signal from an otherwise identical metric-only
experiment. A reflecting boundary evades its hypothesis because its position
depends on $\Phi$.

\begin{proposition}[Reconstruction from three boundary echoes]\label{prop:echo}
In a source-free interval let $a,b,\gamma$ be constant, $\mathcal C=a^2+b\gamma>0$,
and suppose the apparatus can time reflections at the supplied boundary
$\Phi=\Phi_b>\sqrt{\mathcal C}$. Use one calibrated null coordinate frame.
Three reflection pairs $(u_i,v_i)$, with distinct $u_i$ and
$d_i=(a+\Phi_b)u_i-\gamma\ne0$, determine $a,b,\gamma$ uniquely:
\begin{equation}\label{eq:echo-inverse}
 \begin{pmatrix}
 1-u_1v_1&u_1&v_1\\
 1-u_2v_2&u_2&v_2\\
 1-u_3v_3&u_3&v_3
 \end{pmatrix}
 \begin{pmatrix}a\\b\\\gamma\end{pmatrix}
 =\Phi_b
 \begin{pmatrix}1+u_1v_1\\1+u_2v_2\\1+u_3v_3\end{pmatrix}.
\end{equation}
Writing this matrix as $M$, its determinant is
\[
 \det M=
 \frac{2\Phi_b(\mathcal C-\Phi_b^2)
       \prod_{i<j}(u_j-u_i)}{\prod_i d_i}\ne0.
\]
If the $u_i$ are exact and the measured returns are $v_i+\delta v_i$, an
invertible measured matrix $\widetilde M$ yields the exact coefficient error
\begin{equation}\label{eq:echo-error}
 \delta(a,b,\gamma)^T
 =\widetilde M^{-1}(\delta v_i d_i)_{i=1}^3,\qquad
 \|\delta(a,b,\gamma)\|_2
 \le\|\widetilde M^{-1}\|_2\|(\delta v_i d_i)_i\|_2.
\end{equation}
\end{proposition}
\begin{proof}
The boundary is
$v=V(u)=(a+bu-\Phi_b)/[(a+\Phi_b)u-\gamma]$,
with $V'=(\Phi_b^2-\mathcal C)/[(a+\Phi_b)u-\gamma]^2>0$, so it is timelike.
Its equation gives Eq.~\eqref{eq:echo-inverse}. Multiplying row $i$ of $M$ by
$d_i$ leaves evaluations of three quadratic polynomials; their coefficient
determinant is $2\Phi_b(\mathcal C-\Phi_b^2)$, giving the displayed Vandermonde
factor. In the noisy equation the difference between the measured right side
and $\widetilde M(a,b,\gamma)^T$ is exactly $(\delta v_i d_i)_i$.
\end{proof}

Null signals can carry these timestamps back to a receiver before the next
command. They reconstruct the coefficients, hence $\mathcal C$, in a declared
source-free interval. The inverse becomes ill-conditioned near
$\Phi_b=\sqrt{\mathcal C}$ or coincident sample times. Exact reconstruction is
a test-probe limit, not three cost-free finite-energy measurements. A usable
protocol must supply the reflecting boundary, its support, calibrated clocks,
timing resolution, causal waiting time and probe stress; the last changes the
background beyond that limit. Six timestamps stored with $b$ bits each occupy
$6b$ logical bits, not a determined physical energy or complete memory cost.
This gives an explicit sensing interface to implement, rather than granting the
controller direct access to $\Phi$.

\subsection{Probe backreaction and a returned program}
\label{sec:returned-program}

A finite probe contributes stress, including the reflected component. The inverse
can retain both chiral components instead of treating the probe as free. Work in
the mean-source JT equations of Section~\ref{sec:inputs}, with the same trace-anomaly
subtraction. Suppose the complete additional bulk source in the sensing interval
is calibrated as $T_{uu}=f(u)$, $T_{vv}=h(v)$, $T_{uv}=0$. Define
\begin{align}
 A_f(u)&=-\kappa\int_{u_*}^{u}(u-z)f(z)\,dz,&
 B_f(u)&=-\kappa\int_{u_*}^{u}z(u-z)f(z)\,dz,\nonumber\\
 A_h(v)&=-\kappa\int_{v_*}^{v}(v-z)h(z)\,dz,&
 B_h(v)&=-\kappa\int_{v_*}^{v}z(v-z)h(z)\,dz,\nonumber\\
 K(u,v)&=A_f(u)+vB_f(u)+A_h(v)+uB_h(v).
 \label{eq:two-chiral-K}
\end{align}
The matching convention fixes the lower limits and the homogeneous coefficient
data. These are calibrated sources, not a derived reflection law for the wall.
If reflection or support stresses introduce other components or unknown source
terms, the following inverse must be changed accordingly.

\begin{proposition}[Source-compensated echo inverse]\label{prop:compensated-echo}
For a known $K$ as in Eq.~\eqref{eq:two-chiral-K}, the solution is
\begin{equation}\label{eq:two-chiral-phi}
 \Phi(u,v)=\frac{a+bu+\gamma v-auv+K(u,v)}{1+uv}.
\end{equation}
Three points on the supplied boundary $\Phi=\Phi_b$ determine the coefficients
$\boldsymbol\theta=(a,b,\gamma)^T$ whenever their matrix $M$ in
Eq.~\eqref{eq:echo-inverse} is invertible:
\begin{equation}\label{eq:compensated-inverse}
 M\boldsymbol\theta=
 \bigl(\Phi_b(1+u_iv_i)-K(u_i,v_i)\bigr)_{i=1}^3.
\end{equation}
The source-free determinant formula need not hold for these altered echoes;
invertibility persists under sufficiently small perturbations of an invertible
source-free matrix.

Suppose the measured $u_i,v_i$ have errors at most $\delta_u,\delta_v$,
the calibrated $K$ has pointwise error at most $\eta_K$, and the measured
matrix $\widetilde M$ is invertible. On a domain joining true and measured points,
bound the derivatives of
\[
 R(u,v)=(\Phi_b-a)+(\Phi_b+a)uv-bu-\gamma v-K(u,v)
\]
by $|R_u|\le L_u$ and $|R_v|\le L_v$. Then
\begin{align}
 e:=\|\widehat{\boldsymbol\theta}-\boldsymbol\theta\|_2
 &\le\sqrt3\|\widetilde M^{-1}\|_2
       (L_u\delta_u+L_v\delta_v+\eta_K),\label{eq:compensated-error}\\
 |\widehat{\mathcal C}-\mathcal C|
 &\le\sqrt{4a^2+b^2+\gamma^2}\,e+e^2,
 \qquad \widehat{\mathcal C}=\widehat a^2+\widehat b\widehat\gamma.
 \label{eq:invariant-read-error}
\end{align}
If the next projection angle is $\vartheta_{\rm next}=F(\widehat{\mathcal C})$
for a supplied Lipschitz program $F$, its error relative to $F(\mathcal C)$ is
at most $\operatorname{Lip}(F)$ times Eq.~\eqref{eq:invariant-read-error}.
\end{proposition}
\begin{proof}
The $f$ contribution is the retarded solution already obtained in
Section~\ref{sec:source}; interchange $u,v$ for the $h$ contribution. Their
sum satisfies all three JT equations and gives Eq.~\eqref{eq:two-chiral-phi}.
Evaluation on the boundary yields Eq.~\eqref{eq:compensated-inverse}.
The true residual $R(u_i,v_i)$ is zero. The residual at each measured point,
including the calibration error, is bounded by
$L_u\delta_u+L_v\delta_v+\eta_K$. Multiplication by
$\widetilde M^{-1}$ gives Eq.~\eqref{eq:compensated-error}. Expanding
$\widehat a^2+\widehat b\widehat\gamma$ leaves a linear term with gradient
$(2a,\gamma,b)$ and the quadratic term
$(\delta a)^2+\delta b\delta\gamma$, whose magnitude is at most $e^2$.
The program bound follows from its Lipschitz property.
\end{proof}
An apparatus needs a certified parameter range to bound $L_u,L_v$ and the
conditioning of the inverse before use. Neither the assumed calibration of $K$
nor this error bound supplies a timing POVM or the wall's support energy.

Here is a causal one-step return in that declared geometry. Take the selected
projection of Section~\ref{sec:projection-source}, a fresh vacuum field and
the independent bit state $(I+0.8Z)/2$. Use a smooth compact switch
\begin{equation}\label{eq:compact-return-switch}
 g(u)=A\exp\left[-\frac{1}{1-((u-3)/0.25)^2}\right]
 \quad (|u-3|<0.25),\qquad g(u)=0\quad\hbox{otherwise},
\end{equation}
with $A$ fixed by $\int(g')^2/4\,du=0.25$. Unlike the full-line packet of
Section~\ref{sec:source}, this source finishes before the readout begins.
Set $H_0=L=1$, $\kappa=0.05$ and $\Phi_b=3$. The source moments are
$E_0=0.25$, $E_1=0.75$, $E_2=2.257177876$; a coupled population $r$ gives
\[
 (a,b,\gamma)=(1+\kappa rE_1,-\kappa rE_0,\kappa rE_2).
\]
The projections $P_0$ and $P_{\pi/2}$ give $r=0.9$ and $0.5$ respectively,
and hence $\mathcal C=1.067496366$ and $1.037498878$.

Supply a timelike receiver at $v=-\alpha/u$, $\alpha=1/4$. Its proper clock is
$\tau=2L\sqrt\alpha\ln u/(1-\alpha)$, up to an additive constant.
Emit at receiver coordinates $t_i=4,6,9$, with $v_i=-\alpha/t_i$. The outward
ray has fixed $v_i$ and meets the boundary at
\[
 u_i=\frac{a-\Phi_b+\gamma v_i}{(a+\Phi_b)v_i-b}>t_i.
\]
The reflected ray has fixed $u_i$; reception at the receiver therefore occurs
at the same coordinate $u_i$, with $v=-\alpha/u_i>v_i$.
For $r=0.9$ the receive coordinates are $8.18983,12.56501,19.53471$;
for $r=0.5$ they are $8.10371,12.30538,18.81553$. Throughout these intervals
the receiver obeys $\sqrt{\mathcal C}<\Phi<\Phi_b$ and $1+uv>0$, inside the
usual cutoff and outside the stationary horizon. A new compact coupling can
start after $u=20$.
Writing $\mathcal C_{\rm mid}$ for the midpoint of the two invariants, choose
\begin{equation}\label{eq:return-program}
 F(\mathcal C)=\frac\pi4+10(\mathcal C-\mathcal C_{\rm mid}).
\end{equation}
The two returned angles are $0.935385602$ and $0.635410725$. If the inverse
error stays below the corresponding separation margin, the measurement changes
the next coupled projection and its represented algebra $\{P_{F(\mathcal C)}\}'$.
This is a response-dependent return map, not a proof that this program is optimal
or that its interface formed through learning. The controller uses the measured
invariant; computing one from assumed initial data is a different protocol.

The listed times use the source-free geometry. With finite probes, the return
must instead use Eq.~\eqref{eq:compensated-inverse} and the altered boundary.
As a bulk check, two calibrated positive impulsive fluxes of integrals
$10^{-5}$ in $u$ and $2\times10^{-5}$ in $v$ give a coefficient bias of about
$5.7\times10^{-6}$ if ignored; source compensation removes it to numerical
precision. This is not a physical reflection or timing experiment. For a
translation-only readout with the listed proper-delay difference
$\Delta\tau=0.05001338$, the sharp signed-delay bound in Q-RAIF~B requires
at least $11.04\hbar$ in this clock's energy units for binary error $0.1$
\cite{Liu2026QRAIFB}. A real echo can also carry redshift and amplitude information,
so that restriction is not a universal bound on geometric sensing. It shows why
one cannot assign both negligible probe stress and a finite resolution without
constructing the readout. The boundary, receiver, timing instrument, full
probe-induced stress and support costs remain explicit apparatus inputs to this
conditional return; mean-source evolution does not bound stress fluctuations.

\subsection{A quantum boundary record without echo timestamps}
\label{sec:boundary-current-sensor}

A current detector gives a second readback route. It measures fluctuations over
a prescribed interval of proper time, rather than the arrival time of a reflected
signal. The boundary condition still does essential work: the detector is carried
by the supplied timelike curve $\Phi=\Phi_b$. A detector on the same fixed bulk
worldline in both solutions would remain subject to
Proposition~\ref{prop:blind-readback}.

Use a transparent chiral current interface, with the vacuum defined in $u$ and
the current normalization of Section~\ref{sec:source}. Its input on the interval
to be sensed has zero mean and covariance
$\langle j(u)j(u')\rangle=-[4\pi(u-u'-i0)^2]^{-1}$. The apparatus supplies a
proper-time clock, a degenerate detector bit initially in $|0\rangle$, and a real
compact switch $\zeta(\tau)$. The finite interaction is
\begin{equation}\label{eq:boundary-sensor-unitary}
 U_s=e^{-iX_sQ_g},\qquad
 Q_g=\int g(u)j(u)\,du,\qquad
 g(u(\tau))=\zeta(\tau).
\end{equation}
Equivalently its integrated interaction is
$\int\zeta(\tau)X_s\dot u(\tau)j(u(\tau))\,d\tau$.
The fixed detector operator and c-number current commutator make the
time-ordering correction a scalar phase. This is the same gapless-detector
method used for exact field channels~\cite{Landulfo2016}; the finite switch
defines a smeared current, not an unsmeared instantaneous Hamiltonian.
We read $Z_s$ after contact and retain the bit $R=0,1$, with $R=1$ the flipped
outcome. Clock calibration, preparation and recording permissions are supplied
resources.

\begin{proposition}[Invariant-dependent finite detector law]
\label{prop:boundary-sensor-law}
On a source-free background with $0<\mathcal C<\Phi_b^2$, a future boundary
segment has
\begin{equation}\label{eq:boundary-clock-speed}
 \dot u=\frac{bu^2+2au-\gamma}{2L\sqrt{\Phi_b^2-\mathcal C}}>0,
 \qquad
 \lambda_b=\frac{\sqrt{\mathcal C}}{L\sqrt{\Phi_b^2-\mathcal C}}.
\end{equation}
In the test-background calculation, Eq.~\eqref{eq:boundary-sensor-unitary}
has the record law
\begin{align}
 k_b(\mathcal C)&=\frac{1}{4\pi}\int_0^\infty
  \omega\coth\frac{\pi\omega}{\lambda_b}
  |\widehat\zeta(\omega)|^2\,d\omega,
 \qquad \widehat\zeta(\omega)=\int e^{i\omega\tau}\zeta(\tau)\,d\tau,
 \label{eq:boundary-sensor-variance}\\
 \Pr(R=1\mid\mathcal C)&=\frac{1-e^{-2k_b(\mathcal C)}}2.
 \label{eq:boundary-sensor-record}
\end{align}
For every nonzero switch with finite variance, the probability is strictly
increasing in $\mathcal C$. It depends on the invariant, not on the three
coefficients separately. For two equally likely backgrounds the optimal
decision from this one bit has error
$[1-|p_1-p_0|]/2$, where $p_i=\Pr(R=1\mid\mathcal C_i)$.
\end{proposition}
\begin{proof}
Pulling back the metric to the level curve gives
$d\tau/du=2L\sqrt{\Phi_b^2-\mathcal C}/(bu^2+2au-\gamma)$.
The Riccati equation in Eq.~\eqref{eq:boundary-clock-speed} makes $u(\tau)$
a fractional-linear transform of $e^{\lambda_b\tau}$ on each regular segment.
M\"obius invariance of the current two-point function therefore gives
\[
 \langle\dot u(\tau)j(u(\tau))\dot u(\tau')j(u(\tau'))\rangle
 =-\frac{\lambda_b^2}{16\pi
       \sinh^2[\lambda_b(\tau-\tau'-i0)/2]}.
\]
Its symmetrized spectrum yields Eq.~\eqref{eq:boundary-sensor-variance}; the
pullback vacuum is KMS at temperature $\lambda_b/(2\pi)$ in this proper clock.
For a centred Gaussian field,
$\langle e^{2iQ_g}\rangle=e^{-2\langle Q_g^2\rangle}$, hence the detector's
$Z_s$ expectation is $e^{-2k_b}$. The factors $\lambda_b$ and
$\coth(\pi\omega/\lambda_b)$ strictly increase with $\mathcal C$ for
$\omega>0$. A nonzero compact switch has Fourier weight on a set of positive
measure, so the integral strictly increases. The binary decision formula is
the total-variation formula for two Bernoulli laws.
\end{proof}
This response uses the familiar Hawking--Unruh covariance. It does not add a
direct matter--dilaton coupling, and the KMS temperature here is the boundary
proper-time temperature, not an automatically thermal exhaust in the running
rapidity clock.
Without calibration of $L$, $\Phi_b$ and the proper-time scale, these records
identify $\lambda_b$ rather than $\mathcal C$. A clock-scale error can mimic
an invariant change; this instrument alone cannot distinguish them.

\subsection{The sensing interaction changes the background}
\label{sec:driven-boundary-sensor}

The same detector interaction adds the chiral mean stress
$f(u)=[g'(u)]^2/4$. To include it, choose a boundary-source completion:
this current is transmitted into the JT region, the coupling-added bulk source
is precisely $T_{uu}=f(u)$, and the other state-dependent chiral source is zero.
Keep the trace-anomaly subtraction of Section~\ref{sec:inputs}.
This is a specified boundary injection model, not the reflected-probe model of
Section~\ref{sec:echo}. No unknown reflected component is fitted as a known
function. Forces supporting the boundary, an ultraviolet realization of the
contact and the clock's stress require a further apparatus model.

\begin{proposition}[Self-consistent sensing source in mean JT geometry]
\label{prop:driven-boundary-sensor}
In this completion, put $R_b=\sqrt{\Phi_b^2-\mathcal C}$ and let dots denote
proper-time derivatives. Starting with the matched coefficient data, solve
\begin{align}
 \dot u&=\frac{bu^2+2au-\gamma}{2LR_b},&
 f(u(\tau))&=\frac{\dot\zeta^2}{4\dot u^2},\nonumber\\
 (\dot a,\dot b,\dot\gamma)&=
       \frac{\kappa\dot\zeta^2}{4\dot u}(u,-1,u^2).
 \label{eq:boundary-sensor-ode}
\end{align}
On a regular finite segment this gives the boundary trajectory and the complete
coupling-added bulk source together. Its exact balances are
\begin{align}
 \dot{\mathcal C}&=\frac{\kappa L}{2}R_b\dot\zeta^2,&
 R_b(\tau)&=R_b(0)-\frac{\kappa L}{4}
               \int_0^\tau\dot\zeta(s)^2\,ds,\label{eq:sensor-invariant-balance}\\
 E_s^{(\tau)}&:=\int f\dot u^2\,d\tau
       =\frac14\int\dot\zeta^2\,d\tau
       =\frac{R_b(0)-R_b(\tau_f)}{\kappa L},&
 E_s^{(u)}&:=\int f\,du
       =\frac14\int\frac{\dot\zeta^2}{\dot u}\,d\tau.
 \label{eq:sensor-energy-balance}
\end{align}
For the actual, generally nonstationary map $u(\tau)$, the detector variance is
\begin{equation}\label{eq:driven-sensor-variance}
 k[g]=\frac{1}{4\pi}\int_0^\infty\omega|\widehat g(\omega)|^2d\omega
     =\frac{1}{8\pi}\int_{\mathbb R^2}
       \frac{[g(u)-g(v)]^2}{(u-v)^2}\,du\,dv.
\end{equation}
Its record probability remains $(1-e^{-2k[g]})/2$; replacing this variance by
the stationary KMS integral during contact would neglect backreaction.
\end{proposition}
\begin{proof}
The retarded JT solution has running coefficients
$a'=\kappa u f$, $b'=-\kappa f$ and $\gamma'=\kappa u^2 f$, with primes in
$u$. In the derivative of the boundary level equation the source contributions
cancel: its numerator derivative is $b$ and its denominator derivative is
$a+\Phi_b$. Thus its slope is still
$(\Phi_b^2-\mathcal C)/[(a+\Phi_b)u-\gamma]^2$ with running coefficients.
The induced metric gives the first equation in
Eq.~\eqref{eq:boundary-sensor-ode}, while $g'=\dot\zeta/\dot u$ gives the
second. Differentiating $a^2+b\gamma$ and substituting $\dot u$ gives
Eq.~\eqref{eq:sensor-invariant-balance}. The energy balances follow by change
of variable. The source $f$ follows from $U_s^\dagger jU_s=j+X_sg'/2$ and
the initially independent, centred detector with $\langle X_s\rangle=0$.
The incoming current on the deterministic mean trajectory is Gaussian; its
vacuum variance has the Fourier expression in
Eq.~\eqref{eq:driven-sensor-variance}. Parseval's
identity and $\int(1-\cos\omega x)/x^2\,dx=\pi|\omega|$ give the second.
\end{proof}
In this model the mean source fixes a classical trajectory before the detector
outcome is conditioned on. It is not a theory of outcome-conditioned quantum
geometry. The inequality
$\kappa L\int\dot\zeta^2/4<R_b(0)$ prevents closure of the boundary gap;
one must also check $\dot u>0$ and $1+uv>0$ over the chosen window.
The two energies in Eq.~\eqref{eq:sensor-energy-balance} refer to different
clocks. Neither is the total preparation, support, readout or erasure cost.
For any sequence obeying this positive-source completion at unchanged
$\Phi_b,L,\kappa$, the cumulative proper-clock source charge satisfies
\begin{equation}\label{eq:sensor-cumulative-budget}
 \sum_n E_{s,n}^{(\tau)}
 <\frac{\sqrt{\Phi_b^2-\mathcal C_{\rm initial}}}{\kappa L}.
\end{equation}
This follows directly from $R_b>0$ and
Eq.~\eqref{eq:sensor-invariant-balance}. A nonzero charge per contact therefore
rules out infinitely many such contacts at a fixed cutoff. It is an upper
budget for this positive-source boundary family, not a universal lower bound
on sensing energy or a prohibition of protocols with a different net source.

\begin{proposition}[Finite sensing survives its own added source]
\label{prop:weak-boundary-sensor}
Take two source-free initial backgrounds with different invariants, with the
same $\Phi_b,L$ and a common proper-time contact interval on which both
unperturbed curves are regular. Fix a nonzero $C^2$ compact switch $\zeta_0$.
At fixed $\kappa$, the completion in
Proposition~\ref{prop:driven-boundary-sensor}, with
$\zeta=\epsilon\zeta_0$, has distinct record laws for all sufficiently small
$\epsilon>0$. The added energies and invariant changes are finite and of
order $\epsilon^2$.
\end{proposition}
\begin{proof}
On the compact regular segment the trajectory equations are locally Lipschitz,
and their source-dependent terms are $O(\kappa\epsilon^2)$. Continuous
dependence preserves the domain and gives
$u_\epsilon-u_0=O(\kappa\epsilon^2)$ in $C^1$. The corresponding compact
smears obey
$\|g_\epsilon-\epsilon g_0\|_0=O(\kappa\epsilon^3)$:
the inverse clocks and their first derivatives have the same perturbative
order, and the $H^1$ bound controls the vacuum $H^{1/2}$ seminorm.
Thus $k[g_\epsilon]=\epsilon^2 k[g_0]+O(\kappa\epsilon^4)$ and
$p_\epsilon=\epsilon^2 k[g_0]+O((1+\kappa)\epsilon^4)$.
The leading variances differ strictly by
Proposition~\ref{prop:boundary-sensor-law}. Their record difference cannot
vanish for all sufficiently small nonzero $\epsilon$.
Equation~\eqref{eq:sensor-energy-balance} gives the energy order.
\end{proof}

\subsection{A finite record returns to a physical selection}
\label{sec:sensor-selection-return}

The two compact-source branches of Eq.~\eqref{eq:compact-return-switch} can
be sensed by the same contact, without giving their coefficient values to the
program. Start the boundary clock at $u=4$, keep $\kappa=0.05$, and use the
$C^2$ compact profile
\begin{equation}\label{eq:finite-sensor-profile}
 \zeta(\tau)=\frac43\left[1-\left(\frac{\tau-4}{4}\right)^2\right]^3
 \quad (0\le\tau\le8),\qquad \zeta=0\quad\hbox{otherwise}.
\end{equation}
It has finite variance and energy; smooth approximants give the same limiting
quantities. Solving Eq.~\eqref{eq:boundary-sensor-ode} and integrating
Eq.~\eqref{eq:driven-sensor-variance} gives
\begin{center}
\begin{tabular}{@{}ccccc@{}}
\toprule
Original population $r$ & Initial $\mathcal C$ & Final $\mathcal C$
 & $\Pr(R=1)$ & Final $u$\\
\midrule
$0.5$ & $1.037498878$ & $1.1206720$ & $0.31748037$ & $57.481$\\
$0.9$ & $1.067496366$ & $1.1505122$ & $0.31771688$ & $52.458$\\
\bottomrule
\end{tabular}
\end{center}
The contact charge is exactly
$E_s^{(\tau)}=512(4/3)^2/(385\cdot8)=0.29552670$.
The corresponding null energies are $0.06577751$ and $0.06774638$.
These are deterministic numerical evaluations with converged quadratures,
not certified intervals. The single-bit total variation is about
$2.3650\times10^{-4}$, so one contact is a weak discriminator, with
equal-prior error about $0.49988175$. It is a nonzero physical return, not a
high-confidence reconstruction or an optimal sensing strategy. Many repeats
would need a new account of field correlations, preparation and accumulated
source; they cannot be declared independent by multiplying this record law.
Even ignoring all other operations, Eq.~\eqref{eq:sensor-cumulative-budget}
allows at most 190 contacts with this unchanged charge on either initial
branch. This is a necessary budget, not a 190-contact construction.

Retain $R$ and let it select the next coupled projection:
$P_0=(I+Z)/2$ if $R=0$ and $P_{\pi/2}=(I+X)/2$ if $R=1$.
Use the independent bit state $(I+0.8Z)/2$ and a later compact switch
$\zeta_{\rm next}$, scheduled by the same proper clock after $\tau=8$.
Its realised null smear is $g_{\rm next}(u(\tau))=\zeta_{\rm next}(\tau)$.
The program supplies that proper-time profile and selects its projection using
only $R$; it need not know $u(\tau)$ or $\mathcal C$.
The mean population coupled to the next field is
\begin{equation}\label{eq:sensor-next-source}
 r_{\rm next}=0.9-0.4\Pr(R=1),\qquad
 T_{uu,\rm next}^{\rm added}
       =r_{\rm next}\frac{(g_{\rm next}')^2}{4}.
\end{equation}
It is $0.77300785$ and $0.77291325$ in the two branches. Thus the earlier
selection changes a source, that source changes the boundary record law, and
the recorded bit changes both the next pointer algebra and the next mean source.
The controller uses no numerical dilaton input.
For example, use the profile in Eq.~\eqref{eq:finite-sensor-profile} rescaled
to amplitude $1/4$ and the proper-clock interval $9\le\tau\le11$.
Its added proper-clock source charge is
$r_{\rm next}\,512(1/4)^2/(385\cdot2)$. The source kernel in $u$ is the one
on each branch's actual mean trajectory; the population difference alone is
not a comparison of two complete fixed-shape stresses. The same trajectory
equations apply with $\dot\zeta^2$ replaced by
$r_{\rm next}\dot\zeta_{\rm next}^2$, continuing the causal mean solution
beyond readout.

There is a same-current subtlety in this return. On the union of the sensor and
later null intervals, take the input to have the vacuum restriction. A preceding
compact Weyl source in a disjoint earlier interval leaves this restriction
unchanged, since the relevant smeared currents commute. Conditioning on the
sensor bit applies $\cos Q_g$ or $\sin Q_g$ to that restriction. Both resulting
states are invariant under simultaneous current parity, so every later current
has zero conditional mean. The next bit is independently prepared, and the
linear stress cross term vanishes in each record branch, proving
Eq.~\eqref{eq:sensor-next-source} without resetting the field. Its conditional
covariance generally changes. The later field is not a fresh vacuum port, and
a vacuum or thermal gate-error formula cannot be carried over to it.

For a realised switch $\widetilde g$, define the vacuum seminorm
$\|g\|_0=\sqrt{k[g]}$. If $\|\widetilde g-g\|_0\le e$, then
\[
 |\widetilde p-p|\le |k[\widetilde g]-k[g]|
 \le 2\sqrt{k[g]}e+e^2.
\]
This prices the needed control of the actual smeared interaction; clock and
worldline errors must be included in $\widetilde g$, not only in its amplitude.
Two record laws remain distinct if the sum of their probability-error bounds
is below their stated separation. Section~\ref{sec:joint-source-certificate}
adds the independent source requirement, and
Section~\ref{sec:thermal-sensor-renewal} gives a finite thermal replacement.
The clock, contact and continuous exhaust with their complete stress remain
the joint apparatus task of Conjecture~\ref{conj:gravitating}.
The finite run consumes a prepared detector bit and an independently prepared
next bit. Resetting the record without retained side information must export
its binary entropy; it is not supplied by coherent field return.
\subsection{A joint certificate for records and gravitational sources}
\label{sec:joint-source-certificate}

A detector can be accurate as an instrument and inaccurate as a source.
The vacuum seminorm controls a half derivative of its switch, whereas the
injected energy controls a full derivative. This distinction is already
visible in a single contact, before questions of thermal reuse arise.

\begin{proposition}[Record accuracy does not control source energy]
\label{prop:record-source-gap}
Let $g_n(u)=n^{-3/4}e^{-u^2/2}\cos(nu)$, $n\ge1$, and
$E[g]=\int(g')^2du/4$. Then
\begin{align}
 k[g_n]&=\frac{n^{-3/2}}8
       [n\sqrt\pi\,\operatorname{erf}(n)+2e^{-n^2}],\nonumber\\
 E[g_n]&=\frac{n^{-3/2}\sqrt\pi}8
       [n^2+\tfrac12+\tfrac12e^{-n^2}].
 \label{eq:record-source-counterexample}
\end{align}
Its detector flip probability tends to zero while its source energy diverges.
Every member has finite absolute stress moments through order two.
\end{proposition}
\begin{proof}
The Fourier transform is the sum of two shifted Gaussians. Integrating its
squared modulus with weights $\omega$ and $\omega^2$ gives the two formulas.
Consequently $k[g_n]\sim\sqrt\pi/(8\sqrt n)$ and
$E[g_n]\sim\sqrt{\pi n}/8$. Gaussian decay gives the finite moments.
\end{proof}
This is a comparison on a fixed null-coordinate field. A member with too much
energy can leave a chosen gravitational operating domain. A fixed bandwidth
excludes this sequence, but that bandwidth is then part of the source
certificate. The controlled weak-switch result above uses stronger regularity
and is unaffected.

\begin{proposition}[Simultaneous record, source and response control]
\label{prop:joint-source-bound}
Write the realised switch as $g+h$, put $w(u)=1+u^2$, and assume
$g,h\in H^1$ with weighted square-integrable derivatives. Define
\[
 e_0=\sqrt{k[h]},\qquad A=\|\sqrt w\,g'\|_2,\qquad
 e_1=\|\sqrt w\,h'\|_2,\qquad D_s=\tfrac12Ae_1+\tfrac14e_1^2.
\]
For the independent vacuum input of the current detector,
\begin{align}
 |\widetilde p-p|&\le2\sqrt{k[g]}e_0+e_0^2,\nonumber\\
 \int w\,|\widetilde f-f|du&\le D_s,\qquad
 f=(g')^2/4,\quad \widetilde f=(g'+h')^2/4.
 \label{eq:joint-source-bound}
\end{align}
The second line bounds each of the three stress-moment errors by $D_s$.
For matched retarded data, at each common $u$ the running coefficients satisfy
$|\delta a|,|\delta b|,|\delta\gamma|\le\kappa D_s$, hence
\begin{align}
 |\delta\mathcal C|&\le
 \kappa D_s(2|a|+|b|+|\gamma|)+2\kappa^2D_s^2,\nonumber\\
 |\delta\Phi(u,v)|&\le
 \frac{\kappa D_s(1+|uv|+|u|+|v|)}{|1+uv|}.
 \label{eq:source-response-certificate}
\end{align}
For a single contact, even when the detector has a retained reference,
the half trace distance of the complete output from its ideal output is at
most $\sqrt{1-e^{-k[h]}}\le e_0$.
\end{proposition}
\begin{proof}
Cauchy--Schwarz applied to
$\widetilde f-f=g'h'/2+(h')^2/4$ proves the weighted bound.
The inequalities $1,|u|,u^2\le w$ cover both full and prefix moments;
the retarded coefficient equations then give the stated errors.
Expand $(a+\delta a)^2+(b+\delta b)(\gamma+\delta\gamma)$ and the
dilaton numerator for the response bounds.
For the final claim, purify the initial detector and reference.
The Weyl relation makes $U_g^\dagger U_{g+h}$ equal to
$e^{-iX_sQ_h}$ up to one scalar phase, since $X_s^2=I$.
Its vacuum overlap is $e^{-k[h]/2}$ independently of the detector state.
The pure-state trace-distance formula and partial-trace contraction give
the bound. This claim concerns the declared vacuum input, not an arbitrary
field left by earlier instruments.
\end{proof}

The coefficient bounds also control the clock on a regular window.
If both compared fixed-cutoff solutions have
$R_b\ge r_*>0$, $P\ge p_*>0$, and $R_b\le r^*$, put
$\Delta_C=\sup|\delta\mathcal C|$ and
$\Delta_P=\sup\kappa D_s(1+|u|)^2$ on a null interval of length $\ell$.
Matching the initial proper time gives
\[
 \sup|\delta\tau(u)|\le
 2L\ell\left(\frac{\Delta_C}{2r_*p_*}
                 +\frac{r^*\Delta_P}{p_*^2}\right).
\]
An inverse-clock bound follows by dividing by a positive lower bound on
$d\tau/du$. The interval and gap conditions must hold in both solutions.
They cannot be inferred from fidelity alone.

A finite timing tolerance has its own source condition. For a translation
error $h_\delta(u)=g(u-\delta)-g(u)$ and $g$ with the indicated second
derivative norms finite,
\begin{align}
 \sqrt{k[h_\delta]}&\le|\delta|\sqrt{k[g']},\nonumber\\
 \|\sqrt w\,h_\delta'\|_2&\le
 |\delta|\|\sqrt w\,g''\|_2+\tfrac12\delta^2\|g''\|_2.
 \label{eq:clock-source-tolerance}
\end{align}
The first inequality follows from $|e^{i\omega\delta}-1|\le|\omega\delta|$.
For the second, integrate $g''(u-s)$ between $0$ and $\delta$ and use
$\sqrt{1+u^2}\le\sqrt{1+(u-s)^2}+|s|$.
This is a null-clock translation certificate. A proper-clock error must be
pulled through its actual worldline map; a finite quantum clock also needs
control of its tails and correlations. Small arrival error alone does not
supply a continuum stress bound.

\subsection{Boundary energy, work and clock conversion}
\label{sec:boundary-work}

Choose the timelike Dirichlet boundary term, with the normal pointing toward
increasing dilaton,
\begin{equation}\label{eq:finite-boundary-action}
 I_\partial=\frac1\kappa\int_{\partial M}
         \sqrt{-h}\,\Phi(K-L^{-1})\,dt .
\end{equation}
Its lapse variation fixes the proper-time quasilocal energy
\[
 E_b=\frac1\kappa(\Phi_b/L-n^\mu\partial_\mu\Phi).
\]
This is the finite-cutoff JT energy with a stated counterterm, rather than
a total apparatus energy~\cite[Sec.~6]{Griguolo2026Cutoff}.
For example, in Gaussian normal coordinates $ds^2=dr^2-N^2dt^2$,
integrating the bulk curvature term by parts and adding
Eq.~\eqref{eq:finite-boundary-action} leaves the lapse boundary term
$N(\partial_r\Phi-\Phi/L)/\kappa$. Its variation is $-E_b\delta N$.
The dilaton boundary momentum is $(K-L^{-1})/\kappa$.

\begin{proposition}[Energy and work on a running level curve]
\label{prop:boundary-energy-work}
In the traceless single-chiral source completion let
$\Phi_b=\Phi_b(\tau)$ be a smooth prescribed boundary value. On a future
timelike segment put
\[
 R_b^2=\Phi_b^2-\mathcal C>0,\quad
 S_b=\sqrt{R_b^2+L^2\dot\Phi_b^2},\quad
 y=S_b-L\dot\Phi_b>0,\quad q_\tau=f\dot u^2.
\]
Then
\begin{align}
 \dot u&=\frac{P}{2Ly},&
 \dot{\mathcal C}&=2\kappa Lyq_\tau,&
 \dot R_b&=\frac{\Phi_b\dot\Phi_b-\kappa Lyq_\tau}{R_b},\nonumber\\
 E_b&=\frac{\Phi_b-S_b}{\kappa L},&
 K&=\frac{\Phi_b+L^2\ddot\Phi_b+\kappa L^2q_\tau}{LS_b},&
 \dot E_b&=q_\tau+\frac{L^{-1}-K}{\kappa}\dot\Phi_b .
 \label{eq:boundary-work-balance}
\end{align}
At a fixed cutoff these reduce to
$E_b=(\Phi_b-R_b)/(\kappa L)$ and $\dot E_b=q_\tau$.
On the source-free stationary family,
$dE_b=T_b\,dS_h$ at fixed $\Phi_b$, with
$T_b=R/(2\pi LR_b)$.
\end{proposition}
\begin{proof}
Writing the boundary as $v=V(u)$, its dilaton equation and the induced
metric give
$1=4L^2R_b^2\dot u^2/P^2-4L^2\dot\Phi_b\dot u/P$.
Its positive root is the first equation.
Use $\mathcal C'=\kappa fP$ for the next two.
The gradient identity gives $Ln^\mu\partial_\mu\Phi=S_b$.
Along the unit tangent, the JT equation reads
$\ddot\Phi_b=-\Phi_b/L^2-\kappa q_\tau+K S_b/L$.
This gives $K$; differentiating $E_b$ proves the work balance.
The stationary first law follows by differentiating
$E_b=(\Phi_b-\sqrt{\Phi_b^2-R^2})/(\kappa L)$ and
$S_h=S_0+2\pi R/\kappa$.
\end{proof}

For example, retaining a constant gap $R_b=r>0$ under positive flux requires
\[
 y=\frac r{\sqrt{1+2\kappa L^2q_\tau/\Phi_b}},\qquad
 \dot\Phi_b=\frac{\kappa Lyq_\tau}{\Phi_b}.
\]
These equations give a nonempty family on regular finite patches.
Increasing the cutoff carries the work term in
Eq.~\eqref{eq:boundary-work-balance}; it is not an uncharged extension of the
fixed-cutoff budget. The boundary datum is still externally prescribed.
A physical support must produce it and carry its own stress.

The same calculation fixes the conversion to the running horizon clock:
\begin{equation}\label{eq:boundary-clock-conversion}
 z:=\frac{d\tau}{d\eta}=\frac{Ly}{R},\qquad
 q_\eta=z^2q_\tau .
\end{equation}
For a fixed cutoff $y=R_b$, and the stationary temperature is
$T_b=(2\pi z)^{-1}$. During a driven contact this is not a global KMS
identification. If a local thermal-current regime has separately been
justified, entropy flux transforms as $s_\eta=z s_\tau$ and its actual
input temperature must satisfy $zT_\tau=1/(2\pi)$ to use
Proposition~\ref{prop:horizon-floor}. Changing coordinates does not prepare
that input.

\subsection{A thermal detector can be reused without refreshing its field}
\label{sec:thermal-sensor-renewal}

A finite Gibbs swap provides a physical replacement channel for the detector.
Its limitation is equally useful: reset accuracy, retained correlations and
heat cannot be specified independently.

\begin{proposition}[Reference-uniform thermal replacement]
\label{prop:thermal-detector-reset}
During sensing the detector Hamiltonian is degenerate. For reset raise its
excited level to $\Delta>0$ and supply an independent auxiliary bit in
$\tau_\beta=\operatorname{diag}(1-\epsilon,\epsilon)$, with
$\epsilon=(1+e^{\beta\Delta})^{-1}$. A resonant SWAP, followed by lowering
the detector gap, gives for every joint detector--reference state
\[
 \rho_{SA}\longmapsto\tau_\beta\otimes\rho_A .
\]
Its half-diamond distance from pure-blank replacement is exactly $\epsilon$.
If $p=\Pr(S=1)$ just before reset, the auxiliary heat and net gap-modulation
work are both
\begin{equation}\label{eq:thermal-sensor-reset-charge}
 Q_{\rm reset}=W_{\rm gap}=\Delta(p-\epsilon).
\end{equation}
A finite resonant swap can use
$H_{\rm sw}(\tau)=\chi(\tau)(I-\mathrm{SWAP})$ with
$\int\chi\,d\tau=\pi/2$; it commutes with the equal-gap free Hamiltonian.
Its net switch work is zero when $\chi$ vanishes at both endpoints.
\end{proposition}
\begin{proof}
The swap transfers the whole previous detector state and its reference
correlations to the auxiliary, not just its populations.
The distance between the two replacement states is $\epsilon$, for every
reference, proving the channel statement and attainability.
Raising the gap costs $\Delta p$, the resonant swap conserves total bare energy,
and lowering it returns $\Delta\epsilon$.
The outgoing auxiliary has excitation probability $p$.
The expectation of $I-\mathrm{SWAP}$ is constant during the specified
resonant contact, so its switching work integrates to zero.
\end{proof}
This replacement process is the standard swap benchmark for Landauer
accounting~\cite[Example~1 and Sec.~3.3]{ReebWolf2014}.
The auxiliary must be prepared and transported. The old state remains
correlated with retained matter; it has not disappeared.
If a copy of the outcome is kept, that copy is part of the maintained state
in Eq.~\eqref{eq:entropyaccount}.
At finite temperature this finite device never supplies an exact pure blank.

\begin{proposition}[A bounded-heat Gibbs ladder]
\label{prop:thermal-reset-ladder}
Fix $\beta>0$, a target $0<\epsilon<1/2$, and $m$ finite thermal stages.
First swap the degenerate detector with an independent zero-gap Gibbs bit
$I/2$. Put $x_*=\ln[(1-\epsilon)/\epsilon]$,
$x_j=jx_*/m$ and $p_j=(1+e^{x_j})^{-1}$.
At stage $j$ raise the detector gap from $x_{j-1}/\beta$ to $x_j/\beta$
and resonantly swap with an independent Gibbs bit of that gap.
Finally lower its gap to zero. The complete replacement is exactly
$\tau_\beta\otimes\rho_A$, with target excitation $\epsilon$, and
\begin{equation}\label{eq:thermal-ladder-heat}
 \beta Q_{\rm ladder}
 =\sum_{j=1}^m x_j(p_{j-1}-p_j)
 \le\ln2-h(\epsilon)+\frac{x_*}{m}(\tfrac12-\epsilon).
\end{equation}
The net gap-modulation work equals this heat.
Choosing $m\ge x_*/(2\eta)$ gives $\beta Q_{\rm ladder}\le\ln2+\eta$
for every initial detector state and every retained reference.
\end{proposition}
\begin{proof}
The first swap transfers all old correlations to its output.
Each later swap leaves an independent Gibbs detector.
The heat formula follows from the thermal populations, and summation by
parts gives the same net work, including final gap lowering.
The sum is a right-endpoint approximation to
$\int_0^{x_*}x[-p'(x)]dx=\ln2-h(\epsilon)$;
its excess is at most $(x_*/m)(p_0-p_m)$.
\end{proof}
This is an explicit finite swap construction of the type used in
Ref.~\cite[Sec.~6]{ReebWolf2014}, with a distribution-independent starting
step and a bound on its stages. It consumes $m+1$ prepared Gibbs bits,
including the zero-gap port, and raises the maximum gap to $x_*/\beta$.
The finite reset does not hide a pure auxiliary. The first port can leave
with less entropy than it had initially; it must be included along with all
later outputs and their correlations. For a fixed supplied sequence of
operations, $N$ replacements differ from the pure-blank retained-system
process by at most the sum of their $\epsilon_i$, including arbitrary
intervening instruments on the field. Thus high reset accuracy need not
cause divergent heat, but the number of stages and the maximum gap grow.
If the replacement changes the realised geometry or timing of those
operations, their source and clock errors require the separate joint
certificate of Section~\ref{sec:joint-source-certificate}.
The timing, Gibbs preparation, transport and capacity required to execute
these stages are additional resources.


Here is an exact two-read law on one field. Let $g_1,g_2$ have ordered,
disjoint compact supports, put $k_i=\langle Q_i^2\rangle$ and
$c_{12}=\langle Q_1Q_2\rangle$ in the initial vacuum, and reuse the detector
after the above swap. Both contacts start with its state $\tau_\beta$.
Read $Z_s$ and use it for the subsequent command before each reset.
A nondemolition dephasing of the pointer can be implemented by a controlled
$Z_s$ using a separate degenerate Gibbs bit $I/2$. Those ports are supplied
resources, not pure memories. Since $[Q_1,Q_2]=0$, the joint outcomes obey
\begin{equation}\label{eq:correlated-thermal-reads}
\begin{split}
 \Pr(R_1=r,R_2=t)=\frac14\bigl[&
 1+\alpha\{(-1)^re^{-2k_1}+(-1)^te^{-2k_2}\}\\
 &+\alpha^2(-1)^{r+t}e^{-2(k_1+k_2)}\cosh(4c_{12})\bigr].
\end{split}
\end{equation}
Here $\alpha=1-2\epsilon$.
Indeed each conditional pointer probability at a fixed $Q_i$ is
$[1+\alpha(-1)^r\cos(2Q_i)]/2$.
The commuting Gaussian characteristic function then gives the displayed law.
The same field remains conditioned by the first read: thermal replacement of
the detector does not restore a vacuum field. On a stationary test background
the contrast of two invariant-dependent single-contact laws is multiplied
by $\alpha$, so it stays nonzero for $\epsilon<1/2$.
The weak-source proof of Proposition~\ref{prop:weak-boundary-sensor}
also survives this factor.

A retained body makes the two commands physically consequential.
Prepare the body once in $(I+0.8Z)/2$. After read $r$, couple it through
$P_0=(I+Z)/2$ or $P_1=(I+X)/2$ to an independent vacuum control current.
The first source population is $\rho_1=0.9-0.4\Pr(R_1=1)$.
If that contact's realised variance is $k_c$, the first $X$ projection
attenuates the body's $Z$ component by $e^{-k_c/2}$.
With $p_{rt}$ denoting Eq.~\eqref{eq:correlated-thermal-reads}, its second
source population is therefore
\begin{equation}\label{eq:held-body-second-source}
 \rho_2=\tfrac12+0.4\{p_{00}+e^{-k_c/2}p_{10}\}.
\end{equation}
This uses the body left by the first command, not a newly prepared copy.

For a finite boundary run take $L=1$, $\Phi_b=3$, $\kappa=.01$,
initial $(u,a,b,\gamma)=(1,1,0,0)$, $\epsilon=.01$ and $\beta=1$ in
the proper clock. Use the profile
$A[1-(2(\tau-\tau_0)/d-1)^2]^3$ on an interval of duration $d$.
The two sensor contacts have amplitude $.5$ on $[0,2]$ and $[4,6]$;
the two selected control contacts have amplitude $.3$ on $[2.5,3.5]$
and $[6.5,7.5]$. The gaps leave finite windows for readout and reset.
Integrating the same mean boundary equations gives
\[
 k_1=.068080013,\quad k_2=.068082993,\quad
 c_{12}=-.000926643,\quad k_c=.024462001,
\]
\[
 \rho_1=.871049781,\quad \rho_2=.870722307,\quad
 \mathcal C_{\rm final}=1.030570773 .
\]
The proper-clock coupling-source charge is $.540937293$.
The two reset auxiliaries separately absorb $.573257937$, using
$\Delta=\ln99$; this heat has not been inserted as a chiral bulk source.
The two outgoing reset bits have joint entropy $.519488204$, against
$.112003069$ initially, and relative entropy $.165772802$ to the product
Gibbs reference. Thus their heat obeys exactly
$\beta Q=\Delta S(B_1B_2)+D(\rho_{B_1B_2}\Vert\tau_\beta^{\otimes2})$.
Other output ports and retained matter belong to the full entropy account.
The source charge, reset heat and apparatus work are distinct entries.

Starting instead at $a=1.4$ with the other initial coefficients unchanged
changes both recorded laws and returned source populations on the same
proper schedule. The detector and body are reused in both runs.
The completion transmits the sensor current and two independent control
currents into the bulk. The reset ports, clock and support still require a
conserved transport completion; their cost is not removed by the displayed
bulk balance. A long sequence of simple Gibbs swaps is also expensive:
choosing $\epsilon_N=\varepsilon/N$ to certify total blank-replacement error
at most $\varepsilon$ requires
$\beta\Delta_N=\ln(N/\varepsilon-1)$.
If detector excitation probabilities stay above a fixed positive constant,
its heat per reset grows as $\log N$. This particular construction cannot
attain the fixed-rate sharp-exhaust conjecture.
The ladder of Proposition~\ref{prop:thermal-reset-ladder} avoids this heat
divergence, but fitting its stages and transported thermal units into a
fixed capacity and operating rate remains necessary.


\section{Running horizons, source states and clocks}
\label{app:running-clocks}

\subsection{Source sorting and tracking-clock support}
\label{app:source-clock-support}

The finite-carrier sorting statement remains
Proposition~\ref{prop:pointer-sort}, with its source and old-reference
permissions unchanged.
\begin{proof}[Proof of Proposition~\ref{prop:pointer-sort}]
The source enters Q-RAIF~C's proof in three places: concentration
of input eigenlabel strings, the source's fixed neutral charge, and the
final source--carrier entropy identity. Its detailed depolarizing
eigenvalues are not otherwise used. Each $\rho_p$ has rank two in the
same neutral quartet sector and finite information variance at fixed
$q,\epsilon$. Independent eigenlabel
strings have information variance $O(L)$, including a nonidentical
product over the two-symbol alphabet. The sorter's subextensive
typical-label margin, thermal-sector state count and physical
neutral-sector unitary compiler therefore apply with entropy
$\sum_jh_{p_j}$. The source has zero total charge, leaving the carrier's
sector weights and sublinear sector correction unchanged. Completing
the eigenlabel injection to a sector unitary keeps all bad labels and
acts as the identity on external references. The source--carrier entropy
identity and controlled full-carrier energy tails give the deficit and
slack limits through the Gibbs identity. Trace error alone would not
suffice. Uniform local reset error alone does not establish the supplied
product marginal for an adaptive history.
\end{proof}

Equation~\eqref{eq:loop-moving-clock} gives the actual clock, not
a static-wall substitution. For a tracking datum $R_b=\rho R$,
$\Phi_b=\alpha R$, $\alpha=\sqrt{1+\rho^2}$ and
$\delta=R^{-1}dR/d\eta$, it gives
\begin{equation}\label{eq:tracking-clock-exact}
 z=L\sqrt{\rho^2-2\alpha\delta},\qquad \rho^2>2\alpha\delta.
\end{equation}

\subsection{The running horizon and its entropy}
\label{sec:running-horizon}

Long-time entropy limits need an operating geometry, not just a source integral.
Positive steady flux can exhaust a fixed $u$ chart. At the boundary $v=-1/u$
the numerator of the retarded dilaton is
\[
 N_b(u)=2H_0-\frac{\kappa}{u}\int^u(u-z)^2F(z)\,dz.
\]
For a constant flux $F_0>0$ switched on at $u=0$, this becomes
$2H_0-\kappa F_0u^2/3$ and vanishes at
$u_*=\sqrt{6H_0/(\kappa F_0)}$. A device at a fixed exterior position cannot
therefore be assumed to run at fixed $u$ frequency indefinitely. Finite
observation windows and families with larger operating domains are different
from an infinite-time protocol in one fixed background.

There is an exact horizon balance inside the retarded solution. In this subsection
we use $\hbar=k_B=1$ and the classical Wald entropy, not an event-horizon
generalized entropy. Write the running coefficients as
\begin{align}
 a(u)&=H_0+\kappa\int^u zF(z)\,dz,&
 b(u)&=-\kappa\int^u F(z)\,dz,&
 \gamma(u)&=\kappa\int^u z^2F(z)\,dz,\nonumber\\
 \Phi&=\frac{a+bu+\gamma v-auv}{1+uv},&
 \mathcal C(u)&=a^2+b\gamma,& R(u)&=\sqrt{\mathcal C(u)}.
 \label{eq:running-coeff}
\end{align}
The lower limits and any source-free initial coefficients are fixed by the
matching data. Put $P=bu^2+2au-\gamma$ and $\xi=P/(2R)$.

\begin{proposition}[Exact marginal-horizon balance]\label{prop:horizon-balance}
On a patch with $\mathcal C>0$, $a+R\ne0$, $P>0$ and $1+uv_h\ne0$, the curve
\[
 v_h(u)=\frac{b}{a+R}
\]
is a future outer marginal horizon: $\partial_u\Phi=0$,
$\partial_v\Phi<0$ and $\partial_v\partial_u\Phi<0$ there. Its dilaton and Wald
entropy obey
\begin{equation}\label{eq:wald-balance}
 \Phi_h=R,\qquad S_h=S_0+\frac{2\pi R}{\kappa},\qquad
 \frac{dS_h}{du}=2\pi F\xi,\qquad
 \frac{dS_h}{d\eta}=2\pi F\xi^2,\quad d\eta=\frac{du}{\xi}.
\end{equation}
For $F\ge0$ this horizon is spacelike or null and its entropy is nondecreasing.
\end{proposition}
\begin{proof}
The coefficient derivatives $a'=\kappa uF$, $b'=-\kappa F$ and
$\gamma'=\kappa u^2F$ cancel from the first dilaton derivative. With $D=1+uv$,
\[
 \Phi_u=\frac{b-2av-\gamma v^2}{D^2},\qquad
 \Phi_v=-\frac{P}{D^2}.
\]
Substitution of $v_h$ gives $\Phi_h=R$ and
$\Phi_{uv}|_h=-2R/(1+uv_h)^2$. Also
$\mathcal C'=\kappa FP$ and
$v_h'=-\kappa F(1+uv_h)^2/(2R)$. The metric then makes the tangent spacelike
for $F>0$ and null for $F=0$. Finally $R'=\kappa F\xi$ and $\kappa=8\pi G$
give Eq.~\eqref{eq:wald-balance}.
\end{proof}

The identification of this classical marginal horizon and its area law is a
standard JT result~\cite[Sec.~2.3.1]{Moitra2020}. Here it assigns an actual
horizon observable to the running invariant used in the source account.
Negative quantum excess stress need not be monotone, and replacing $S_h$ by
$S_{\rm gen}$ would require the matter entanglement contribution and a separate
quantum argument.

\begin{proposition}[Conditional entropy-export cost at the horizon]
\label{prop:horizon-floor}
Suppose the charged heat in the rapidity clock is exactly
$\dot Q_\eta=F\xi^2$, and the exhaust satisfies
Proposition~\ref{prop:chiral-floor} in that same clock at $T_H=1/(2\pi)$.
On a window of rapidity duration $\tau$, let $X_\tau$ be the net exported
maintenance entropy. In the extensive rate regime,
\[
 \frac{\Delta S_h}{\tau}\ge\frac{X_\tau}{\tau}
       +\frac{6}{c_{\rm ex}}\left(\frac{X_\tau}{\tau}\right)^2+o(1).
\]
Writing bars for the corresponding mean rates gives
\begin{equation}\label{eq:horizon-floor}
 \overline{\dot S_h}\ge \dot s+\frac{6}{c_{\rm ex}}\dot s^2.
\end{equation}
The same pointwise formula requires a justified local-current limit; it does
not follow for arbitrary fluctuating packets. A mean horizon-entropy budget
$\overline{\dot S_h}\le B$ permits at most
\[
 \dot s\le \frac{c_{\rm ex}}{12}
       \left(\sqrt{1+\frac{24B}{c_{\rm ex}}}-1\right).
\]
\end{proposition}
\begin{proof}
Multiply Eq.~\eqref{eq:chiral-floor} at $T=T_H$ by $2\pi$, use
the integrated Eq.~\eqref{eq:wald-balance} and the long-window Gibbs account,
then invert the increasing quadratic.
\end{proof}

This transfers the entropy removed from the full maintained state to a specified
gravitational entropy, including the finite-rate excess. The heat-routing
condition is substantive: incoming nonequilibrium energy, other stress sources,
charge transfer and the clock's own entropy cannot be omitted. Nor does a change
of clock alone prepare a thermal reservoir.

\begin{proposition}[Schwarzian of the running clock]
\label{prop:clock-schwarzian}
On a smooth running patch put $q=F\xi^2$ and
$\delta=R^{-1}dR/d\eta=\kappa q/R$. Then
\begin{equation}\label{eq:clock-general}
 \{u;\eta\}=-\frac12-\frac{d\delta}{d\eta}-\frac{\delta^2}{2}.
\end{equation}
For constant $q>0$, $R(\eta)=R_0+\kappa q\eta$ and
\begin{equation}\label{eq:clock-schwarzian}
 \{u;\eta\}=-\frac12+\frac{(\kappa q)^2}{2R(\eta)^2}.
\end{equation}
\end{proposition}
\begin{proof}
Write $j=(bu+a)/R$, and let a dot denote $d/d\eta$.
The coefficient equations and $h^2-2bR\xi=R^2$, with $h=bu+a$, give
\[
 \frac{\dot\xi}{\xi}=j-\delta,\qquad
 \dot j=\frac{j^2-1}{2}-j\delta.
\]
Substitute these into $\{u;\eta\}=(\dot\xi/\xi)^{\displaystyle\cdot}
-(\dot\xi/\xi)^2/2$. For constant $q$, $\dot\delta=-\delta^2$,
which recovers Eq.~\eqref{eq:clock-schwarzian}.
\end{proof}

Constant positive heat current changes the vacuum clock correlation, not merely
the horizon size. Positive growth itself, however, does not rule out an exact
thermal reference. The following two histories solve the clock condition.

\begin{proposition}[Growing horizons with exact geometric KMS clocks]
\label{prop:running-kms}
Use natural units and the prescribed-mean-source JT equations. An independent
chiral conformal reference sector is in its vacuum; its state is not the state
of the driven matter supplying $F$. On every finite smooth operating window,
the following positive-flux profiles admit the initial matching
$(a,b,\gamma)=(R_0,0,0)$ at $\eta=0$, $u=u_0>0$:
\begin{enumerate}
\item Constant fractional growth $\delta_0>0$:
$R=R_0e^{\delta_0\eta}$, $q=\delta_0R/\kappa$. The reference pulled back
to the $\eta$ clock is geometric KMS at $T=\sqrt{1+\delta_0^2}/(2\pi)$.
\item Fixed rapidity temperature: for $\eta_0>0$,
\begin{equation}\label{eq:fixed-temperature-growth}
 R=R_0(1+\eta/\eta_0)^2,\qquad
 q=\frac{2R_0(\eta+\eta_0)}{\kappa\eta_0^2}.
\end{equation}
The same reference is exactly geometric KMS at $T_H=1/(2\pi)$ despite
the growing horizon.
\end{enumerate}
Here KMS refers to the pulled-back reference algebra and its complete
$\eta$-translation dynamics. A finite experiment samples its restriction.
\end{proposition}
\begin{proof}
Both profiles make Eq.~\eqref{eq:clock-general} constant, $-s^2/2$,
with $s=\sqrt{1+\delta_0^2}$ in the first case and $s=1$ in the second.
Let $\delta(0)=\delta_0$ also in the second case, where
$\delta_0=2/\eta_0$. An explicit solution with the required initial
coordinate and derivatives is
\begin{equation}\label{eq:running-mobius-clock}
 u(\eta)=u_0\left[1+\frac{e^{s\eta}-1}
 {s+m(e^{s\eta}-1)}\right],\qquad
 m=\frac{s-1+\delta_0}{2}.
\end{equation}
It has $\xi=\dot u>0$, $\xi(0)=u_0$ and
$\dot\xi(0)/\xi(0)=1-\delta_0$. Set $j=\dot\xi/\xi+\delta$ and
\[
 b=\frac{R(j^2-1)}{2\xi},\qquad a=Rj-bu,\qquad
 \gamma=bu^2+2au-2R\xi.
\]
The constant-Schwarzian equation gives
$\dot j=(j^2-1)/2-j\delta$.
Differentiation then gives
$\dot b=-\kappa q/\xi$, $\dot a=\kappa uq/\xi$ and
$\dot\gamma=\kappa u^2q/\xi$.
Thus these are the retarded coefficient equations with
$F=q/\xi^2$, $a^2+b\gamma=R^2$ and $P=2R\xi$.
At zero they give $(R_0,0,0)$; $a+R$ and the horizon denominator remain
nonzero on a sufficiently short initial operating patch, which can be
continued while those patch conditions hold.

Equation~\eqref{eq:running-mobius-clock} is a M\"obius transform of
$e^{s\eta}$. M\"obius invariance of the reference vacuum and the geometric
KMS construction therefore give inverse temperature $2\pi/s$
\cite[Sec.~2.8, Proposition~2.4 and Theorem~2.5]{Camassa2012KMS}.
For a free current the complete Gaussian correlation test is explicit:
if $\langle j(u)j(u')\rangle=-C_j/(u-u'-i0)^2$, then
\[
 \langle \xi(\eta)j(u(\eta))\,\xi(\eta')j(u(\eta'))\rangle
 =-\frac{C_js^2}{4\sinh^2[s(\eta-\eta'-i0)/2]}.
\]
This is stationary and obeys the KMS boundary condition; the conformal-net
argument establishes the state on bounded observables, not just on one
stress expectation.
\end{proof}

These are actual nonstationary solutions of the mean-source equations, rather
than frozen-background substitutions. The reference sector's thermality is a
kinematic input that can now be made compatible with growth. It is not a
thermal exhaust realization: the driving sector can be coherent, its state
is not KMS merely because the independent reference is, and preparation,
heat transport and the physical clock still need their own account. In the
second profile $q$ grows linearly, so it does not realize the fixed mean
export in Conjecture~\ref{conj:gravitating}. As $\eta\to\infty$, $u$ approaches
a finite endpoint and $F=q/\xi^2$ diverges. Each finite smooth window has
finite source energy; these formulas assert no finite-resource infinite run.


\subsection{The actual thermal source changes its clock reference}
\label{sec:thermal-source-clock}

A thermal spectator field does not make the exhaust thermal. Here the
state condition is imposed on the source field itself. Work with the
vacuum-subtracted chiral flux of the preceding JT construction, with the
state-independent trace background absorbed consistently into the
renormalized JT parameters. A complementary chirality may remain in its
reference state. This is a semiclassical conformal-matter completion,
not a microscopic construction of a thermal emitter.

For the inner patch let $u=f(v_-)$ be its Poincar\'e null coordinate.
The coordinate change
$\rho=\Phi/f'-L^2f''/(f')^2$ gives
\[
 \{u;v_-\}=-\frac{\mathcal C_-}{2L^4}.
\]
Specify the source state as the pullback of a geometric KMS state at
temperature $T_o$ in $\eta$, with $T_H=1/(2\pi)$ as the comparison
temperature. Define the charged excess current and entropy current by
\begin{equation}\label{eq:thermal-source-excess}
 q=\frac{\pi c}{12}(T_o^2-T_H^2)>0,\qquad
 s=\frac{\pi c}{6}(T_o-T_H),\qquad
 2\pi q=s+\frac{6s^2}{c}.
\end{equation}
This state condition is imposed on the finite operating patch.
Section~\ref{sec:local-source-contact} gives a local preparation route
from supplied thermal reservoirs and also treats nonthermal outputs;
the reservoir and contact's complete material realization remain inputs.

The physical stress is a tensor. The Schwarzian enters when the state is
expressed relative to the Poincar\'e vacuum, through the chiral stress
transformation law~\cite[Eq.~(2.14)]{Hollands2019}:
\begin{align}\label{eq:actual-thermal-source}
 F_{vv}&=\left(\frac{d\eta}{dv_-}\right)^2
              \left(\frac{c}{48\pi}+q\right)
       -\frac{c}{24\pi}\{\eta;v_-\}
       +\frac{c}{24\pi}\{u;v_-\}\nonumber\\
       &=\frac{R^2q}{L^4}-\frac{c}{24\pi}\{\eta;v_-\}.
\end{align}
The last term in the first line is essential. On a static patch, the
$T_H$ thermal state is the zero-flux vacuum reference; retaining only its
positive thermal energy would falsely insert a source. Conversely, on a
growing patch that same running-clock reference need not have zero
physical stress.

\begin{proposition}[State--clock--source consistency]
\label{prop:thermal-source-clock}
Put $a=\kappa c/(24\pi)$ and $\delta=R'/R$, where primes in this
proposition denote $\eta$ derivatives. The actual source state and the
JT equation are jointly satisfied by
\begin{equation}\label{eq:thermal-source-clock}
 R'=R\delta,\qquad
 \delta'=\frac{\kappa q-R\delta}{a}-\frac{\delta^2}{2}.
\end{equation}
The physical flux in this clock and the horizon entropy account are
\begin{align}
 F_\eta&=q+q_{\rm ref}=R'/\kappa,&
 q_{\rm ref}&=-\frac{c}{24\pi}(\delta'+\delta^2/2),\nonumber\\
 \Delta S_{h,-}&=2\pi q\Delta\eta
       -\frac c{12}\Delta\delta
       -\frac c{24}\int\delta^2d\eta.
                    \label{eq:thermal-clock-entropy}
\end{align}
For $R_0^2>2a\kappa q$ and
$\kappa q/(2R_0)\le\delta_0\le2\kappa q/R_0$, the source--clock ODE
has a global positive solution in $\eta$ with
\begin{equation}\label{eq:thermal-clock-growth}
 \frac{\kappa q}{2R}\le\delta\le\frac{2\kappa q}{R},\quad
 R=R_0+\kappa q\eta+O(1),\quad
 \int_0^\infty\delta^2d\eta\le\frac{8\kappa q}{R_0}.
\end{equation}
Thus $\Delta S_{h,-}/\Delta\eta\to2\pi q=s+6s^2/c$.
This global statement concerns the source and clock, not the lifetime
of a finite material energy reserve.
\end{proposition}
\begin{proof}
Since $dv_-/d\eta=L^2/R$, its Schwarzian is
$\{v_-;\eta\}=-\delta'-\delta^2/2$.
The chain rule in Eq.~\eqref{eq:actual-thermal-source} gives
$F_\eta=q-c(\delta'+\delta^2/2)/(24\pi)$.
Together with $\partial_v\mathcal C_-=2\kappa L^2F_{vv}$ this is
Eq.~\eqref{eq:thermal-source-clock}; integration gives
Eq.~\eqref{eq:thermal-clock-entropy}.
At the upper boundary $\delta=2\kappa q/R$, the derivative of
$\delta-2\kappa q/R$ is
$-\kappa q/a+2(\kappa q)^2/R^2<0$.
At the lower boundary the derivative of
$\delta-\kappa q/(2R)$ is
$\kappa q/(2a)+(\kappa q)^2/(8R^2)>0$.
The strip is invariant and $R'\ge\kappa q/2$.
Its upper bound gives
$\int\delta^2\le\int_0^\infty4(\kappa q)^2
 (R_0+\kappa q\eta/2)^{-2}d\eta$.
Finally $R'+a\delta'+a\delta^2/2=\kappa q$ gives the bounded correction
to linear growth. All variables stay finite on finite $\eta$ intervals.
\end{proof}

The charged heat $q$ and physical source $F_\eta$ differ by an explicitly
accounted reference stress. Proposition~\ref{prop:horizon-floor} remains
conditional on their equality; it cannot be applied pointwise to $q$
after discarding $q_{\rm ref}$. Here the reference correction is
nonextensive, so the sharp long-time rate is recovered by the stronger
state--clock calculation. A finite patch does not establish a globally
finite-energy state across a limiting Poincar\'e coordinate endpoint.

For any fixed finite proper-time window, nonempty operating solutions
also follow without relying on the numerical table. At $q=0$, start
from the positive static support, $e>0$, $R_0>0$ and $\delta=0$.
This is a regular equilibrium of the combined equations. Continuous
dependence on $q$, with $\delta_0=\kappa q/R_0$, preserves positive
mass, energy reserve and exterior domain on that window for sufficiently
small $q>0$. This is not a fixed-rate infinite-lifetime statement.

Combining Eq.~\eqref{eq:thermal-source-clock} with
Eq.~\eqref{eq:radiating-wall} uses
$P_{\rm out}=R'/[\kappa z^2]$ and $P_{\rm in}=0$.
For the support above, initial $e=2.34$ and
$\delta_0=\kappa q/R_0$, all three examples below run for proper
duration 80 with positive reserves and $\Phi>R$:
\begin{center}
\begin{tabular}{cccc}
\hline
$q$ & $\eta_{\rm end}$ & $\Phi_{\rm end}$ & $e_{\rm end}$\\
\hline
.005 & 85.53862 & 3.003858 & 1.882696\\
.010 & 85.60190 & 3.007607 & 1.424039\\
.020 & 85.73040 & 3.015345 & .502569\\
\hline
\end{tabular}
\end{center}
The mass-junction residual is below $1.2\times10^{-12}$ on the checked
grid. These solve the moving mean support and the state condition together.
The internal reserve's decline is part of the solution. It has not been
derived from the seven-qubit core's local release Hamiltonian.

\subsection{A local thermal contact supplies the actual source}
\label{sec:local-source-contact}

The source condition can be replaced by a scattering calculation. Use
two complex chiral fermion ports, each with $c=1$, at a local conformal
contact. This is the doubled free-fermion defect of
Bernard, Doyon and Viti~\cite[Sections~3 and 5]{BernardDoyonViti2015}.
The zero-charge input streams are jointly independent and thermal at
$T_H$ and $T_o$ in the same supplied $\eta$ frame. Their fields mix by
$S_\theta=\cos\theta I-i\sin\theta\sigma_x$ at the contact. Put
$\chi=\sin^2\theta$. For positive frequency, the source occupation is
\begin{equation}\label{eq:local-source-occupation}
 n_\chi(\omega)=(1-\chi)n_{T_H}(\omega)+\chi n_{T_o}(\omega),\qquad
 n_T(\omega)=\frac1{1+e^{\omega/T}}.
\end{equation}
The actual output is hot KMS for $\chi=1$. For $0<\chi<1$ it is
stationary but nonthermal: its high-frequency tail has coefficient
$\chi$ rather than one. Its return-port cross covariance is nonzero.
These are retained field correlations, not fresh independent thermal
maps. Q-RAIF~C derives the complete contact law and its energy
account~\cite[Section on a local thermal contact]{QRC}.

\begin{proposition}[Stationary transfer without a KMS output]
\label{prop:nonthermal-source-clock}
For a constant contact angle set
\begin{equation}\label{eq:local-source-charge}
 q_\chi=\frac{\pi\chi}{12}(T_o^2-T_H^2),\qquad
 s_\chi=\frac1\pi\int_0^\infty
       [h(n_\chi)-h(n_{T_H})]d\omega.
\end{equation}
Here $h(x)=-x\ln x-(1-x)\ln(1-x)$.
The state--clock equations, positive growing family and finite reference
correction in Proposition~\ref{prop:thermal-source-clock} hold with
$c=1$ and $q=q_\chi$. Their asymptotic horizon rate is
\begin{equation}\label{eq:nonthermal-horizon-rate}
 \frac{\Delta S_{h,-}}{\Delta\eta}\longrightarrow2\pi q_\chi
       \ge s_\chi+6s_\chi^2.
\end{equation}
The inequality is strict for $0<\chi<1$ and $T_o>T_H$.
\end{proposition}
\begin{proof}
The stationary output's stress in $\eta$ is
$1/(48\pi)+q_\chi$. Its coordinate/state transformation has the same
central charge and projective-vacuum subtraction as
Eq.~\eqref{eq:actual-thermal-source}. That derivation uses only this
stress expectation, not the KMS property. The source--clock ODE and
its invariant strip therefore follow unchanged. At fixed energy the
Gibbs state maximizes the fermion entropy current. Its temperature is
$T_{\rm eff}=\sqrt{T_H^2+12q_\chi/\pi}$ and its entropy excess is
$\pi(T_{\rm eff}-T_H)/6$. Since the partially transmitted state is not
that Gibbs state, the entropy inequality is strict. Eliminate
$T_{\rm eff}$ and use $T_H=1/(2\pi)$.
\end{proof}

For a hot-input excess $.02$ and $\chi=.5$, the output has
$q_\chi=.01$, $s_\chi=.04818703$ and thermal-floor excess
$q_\chi-T_Hs_\chi-3s_\chi^2/\pi=.000113460$.
A KMS output with the same heat has $s=.04863795$.
Their mean stress, geometry and clock can be identical despite their
different states and entropy currents. A geometric response alone does
not establish thermalization or identify the source's complete state.

A program that changes the contact also changes its stress. For a
smooth prescribed angle, point splitting the full unequal-time
covariance gives
\begin{equation}\label{eq:local-source-pump}
 q_{\rm eff}(\eta)=\sin^2\theta(\eta)q_{\rm hot}
                  +\frac{\theta'(\eta)^2}{4\pi},\qquad
 F_\eta=q_{\rm eff}+q_{\rm ref}.
\end{equation}
Each output port receives the same pump term. The total field work is
$\int\theta'^2d\eta/(2\pi)$ and is at least
$(\Delta\theta)^2/(2\pi\Delta\eta)$ over a switching interval.
Thus switching the selected interface has a source cost in the same
interaction, rather than triggering an unrelated emitter. The running
clock equations use $q_{\rm eff}(\eta)$, retaining their reference
stress. They also permit a stronger statement for changing programs.

\begin{proposition}[Positive growth under bounded changing sources]
\label{prop:changing-source-clock}
Suppose $q_{\rm eff}$ is smooth,
$0<q_{\min}\le q_{\rm eff}(\eta)\le q_{\max}$ and
$R_0^2>2a\kappa q_{\max}$, with
$\kappa q_{\min}/(2R_0)\le\delta_0\le2\kappa q_{\max}/R_0$.
Then the source--clock ODE with $q_{\rm eff}$ has a global positive
solution satisfying
\begin{align}\label{eq:changing-source-growth}
 \frac{\kappa q_{\min}}{2R}&\le\delta\le
                 \frac{2\kappa q_{\max}}R,&
 R'&\ge\frac{\kappa q_{\min}}2,\nonumber\\
 \int_0^\infty\delta^2d\eta&\le
       \frac{8\kappa q_{\max}^2}{q_{\min}R_0},&
 R-R_0&=\kappa\int q_{\rm eff}d\eta
             -a\Delta\delta-\frac a2\int\delta^2d\eta.
\end{align}
The reference correction stays bounded. If the actual source has a
long-time mean, the mean horizon entropy rate is $2\pi$ times that
mean. This concerns the source--clock sector, not an unlimited body
energy reserve.
\end{proposition}
\begin{proof}
At the upper barrier, the relative derivative is at most
$-\kappa q_{\max}/a+2(\kappa q_{\max})^2/R^2<0$.
At the lower barrier it is at least
$\kappa q_{\min}/(2a)+(\kappa q_{\min})^2/(8R^2)>0$.
Thus the strip is invariant. Integrate its upper bound using
$R\ge R_0+\kappa q_{\min}\eta/2$ to obtain the integral estimate.
Integration of $R'+a\delta'+a\delta^2/2=\kappa q_{\rm eff}$ gives
the last identity. The strip also prevents finite-time blowup.
\end{proof}

A bounded clock-reference correction does not remove the program's
work. If $N$ disjoint contact changes each have angle increment at least
$\Delta_{\min}$ and duration at most $\nu^{-1}$, their total pump
work is at least $N\nu\Delta_{\min}^2/(2\pi)$.
At $N/\Delta\eta=\alpha\nu$ the mean work current is at least
$\alpha\nu^2\Delta_{\min}^2/(2\pi)$, half in each output.
This is an additional switching-charge requirement of this contact
architecture, not a universal heat bound to be added again when that
same output energy has already been counted in the entropy account.

The source is prepared by the contact, but the input reservoir is still
prepared in a specified clock frame. The contact's local scattering law
does not supply a microscopic moving material body or its controller.
Its other port, reservoir preparation, switching work, clock conversion
and recoil must be included before it can replace the complete source
in the mass-10 table. Nor may a stationary entropy formula be applied
instantaneously during the switch. For a stationary frequency cutoff
$\Lambda$ with the common vacuum retained, a separate current bound is
\[
 |\Delta q|\le\frac1\pi\sum_{T=T_H,T_o}
                  T(\Lambda+T)e^{-\Lambda/T}.
\]
This controls omitted thermal source energy, not the field's global
trace distance or the stress of turning the contact on.


\section{Finite preparation and quantum material source certificates}
\label{app:material-certificates}

\subsection{Finite preparation on a predicted running clock}
\label{sec:finite-thermal-preparation}

The finite interval need not be supplied by an infinite-energy KMS
state. Q-RAIF~C constructs a compact diffeomorphism of the vacuum
that agrees with $\exp[2\pi T\eta(v_-)]$ on the regular operating
interval~\cite[Section on thermal correlations from a finite-energy
preparation]{QRC}. Its full local covariance equals the thermal
$\eta$ covariance, while the global state is pure and has finite
energy and its second moment. Two independent copies give the local
contact's required inputs, retaining each copy's outside purification.
Their stress is exactly Eq.~\eqref{eq:actual-thermal-source} after
the same Poincar\'e subtraction. A local bounded regularization and
positive history clock approximate this preparation and its smooth
stress to any fixed finite tolerance. The field-energy gain is paid
by the clock-interaction energy; this is not a free reservoir renewal.

The prepared states and inertial contact now have a common fermion
carrier~\cite[Proposition on a positive contact on the full CAR
carrier]{QRC}. A compact-gauge compensation proves the single-contact
Hamiltonian positive and self-adjoint even when its static dressing
is not implemented. The flavor rotation retains the original two
fields' correlations instead of replacing them by independent current
baths. On this carrier the actual reaction and the exterior
smooth-source certificate apply together. Matching that quantum
contact to the moving material action remains a separate requirement.

A record selected by earlier local operations can also choose the
next thermal preparation on this same field~\cite[Proposition on
record-selected local thermal preparation]{QRC}. If its past field
protocol lies in an interval of length $a$, the next preparation
support has length $b$, and their gap is $d$, then the two complex
chiral fields give
$I_0=\frac13\log[1+ab/(d(d+a+b))]$.
The probability-weighted conditional thermal error is at most
$\sqrt{I_0/2}$. A smooth source test with actual and comparison
second-moment norms $M_a,M_i$ has error at most
$\sqrt{I_0}(M_a+M_i)$, in addition to the preparation-gate and clock
errors. The outside purification and prior records remain in the
device. The same relative-entropy bound includes a controller and
its initially retained references obtained by the local past protocol;
it passes through later joint channels on that controller and window.
Additional incoming fields must be included rather than treated as
independent. Separation, physical record transport and a new preparation
are charged; the comparison state is not a physical reset of a
correlated field. This permits a conditional finite update, while
causal transport in the moving material and renewal of its clock,
blanks and reference remain to be realized together.

The field contact can also acquire a finite-mass quantum position
\cite[Section on the contact position]{QRC}. The same full-CAR
Hamiltonian then conserves body-plus-field momentum; a finite anchor
contributes its own recoil. Its moving-frame term is
$(p-H_f)^2/(2M)$, so the fixed contact is a fixed-momentum large-mass limit, not
a free prescribed trajectory. Position tails require a generator-weighted
locality error before using the record-controlled thermal certificate.
This support is nonrelativistic and does not yet provide the covariant
material stress or the physical clock required below.
Prescribing a constant velocity is a different limit: the harmonic
support has a stiffness-dependent stability boundary, while the full
finite-mass recoil Hamiltonian remains positive.

Energy conservation in the unfolded carrier does not identify its
translation generator with bulk momentum. Folding its two flavors into
right- and left-moving fields makes the finite-width mixing a mirror
coupling. A local bidirectional current instead needs two flavors per
direction, with $P=H_R-H_L$ and recoil $(p-P)^2/(2M)$
\cite[Proposition on bidirectional current and physical recoil]{QRC}.
Its harmonic prescribed-velocity threshold is
$\sqrt{a_0/(a_0+d_R+d_L)}$, while the full finite-mass energy remains
positive. The added preparations and purifications enter the source
account, and each exhaust capacity counts its actually routed fields.
For the same interval lengths $a,b$ and gap $g>0$, their product free
vacuum gives $I_0=\frac23\log[1+ab/(g(g+a+b))]$, twice the
two-field renewal value above; the conditional errors use that value.
This supplies a physical momentum and local field density; deriving
the propagating material's renormalized quantum stress is still needed
before inserting that density as the complete matter in a JT apparatus.

For a flat-coordinate interval of length $\ell$, the compact
vacuum-orbit preparations at temperature $T$ have infimum energy
\[
 E_{\min}=\frac{\pi cT^2\ell}{12}
       -\frac{cT}{6}\tanh\frac{\pi T\ell}{2}.
\]
On a changing clock the actual map, including its compact completion,
sets the physical preparation cost. Its waveform cannot be redefined
after a read changes the future trajectory.

\begin{proposition}[Source error from a predicted thermal clock]
\label{prop:predicted-thermal-clock}
Prepare with $\widehat\eta$ while the actual clock is $\eta$.
On a fixed regular source interval suppose both first derivatives
are at least $m>0$, both $j$th derivatives have absolute value at
most $M_j$, and
$\|\eta^{(j)}-\widehat\eta^{(j)}\|_\infty\le\epsilon_j$ for
$j=1,2,3$. Put $J_T=\pi cT^2/12$. The prepared one-point stress
differs from the exact $\eta$-thermal stress by at most
\begin{equation}\label{eq:predicted-clock-source-bound}
 B_{\rm clock}=2J_TM_1\epsilon_1+\frac{c}{24\pi}
 \left(\frac{\epsilon_3}{m}+\frac{M_3\epsilon_1}{m^2}
       +\frac{3M_2\epsilon_2}{m^2}
       +\frac{3M_2^2\epsilon_1}{m^3}\right).
\end{equation}
A constant two-port contact combines its two input bounds with
weights $1-\chi,\chi$. A smooth source test $f$ has additional
clock error at most $\|f\|_1 B_{\rm clock}$, separate from the
autonomous execution's state and stress error.
\end{proposition}
\begin{proof}
Apply $|a^2-b^2|\le2M_1|a-b|$ to the squared Jacobian.
The ratio of third to first derivative changes by at most
$\epsilon_3/m+M_3\epsilon_1/m^2$.
Factoring the difference of the squared second-to-first ratios gives
$2M_2\epsilon_2/m^2+2M_2^2\epsilon_1/m^3$.
The Schwarzian's factor $3/2$ yields
Eq.~\eqref{eq:predicted-clock-source-bound}.
The same coordinate-vacuum subtraction cancels between the two
preparations, not between different geometric backgrounds.
\end{proof}

This tolerance makes a record-selected prediction testable without
an exactly known future geometry. It controls mean source, not complete
state thermality or entropy. Learning a predictor that maintains the
stated derivative tube, and embedding the local gate, clock and
two-port stress in the moving material action, remain physical tasks.
The thermal restriction has an exterior purification; it cannot be
inserted unchanged into a finite-window entropy-reset argument.

\subsection{A matched quantum material source}
\label{sec:material-vacuum-source}

Positivity of an excitation Hamiltonian leaves its vacuum contribution
to the source undetermined. A local Gaussian material makes the required
comparison explicit. On a finite interval with Dirichlet boundaries,
let $a\in C_c^\infty$ vanish near the endpoints and $m>0$. The spatial
form for an active field $\phi$ and a normalized angle $\Theta$ is
\[
 q_a=\int\bigl[|\phi'+a\Theta|^2+|\Theta'|^2+m^2|\Theta|^2\bigr]dx.
\]
It is closed and coercive. Its operator, with the same boundary
conditions for $a=0$, is
\[
 K_a=\begin{pmatrix}-\partial_x^2&-\partial_x a\\
 a\partial_x&-\partial_x^2+a^2+m^2\end{pmatrix}.
\]
Here $-\partial_xa$ denotes $-\partial_x\circ a$, including the
derivative of the field. A constant positive angle--storage mass
matrix and free spectator fields leave the following coefficient unchanged.

\begin{proposition}[Matched material vacuum coefficient]
\label{prop:material-vacuum-coefficient}
For the proper-time cutoff comparison
\[
 \Delta E_\Lambda=-\frac1{4\sqrt\pi}
 \int_{\Lambda^{-2}}^\infty t^{-3/2}
 \operatorname{Tr}(e^{-tK_a}-e^{-tK_0})\,dt,
\]
and a fixed reference scale $\mu>0$,
\begin{equation}\label{eq:material-vacuum-log}
 \Delta E_\Lambda=\frac{\log(\Lambda/\mu)}{8\pi}
                         \int a^2dx+C_\mu[a]+o(1).
\end{equation}
The finite part requires a common material normalization condition.
\end{proposition}
\begin{proof}
Write $K_a=-\partial_x^2+B\partial_x+C$, where
$B=\left(\begin{smallmatrix}0&-a\\a&0\end{smallmatrix}\right)$ and
$C=\left(\begin{smallmatrix}0&-a'\\0&a^2+m^2\end{smallmatrix}\right)$.
The antisymmetric connection $A=-B/2$ gives
$K_a=-(\partial_x+A)^2+P$, with
\[
 P=C-B'/2+B^2/4=
 \begin{pmatrix}-a^2/4&-a'/2\\-a'/2&m^2+3a^2/4\end{pmatrix},
 \qquad \operatorname{tr}(P-P_0)=a^2/2.
\]
The heat coefficient for $-(\nabla^2+E)$ uses $E=-P$
\cite[Eq.~(4.27)]{Vassilevich2003}. Matched boundary jets cancel,
so the trace difference is
$-\sqrt t\int a^2/(4\sqrt\pi)+O(t^{3/2})$.
Its proper-time integral gives Eq.~\eqref{eq:material-vacuum-log}.
Coercivity controls the large-$t$ end of that integral.
\end{proof}

The coefficient can be assigned to a covariant material action. For
the positive target metric
\[
 ds_{\rm mat}^2=(d\phi+\lambda\theta f(Z)dZ)^2
                  +m_Z^2dZ^2+I d\theta^2,
\]
linearization about a static texture $Z_0$, with $\theta=0$, gives
$a=\lambda f(Z_0)Z_0'/\sqrt I$. Direct curvature calculation yields
$\operatorname{Ric}_{ZZ}=-\lambda^2f(Z)^2/(2I)$ there.
The usual one-loop kinetic counterterm in the Euclidean action
\cite[Section~7.2]{Friedan1985} is, in the present normalization,
\[
 S_{{\rm ct},E}=\frac{\log(\Lambda/\mu)}{4\pi}
 \int\sqrt g\,\operatorname{Ric}_{AB}
                         \partial_\mu X^A\partial^\mu X^B.
\]
On the texture it cancels the logarithm in
Eq.~\eqref{eq:material-vacuum-log}. Other potential, boundary and
metric counterterms must be matched as well. The finite Gaussian model
is not the full charged-fermion contact or an all-order continuum construction.

The charged carrier can nevertheless be retained at finite material
resolution. Q-RAIF~C's positive material construction~\cite{QRC}
uses the full fixed-spin charge/current basis, including the unused
oscillator modes, and keeps the material momentum squares. Its
boson--fermion correspondence is an operator equivalence, not a
neutral-sector replacement. This supplies a positive nonlinear
regulated Hamiltonian. The local continuum, its thermal states and
its renormalized source remain to be matched; the correspondence
does not determine a new vacuum subtraction for each contact.
Smoothing a permitted bounded gate under this interacting Hamiltonian
also supplies the energy regularity needed for autonomous
record-controlled operation. Its positive regulated energy tests
have a source certificate, while their covariant continuum stress
still requires a separate limit.
A local positive binding potential also produces finite normalizable
angle and storage modes with the finite contact's mechanical
coefficients~\cite{QRC}. Their omitted material modes give an explicit
self-energy and a filtered causal-response error. This avoids an
infinite uniform angle, but an exact local light front and quantum
thermal state are not supplied by the collective-mode projection.
Local record-selected drives can instead be compared on the same
complete Gaussian Hamiltonian and initial state~\cite{QRC},
Proposition on local actuation with the initial state retained.
The actual field's characteristic function controls the joint state
error with references retained, and a centred Gaussian source changes
by the classical full retarded response. Its force work is charged;
an old mean contributes a cross term. This supplies a quantum
actuation/source bridge in the fixed-texture sector, not a new vacuum
for the projected body or a proof of the moving clock's stress and
thermal renewal.

A source can be derived from this driven material rather than specified
as a separate emitter. In the corresponding flat Gaussian sector,
$\xi_\Theta$ is the complete record-driven angle response, including
its coupling to the carrier. Outside the compact contact, the actual
right and left profiles satisfy
\begin{align}
 \xi_R(u)&=-\frac12\int a(y)\xi_\Theta(y,u+y)dy,
 &\xi_L(v)&=\frac12\int a(y)\xi_\Theta(y,v-y)dy,\nonumber\\
 \Delta T_{uu}&=(\partial_u\xi_R)^2,
 &\Delta T_{vv}&=(\partial_v\xi_L)^2.
 \label{eq:material-emitted-source}
\end{align}
Here $u=t-x$, $v=t+x$, the past displacement is zero and the
initial Gaussian mean is zero. The carrier equation is
$(\partial_t^2-\partial_x^2)\xi_\phi=\partial_x(a\xi_\Theta)$.
Its retarded Green function, integrated once in $y$, proves the two
profiles. Unchanged covariance and common vacuum subtraction prove the
stress identities. The integrated two-side carrier radiation is at
most the actual force work; the rest is retained material energy.
Different finite drive records can therefore change the source through
actual material motion. This is not injective: opposite drives have
equal quadratic stress on centred input, and $a=0$ emits no carrier
source while material energy can remain. Putting these flat-background
data into JT requires a common quantum-state and coordinate map plus
the driver, moving support and clock's stress; the formulas alone do
not supply that map.

Finite autonomous execution also needs a source certificate for this
material, rather than the free-carrier estimate. On a fixed Dirichlet
interval, Q-RAIF~C derives a relative form bound for each compact-time
quadratic energy or exterior null-stress test $O$:
\[
 \|(H+\Lambda)^{-1/2}O(H+\Lambda)^{-1/2}\|\le C_O,
 \qquad H=d\Gamma(\sqrt K),\quad\Lambda>0.
\]
The complete Gaussian mode tower and its common vacuum subtraction
are retained. If actual and target output purifications have
$e=\|\psi-\varphi\|$, $g_H=\|H(\psi-\varphi)\|$, then
\begin{align}
 |\langle O\rangle_\psi-\langle O\rangle_\varphi|
 &\le C_O\sqrt{\Lambda e^2+eg_H}\nonumber\\
 &\quad\times\left[\sqrt{\Lambda+\|H\psi\|}
                    +\sqrt{\Lambda+\|H\varphi\|}\right].
 \label{eq:Gaussian-material-source-error}
\end{align}
Energy-regular return under this same $H$ therefore certifies a common
output source reading, with old references retained. Uniform physical
rms bounds $M_a,M_i$ give $g_H\le M_a+M_i$, so the certificate
converges with its state error despite the growing history clock.
The bound does
not use ordinary trace distance on an unbounded stress. It is an
output comparison under the agreed subsequent evolution, not an
account of the driver and clock stress during their interaction.
Its constants depend on the finite interval. Matched source tests
can then enter the response estimate of this paper only after the
state, coordinate and total-source identification is supplied;
neither a uniform real-line limit nor the full moving JT apparatus
follows from this Gaussian certificate.
The unbounded-time weighted moments in the JT source hypothesis
also need their own tail control; finite smooth tests alone do not
certify them.

A physical change of the coupling or texture changes this fluctuation
operator. It must use the same measured material parameters and vacuum
normalization as the original configuration. Independently setting each
contact's vacuum energy to zero would also discard its response to those
changes. Conversely, the bare logarithm supplies no finite universal
gravitational price. A covariantly renormalized material functional
$\Gamma$, with $T_{\mu\nu}=-2\delta\Gamma/(\sqrt{-g}\,\delta g^{\mu\nu})$
and $E_i=\delta\Gamma/(\sqrt{-g}\,\delta X^i)$ for its scalar arguments, obeys
\[
 \nabla_\mu T^\mu{}_{\nu}=E_A\partial_\nu X^A+E_\Phi\partial_\nu\Phi,
\]
where $E_A,E_\Phi$ are its variational equations. On the material equations
this retains the compensating dilaton force. The same functional must
be varied for the material and gravity; a static Euclidean vacuum
comparison cannot replace the causal, initial-state response of a controller.


\section{Supported apparatus and an acquired-record source return}
\label{app:material-return}

\subsection{A material support with its own conserved core}
\label{sec:material-support}

A boundary action gives its conjugate energy but not the matter that holds
the boundary. A metric-only local apparatus is also blind to the dilaton
on a fixed AdS$_2$ patch. We now change the matter model explicitly,
without applying the preceding boundary readout laws to this new apparatus.
In natural units add a timelike worldline $\Gamma$ to the JT action:
\begin{align}\label{eq:material-action}
 I=\frac1{2\kappa}\int\sqrt{-g}\,\Phi(R+2/L^2)\,d^2x
       +I_{\rm outer}+I_{\rm fields}-\int_\Gamma M(\Phi)\,d\tau,&\nonumber\\
 M=m_{\rm rest}+U(\Phi)+\langle H_{\rm core}\rangle.&
\end{align}
The usual outer boundary terms are included. Take $m_{\rm rest}>0$ and
$U\ge0$. The core has worldline quantum action
$\int d\tau\,[i\psi^\dagger D_\tau\psi-
\psi^\dagger(m_{\rm rest}+U+H_{\rm core})\psi]$; its expectation
provides a semiclassical worldline source; the apparatus's internal spatial
geometry is not resolved in this approximation. Q-RAIF~C constructs a finite
time-independent, nonnegative $H_{\rm core}$ containing the battery,
clock-dependent interactions, detector, retained body and every thermal
output cell~\cite{QRC}. Its mean energy is conserved. Output cells remain
inside the apparatus until a release operation is supplied.

Let $\delta_\Gamma$ be the covariant line delta and
$M_\Phi=\partial M/\partial\Phi$. Variation gives
\begin{align}
 R+2/L^2&=2\kappa M_\Phi\delta_\Gamma,
       \label{eq:material-curvature}\\
 E_{\mu\nu}:=g_{\mu\nu}\Box\Phi-\nabla_\mu\nabla_\nu\Phi
                     -g_{\mu\nu}\Phi/L^2
       &=\kappa T^{\rm total}_{\mu\nu}.
       \label{eq:material-metric}
\end{align}
The material stress is not separately divergence-free:
\begin{equation}\label{eq:material-ward}
 \nabla_\mu T^{{\rm total}\,\mu\nu}
       =-M_\Phi\delta_\Gamma\nabla^\nu\Phi,
 \qquad
 \nabla_\mu E^{\mu\nu}
       =-\tfrac12(R+2/L^2)\nabla^\nu\Phi.
\end{equation}
The identities agree by Eq.~\eqref{eq:material-curvature}, including
energy exchange with the dilaton. Before radiation forces are added,
$M a^\mu=-(g^{\mu\nu}+u^\mu u^\nu)\nabla_\nu M$.
Direct dilaton coupling therefore changes the curvature equation at the wall.

\begin{proposition}[Positive static support]
\label{prop:material-support}
In a symmetric thin-wall limit choose a common normal toward increasing
dilaton. At a static wall let $\Phi=\phi>0$, with derivatives
$d_\pm=\partial_n\Phi|_\pm$ and extrinsic curvatures $K_\pm$.
The vacuum junction and constraint equations are
\begin{equation}\label{eq:material-junction}
 [d]=-\kappa M,\quad [K]=-\kappa M_\Phi,\quad
 K_\pm d_\pm=\phi/L^2,\quad M\overline K+M_\Phi\overline d=0.
\end{equation}
For $M>0$ and $D>\kappa M/2$, they have the solution
\[
 d_\pm=D\mp\kappa M/2,\qquad
 K_\pm=\frac{\phi}{L^2d_\pm},\qquad
 M_\Phi=-\frac{\phi M}{L^2d_+d_-}.
\]
A positive quadratic material potential realizes these data whenever
its stiffness $a$ obeys $M-M_\Phi^2/(2a)>\langle H_{\rm core}\rangle$.
If a positive inner invariant is required, also choose $\phi>Ld_-$.
For every finite conserved core mean energy, a choice of $M$, $D$ and
sufficiently large $a$ gives positive rest mass and linear radial stability.
\end{proposition}
\begin{proof}
In Gaussian normal coordinates $ds^2=dr^2-N(r)^2dt^2$,
$R=-2N''/N$ gives the curvature jump. The singular $tt$ metric equation
gives $[d]=-\kappa M$, and the regular $rr$ equation gives
$Kd=\phi/L^2$. Its jump gives the averaged force equation.
The displayed data satisfy all four conditions. Choose
\[
 \phi_* =\phi-M_\Phi/a,\quad
 U(\Phi)=\tfrac a2(\Phi-\phi_*)^2,\quad
 m_{\rm rest}=M-\frac{M_\Phi^2}{2a}-\langle H_{\rm core}\rangle>0.
\]
The regular patches near the wall are
$N_\pm=\cosh(r/L)+LK_\pm\sinh(r/L)$ and
$\Phi_\pm=\phi\cosh(r/L)+Ld_\pm\sinh(r/L)$.
They have positive lapse there and solve the vacuum equations.
\end{proof}

For $L=1$, $\kappa=.01$, $M=10$, $\phi=3$, $D=2$ and $a=100$,
\[
 d_+=1.95,\quad d_-=2.05,\quad
 \mathcal C_+=5.1975,\quad \mathcal C_-=4.7975,
 \qquad \mathcal C_\pm=\phi^2-L^2d_\pm^2.
\]
Here $M_\Phi=-7.50469043$, $\phi_*=3.07504690$, and
$m_{\rm rest}+\langle H_{\rm core}\rangle=9.71839811$.
The two-cycle core with 39 battery levels has mean $2.34$, including
its clock-dependent interaction, so $m_{\rm rest}=7.37839811>0$.
The 135-level core has mean $7.14$ and still fits with positive rest mass.
Support and a running finite internal circuit thus have jointly satisfiable
semiclassical source data, with clock and battery included in the mass.

The reference support is locally stable in radial vacuum dynamics.
On the positive-normal branch let
$\beta_\pm=\sqrt{\dot\Phi^2+(\Phi^2-\mathcal C_\pm)/L^2}$.
The moving-wall junction $\beta_--\beta_+=\kappa M(\Phi)$ gives
\begin{equation}\label{eq:material-potential}
 \dot\Phi^2+V(\Phi)=0,\qquad
 V=\frac{\Phi^2-\mathcal C_+}{L^2}
 -\left[\frac{\mathcal C_+-\mathcal C_-}{2\kappa L^2M(\Phi)}
                    -\frac{\kappa M(\Phi)}2\right]^2.
\end{equation}
At the reference, $V=V'=0$, $V''=388805871/5681780>0$, and the proper-time
linear frequency is $5.84937147$. Increasing $\mathcal C_+$ by $10^{-5}$
gives a small oscillation between approximately $2.99755$ and $3.00244$.
This proves linear radial stability at one vacuum support, not stability
through an arbitrarily long export history.

The stability construction extends beyond the numerical example.
At the supported point, differentiation of Eq.~\eqref{eq:material-potential}
with fixed $\mathcal C_\pm$ gives
\begin{equation}\label{eq:material-stiffness}
 V''(\phi)=\frac2{L^2}-\frac{2\phi^2}{L^4d_+^2}
 -\frac{4D\phi^2}{L^4d_+d_-^2}
 +\frac{2d_+d_-}{M}\,a .
\end{equation}
Choose $M>\langle H_{\rm core}\rangle$ and $D>\kappa M/2$ first.
The coefficient of $a$ is positive, so finite stiffness can satisfy both
$m_{\rm rest}>0$ and $V''>0$. This supports any fixed finite compiled core;
increasing accuracy or operating length may require a larger material device.

\subsection{An output pulse must recoil against its source}
\label{sec:material-recoil}

A thermal cell becomes radiation only through a release mechanism.
Even an idealized event must conserve momentum as well as energy.
Suppose the initially static wall emits an inward null pulse of energy
$q$ in its initial local rest frame. Take an isolated instantaneous event
with $0<q<M/2$, with no external holding impulse. Local four-momentum
conservation fixes
\begin{equation}\label{eq:wall-recoil}
 M_f=\sqrt{M^2-2Mq},\quad v_f=\frac q{M-q},\quad
 \dot\Phi_f=\frac{d_+q}{M_f},\quad
 \beta_{+,f}=\frac{d_+(M-q)}{M_f}.
\end{equation}

\begin{proposition}[The source--support charge balance]
\label{prop:wall-recoil}
Keeping $\mathcal C_+$ fixed and imposing the moving-wall junction after
the event gives
\begin{equation}\label{eq:wall-source-invariant}
 \mathcal C_{-,f}-\mathcal C_-=2\kappa L^2d_-q.
\end{equation}
The wall's Killing charge
$E_\Gamma=(\mathcal C_+-\mathcal C_-)/(2\kappa L^2)$ loses exactly
$d_-q$, the pulse's initial interior Killing energy.
For an internal mass drop $\Delta=M-M_f$,
$q=\Delta-\Delta^2/(2M)$ and recoil energy is $\Delta^2/(2M)$.
\end{proposition}
\begin{proof}
Initially $E_\Gamma=Md_++\kappa M^2/2$; after emission it is
$M_f\beta_{+,f}+\kappa M_f^2/2$.
Equation~\eqref{eq:wall-recoil} gives their difference
$(d_++\kappa M)q=d_-q$. Solving the mass shell for $q$ gives the last claim.
\end{proof}

At $q=.06$ in this support, $M_f=9.93981891$ and
$\Delta\mathcal C_-=.00246$. Reducing mass while pinning the wall
would introduce the spurious term $-\kappa^2L^2q^2$.
The inner horizon entropy has infinitesimal response
\[
 \frac{dS_{h,-}}{dq}\bigg|_{q=0}
       =\frac{2\pi L^2d_-}{\sqrt{\mathcal C_-}},\qquad
 S_{h,-}=\frac{2\pi}{\kappa}\sqrt{\mathcal C_-}.
\]
The number $q=.06$ is an emission example, not the finite core's net heat
automatically sent into the bulk. Its cells contain gross energy and an
initial thermal contribution. The reference and all correlated ports must
be transported before an excess source can be identified with charged heat.
The finite energy-sector gates do not establish a local field release
Hamiltonian or prepare a KMS exhaust on a running horizon clock.
This branch supplies a positive support and its charge balance while
leaving those transport operations explicit.

\subsection{Continuous radiation moves the supported apparatus}
\label{sec:radiating-support}

The mass, recoil and radiation can be evolved together. On each side of
the material wall use an ingoing Eddington--Finkelstein coordinate:
\[
 ds_\pm^2=-f_\pm dv_\pm^2+2dv_\pm d\Phi,\qquad
 f_\pm=\frac{\Phi^2-\mathcal C_\pm(v_\pm)}{L^2}.
\]
Keep the normal toward increasing $\Phi$ and write
$w=d\Phi/d\tau$, $\beta_\pm=\sqrt{w^2+f_\pm}$.
Let $P_{\rm in}$ be the proper power of inward radiation absorbed from
the outer side, and $P_{\rm out}$ the inward power sent into the inner
side. These are the actual renormalized source powers; an excess
over a thermal reference is specified separately below. Put
$M=m_{\rm rest}+U(\Phi)+e$ and $A=U'(\Phi)$, where $e$ is usable internal
mean energy. The mean worldline approximation is that of
Eq.~\eqref{eq:material-action}; junctions with dilaton-coupled surface
matter are standard gravitational constructions~\cite[Sec.~V]{Shen2024}.

\begin{proposition}[Radiating material junction]
\label{prop:radiating-junction}
For smooth supplied powers and initial data with $M>0$,
$\beta_--\beta_+=\kappa M$ and $\beta_\pm-w>0$, the coupled equations are
\begin{align}
 \dot e&=P_{\rm in}-P_{\rm out},&
 \dot M&=Aw+P_{\rm in}-P_{\rm out},\nonumber\\
 \dot w&=\frac{\beta_+P_{\rm out}-\beta_-P_{\rm in}
                         -A\beta_+\beta_-}{M}-\frac\Phi{L^2},
                     \label{eq:radiating-wall}\\
 \dot{\mathcal C}_+&=2\kappa L^2P_{\rm in}(\beta_+-w),&
 \dot{\mathcal C}_-&=2\kappa L^2P_{\rm out}(\beta_--w),\nonumber\\
 \dot v_\pm&=(\beta_\pm-w)^{-1}.&&\nonumber
\end{align}
Dots denote proper time. These equations preserve the mass junction and
obey $[K]=-\kappa A$. They have a local solution while the stated domain
and energy reserves hold.
\end{proposition}
\begin{proof}
The regular JT equation gives
$T_{v_\pm v_\pm}=\partial_{v_\pm}\mathcal C_\pm/(2\kappa L^2)$.
Its contraction with the wall velocity gives the two invariant equations.
The tangent and normal projections of the material force balance give
$\dot M-Aw=P_{\rm in}-P_{\rm out}$ and
$M\overline K+A\overline\beta=P_{\rm out}-P_{\rm in}$.
Direct evaluation of the extrinsic curvature gives
\[
 K_+=\frac{\dot w+\Phi/L^2+\kappa P_{\rm in}}{\beta_+},\qquad
 K_-=\frac{\dot w+\Phi/L^2+\kappa P_{\rm out}}{\beta_-}.
\]
Their jump and the mass junction give Eq.~\eqref{eq:radiating-wall}.
Differentiating $\beta_--\beta_+-\kappa M$ using these equations gives
zero. The right-hand sides are smooth inside the domain, giving local
existence and uniqueness. Positivity of quantum excess stress is not used
in this algebraic consistency statement.
\end{proof}

For $P_{\rm in}=0$, $R=\sqrt{\mathcal C_-}>0$ and
$d\eta/dv_-=R/L^2$, the actual clock conversion is
\begin{equation}\label{eq:material-clock-conversion}
 z:=\frac{d\tau}{d\eta}=\frac{L^2(\beta_--w)}R,\qquad
 \frac{dS_{h,-}}{d\eta}=2\pi z^2P_{\rm out}.
\end{equation}
Recoil enters both the source and the clock; keeping the wall fixed while
its mass decreases does not solve these equations.


\subsection{Stored mechanical energy, radiation and the moving support}
\label{sec:inertial-scalar-support}

The angle need not be specified externally. Q-RAIF~C gives a positive
inertial current Hamiltonian with a storage coordinate, its retarded
reaction and force spectrum~\cite[Section on an inertial contact]{QRC}.
A displacement selected by a record changes the actual excess current
stress, not merely a reconstructed parameter. That local rest-frame
model still needs its infrared representation, moving ports and complete
quantum stress to be matched to the material action. In particular,
$\gamma\langle\dot\theta^2\rangle$ alone is not its net heat: the
noise--velocity correlation contributes to the work.

A distinct classical branch makes a common moving action explicit.
Put a real massless scalar on each side of the wall, with actions
$-\frac12\int_{\mathcal M_\pm}\sqrt{-g}(\nabla\varphi_\pm)^2$.
Append the mechanical worldline action
\begin{equation}\label{eq:scalar-mechanical-action}
 S_\Gamma=\int d\tau\left[-m_{\rm rest}-U(\Phi)
       +\frac I2\dot\theta^2+\frac J2\dot y^2
       -\frac k2(\theta-y)^2-\frac K2 y^2\right],
 \qquad I,J,k,K>0,
\end{equation}
and impose $\varphi_\pm|_\Gamma=g_\pm\theta$ with worldline Lagrange
multipliers. The wall, storage and angle are all varied. These are full
nonchiral scalar fields, not two auxiliary right movers reinterpreted as
an anomaly-free bulk theory. With no classical incoming waves, their
outgoing solutions are $\varphi_\pm(\xi_\pm)=g_\pm\theta(\tau_{\rm ret})$.
Their proper powers and the mechanical equations are
\begin{align}\label{eq:scalar-mechanical-source}
 P_\pm&=g_\pm^2\dot\theta^2,&
 \gamma&=g_+^2+g_-^2,\\
 I\ddot\theta+k(\theta-y)&=-\gamma\dot\theta,&
 J\ddot y+k(y-\theta)+Ky&=0,\nonumber\\
 e&=\tfrac I2\dot\theta^2+\tfrac J2\dot y^2
              +\tfrac k2(\theta-y)^2+\tfrac K2y^2,&
 \dot e&=-P_+-P_- .\nonumber
\end{align}
Boundary variation gives the mechanical force; characteristics give
the powers. There is no independent supplied radiation function.
Incoming waves would add their force and work and are excluded in
this branch. Quantizing a moving boundary also requires its vacuum
stress and fluctuations; the classical formulas do not supply them.

The outer radiation now moves outward, unlike the outer input of
Eq.~\eqref{eq:radiating-wall}. Use outgoing EF coordinate $u_+$ outside
and ingoing $v_-$ inside. Set $M=m_{\rm rest}+U+e$, adding any fixed
probe offset, and keep $\beta_--\beta_+=\kappa M$. The same action gives
\begin{align}\label{eq:two-outgoing-wall}
 \dot w&=\frac{\beta_+P_- -\beta_-P_+ -U'\beta_+\beta_-}{M}
                       -\frac\Phi{L^2},\\
 \dot{\mathcal C}_+&=-2\kappa L^2P_+(\beta_++w),&
 \dot{\mathcal C}_-&= 2\kappa L^2P_-(\beta_--w),\nonumber\\
 \dot u_+&=(\beta_++w)^{-1},&
 \dot v_-&=(\beta_--w)^{-1}.\nonumber
\end{align}
The opposite invariant signs follow from
$T_{u_+u_+}=-\partial_{u_+}\mathcal C_+/(2\kappa L^2)$ and
$T_{v_-v_-}=\partial_{v_-}\mathcal C_-/(2\kappa L^2)$.
The tangent Ward equation is $\dot M=U'w-P_+-P_-$ and the normal
equation is $M\overline K+U'\overline\beta=P_--P_+$.
Here $K_\pm=(\dot w+\Phi/L^2+\kappa P_\pm)/\beta_\pm$.
These identities give both $[K]=-\kappa U'$ and conservation of the
mass junction. They do not follow by putting a negative outer input
into Eq.~\eqref{eq:radiating-wall}; the null direction also changes.

\begin{proposition}[A finite mechanical-source window]
\label{prop:mechanical-source-window}
Equations~\eqref{eq:scalar-mechanical-source} and
\eqref{eq:two-outgoing-wall} have a unique local solution in the domain
$M>0$, $\beta_\pm>0$, $\beta_++w>0$, $\beta_--w>0$, $\mathcal C_->0$.
They preserve the mass junction, $e\ge0$, and
\begin{equation}\label{eq:mechanical-source-budget}
 \int_0^\tau(P_++P_-)\,ds=e(0)-e(\tau)\le e(0).
\end{equation}
For every fixed finite proper-time window, positive static-support data
admit a sufficiently small positive coupling with nonzero radiation
throughout a nontrivial subinterval and a solution inside that domain.
\end{proposition}
\begin{proof}
The right-hand side is smooth in the stated domain. Direct differentiation
of the junction gives zero, and the positive quadratic mechanical energy
obeys the displayed budget. At $g_\pm=0$ the supported geometry is static
and the internal oscillator energy is conserved. Continuous dependence
on the couplings preserves its strict domain margins on any specified
compact time interval for sufficiently small positive $g_\pm^2$.
Nonstationary mechanical data give nonzero angle velocity.
\end{proof}

Balanced outgoing power also gives a necessary operating-lifetime
bound for this same apparatus. If its fixed operating domain satisfies
$0<\Phi\le\Phi_{\max}$, $\Phi^2>\mathcal C_-$ and
$\beta_--\dot\Phi\ge b_{\min}>0$, then
\begin{align*}
 \mathcal C_-(\tau)-\mathcal C_-(0)
   &\ge\kappa L^2b_{\min}Q_{\rm out}(\tau),\\
 Q_{\rm out}(\tau)&<
 \frac{\Phi_{\max}^2-\mathcal C_-(0)}{\kappa L^2b_{\min}}.
\end{align*}
Equation~\eqref{eq:two-outgoing-wall} integrates this inequality for
$P_-=P_+=P/2$. If each complete episode exports at least
$q_{\min}>0$, its count is below the latter bound divided by
$q_{\min}$, as well as limited by the installed program and prepared
bank. The two acquisitions in Foundation~III are one episode, not two
renewals. Enlarging the domain or changing its null-clock margin changes
this necessary bound; it is not a general impossibility theorem.

The common energy and geometry can be held fixed across records.
Use $I=10,J=20,k=1,K=.1$, $g_+^2=g_-^2=1/(4\pi)$ and $e(0)=.95$.
Install $\theta(0)=\dot\theta(0)=0$, $y(0)=A_R$ and
$\dot y(0)=\sqrt{2[.95-.55A_R^2]/J}$ for $A_R=.5,1.3$.
Both records then use the same mass-10 support, potential and initial
invariants of Section~\ref{sec:material-support}. With the positive
probe offset $.02$ below, their common $m_{\rm rest}$ is $8.74839811$.
At proper time 80 the two solutions give
\begin{equation}\label{eq:mechanical-source-numbers}
\begin{array}{c|cc}
 A_R&.5&1.3\\\hline
 \int(P_++P_-)d\tau&.428637558&.522051307\\
 e(80)&.521362442&.427948693\\
 \mathcal C_-(80)-\mathcal C_-(0)&.008793748&.010712134\\
 \Phi(80)&3.003737133&3.004545378
\end{array}
\end{equation}
Both remain outside the inner horizon with positive reserves. Integrating
both invariants independently preserves the junction to $8.3\times10^{-15}$;
another integrator agrees to $3.7\times10^{-12}$. The phase probe below,
at $\lambda=.02$, yields $Y$-outcome probabilities $.503428742$ and
$.504187486$. The difference is created during evolution, not by a
different initial mass or an externally selected source. The numerical
80-window check is not an interval-arithmetic domain proof.

The inner rapidity clock is determined by the same solution:
$\dot\eta=\sqrt{\mathcal C_-}/[L^2(\beta_--w)]$ and
$dS_{h,-}/d\eta=2\pi(d\tau/d\eta)^2P_-$.
This source is coherent classical radiation, not an $\eta$-KMS heat
output, so its horizon response cannot be substituted into the thermal
entropy bound. Installing the mechanical record has a preparation cost.
The finite return in Section~\ref{sec:record-actuated-source} instead lets
an acquired quantum record exert force on the same storage coordinate.
Preparing these initial mechanical data and renewing the whole apparatus
are distinct from that return.

\begin{proposition}[A rest-frame mean-stress bridge]
\label{prop:inertial-mean-bridge}
Use Q-RAIF~C's quadratic current--mechanical contact on a fixed inertial
line, with the same $I,J,k,K$ and a common centered Gaussian state for
mechanics and incoming currents. Displace only the mean mechanical data;
take the mean initial angle to be zero. Keep all centered correlations
and old references. For independent physical outgoing directions
$s=\pm$, let the contact coefficients obey
$\kappa_s=\sum_i q_{si}^2$ and $\gamma_s=\kappa_s/(4\pi)$, where
$q_{si}$ are the effective angle-coupling coefficients, not a change
of the carrier's physical charge sectors.
Assume the renormalized current products exist on the common domain.
In the zero-width limit of the retarded mean solution, away from
initial/contact singularities, the outgoing stress difference and mean
mechanical equations are
\begin{align}
 \Delta\langle T_{{\rm out},s}\rangle
    &=\gamma_s v(\tau_{\rm ret})^2, &
 I\ddot\theta+k(\theta-y)&=-(\gamma_++\gamma_-)\dot\theta,
 \quad v=\dot\theta. \label{eq:mean-contact-bridge}
\end{align}
Thus $g_s^2=\gamma_s$ identifies this excess mean stress and its
mechanical energy loss with Eq.~\eqref{eq:scalar-mechanical-source}
in that fixed local frame. The balanced scalar example requires
$\kappa_+=\kappa_-=1$. Q-RAIF~C's unscaled two-flavor bidirectional
contact has $\kappa_+=\kappa_-=2$ and twice that total damping;
scaling its angle coupling by $1/\sqrt2$ gives the balanced normalization.
\end{proposition}
\begin{proof}
Take $f_a(x)=a^{-1}f(x/a)$ with nonnegative $f\in C_c^\infty(\mathbb R)$
and $\int f=1$. At finite width the displacement adds the classical
linear solution without changing the common covariance. In direction $s$
its current mean shift is $q_{si}h_a/(2\pi)$, where $h_a$ is the retarded
convolution of $v$ with $f_a$ along that characteristic. The Sugawara
stress $\pi\sum_i:j_{si}^2:$ has zero cross term with the centered
baseline. Its difference is therefore $\gamma_s h_a^2$.
Eliminating the mean currents gives the retarded kernel
$\Gamma_{a,s}(t)=\kappa_s\int f_a(x)f_a(x-t)dx/(2\pi)$, whose
one-sided integral is $\gamma_s$. For smooth mean histories its
convolution error for a smoothly extended velocity is bounded by
$\|v'\|_\infty\int_0^\infty t\Gamma_{a,s}(t)dt=O(a)$.
For the initial-value solution, the initial kernel tail has integrated
size $O(a)$ as well. The common kernel mass bounds the linear Volterra
solution uniformly on fixed windows; Gronwall then gives convergence
to the displayed damped equation. Beyond the contact, $h_a$ tends to
the retarded velocity, proving the stress statement. The resulting
positive quadratic displacement energy obeys
$\dot e=-(\gamma_++\gamma_-)v^2$, the scalar-source balance.
Incoming noise--velocity work and centered quantum stress remain in
the common baseline, rather than being set to zero.
\end{proof}

This identifies an excess mean radiation channel. Absolute quantum
stress, a quantum zero-width channel limit, dilaton-dependent mass,
moving-frame work, physical support and recoil, acquired-record spring
forces and stress fluctuations require further matching. In particular
a statistical mixture of mechanical records has branch covariance
and cannot be replaced by one Gaussian displacement in this proposition.

\subsection{A quantum phase reads the moving dilaton}
\label{sec:material-dilaton-probe}

The material model also permits a direct finite sensor. For $\lambda>0$, add a probe with
\begin{equation}\label{eq:material-probe}
 H_P(\Phi)=\epsilon_P I+\lambda(\Phi-\phi_0)Z_P,
 \qquad \epsilon_P>\lambda\sup_\Gamma|\Phi-\phi_0|,
\end{equation}
initially in $|+\rangle$. This is the same kind of explicit dilaton
coupling as the support, now in an internal quantum degree of freedom.
It does not follow from metric-only matter.
During the reading, $\langle Z_P\rangle=0$ and $Z_P$ is conserved.
In the mean worldline approximation its energy is $\epsilon_P$ and its
mean additional dilaton force is zero. Count $\epsilon_P$ in the mass,
reducing $m_{\rm rest}$ by that amount in the displayed examples.
The probe's phase gives
\begin{equation}\label{eq:material-probe-record}
 \Pr(Y_P=+1)=\frac{1+\sin[2\lambda\int(\Phi-\phi_0)d\tau]}2.
\end{equation}
With $\lambda=\epsilon_P=.02$, the $q=.01$ and $.02$ histories give
$.50617417$ and $.51230214$ at the same proper duration.
Their actual bit laws have total variation distance $.00612797$.
No numerical dilaton value is supplied to the controller.

The terminal read itself can be implemented. Q-RAIF~C gives a positive
autonomous clock whose last two gates write the actual $Y$ result into
a record, then apply a true SWAP with a paid coherent blank~\cite{QRC}.
The probe is renewed; all old correlations move into the blank and
record. For a supplied finite $\Phi(\tau)$ trajectory, its joint error
relative to the actual phase at $\tau_0<T$ is bounded by
\begin{equation}\label{eq:material-terminal-read}
 d_{\rm joint}\le2A\tau_0\sqrt{P_0}+2\sqrt{P_0}
                    +\lambda D(T-\tau_0),\quad
 D=\sup|\Phi-\phi_0|,
\end{equation}
where $P_0$ is the clock's terminal binomial tail and $A$ is the sum
of its last two edge strengths. Finite resources suffice for each
specified error; the clock mean energy is $Jm$, and must enter the
supported mass. The sufficient $m=10^5$, $T=80$, $\tau_0=78$ example
uses clock mean energy $1963.4954$, not the mass-10 apparatus above.
During writing and renewal the mean additional force need not vanish:
the material equation uses $A_{\rm mat}=U'+\lambda\langle Z_P\rangle$,
with $|A_{\rm mat}-U'|\le\lambda$. The probe, clock, blank and record
belong in the same matter action. If $e$ now includes their
$\Phi$-dependent internal energy, its equation is
$\dot e=\lambda\dot\Phi\langle Z_P\rangle+P_{\rm in}-P_{\rm out}$;
then $\dot M=A_{\rm mat}\dot\Phi+P_{\rm in}-P_{\rm out}$ as required
by the material Ward balance.

A smaller clock can be included directly in this mean material action,
instead of using the large worst-case certificate. Take $m=128$,
$J=\pi/160$, $\epsilon_P=\lambda=.02$, and retain the complete
129-state clock, probe, record and blank throughout the proper-duration
80 evolution. Its initial internal mean energy is $2.5532741$.
Choose $m_{\rm rest}=4.8251240$ so that the initial mass 10 already
includes it and the declared source reserve 2.34. Evolve the joint
Schr\"odinger equation for
\[
 H_{\rm int}(\Phi)=H_{\rm hist}+2\epsilon_P I
                       +\lambda(\Phi-\phi_0)Z_P
\]
with the radiating junction and its actual $\eta$ clock. The two contact
outputs $\chi=.5$ and $1$, at hot-input excess $.02$, give
\begin{center}
\begin{tabular}{cccc}
\hline
$q_\chi$ & actual $\Pr(R=1)$ & final source reserve & probe reset distance\\
\hline
.010 & .50529766 & 1.42403868 & .00075428\\
.020 & .51055894 & .50256908 & .00150041\\
\hline
\end{tabular}
\end{center}
The last column is the distance to the closest pure product of a
renewed probe and a normalized retained internal reference, not an error
certificate against the ideal meter's entire reference state. The actual record
laws differ by $.00526128$. Norm, mass-junction and internal dilaton-work
residuals are below $4\times10^{-11}$ in this integration; a second
integrator agrees on the displayed trajectory and record values within
$1.3\times10^{-9}$. These are numerical checks, not interval-certified
error bounds. For these gates and initial data, the vanishing mean force
has an exact reason.

\begin{proposition}[A force-neutral read and replacement]
\label{prop:force-neutral-meter}
Let $K$ denote conjugation in the clock and computational bases and set
\[
 \mathsf C=(-1)^{N_{\rm clock}}X_PX_B K,\qquad
 H_0=H_{\rm int}-(Jm+2\epsilon_P)I.
\]
For the identity-edge clock with terminal $C_Y$ and true SWAP,
$\mathsf C H_0(\Phi)\mathsf C^{-1}=-H_0(\Phi)$ and
$\mathsf C Z_P\mathsf C^{-1}=-Z_P$.
With initial clock state $0$ and data $|+\rangle_P|0\rangle_R|+\rangle_B$,
every real dilaton trajectory has
$\langle Z_P\rangle=0$ and
$\langle H_{\rm int}\rangle=Jm+2\epsilon_P$ throughout the read.
\end{proposition}
\begin{proof}
Conjugation exchanges the two $Y$ projections, and $X_P$ exchanges
them again. Thus $X_PX_B C_Y^*X_PX_B=C_Y$; the real SWAP commutes
with $X_PX_B$. Clock parity reverses every off-diagonal edge, while
$X_P$ reverses $Z_P$. This proves both operator identities.
Remove the constant Hamiltonian phase. Since $\mathsf C$ is
antiunitary, it reverses both $i$ and $H_0$ in the Schr\"odinger
equation, and so maps its solution to another solution with the same
initial state. Uniqueness makes that solution $\mathsf C$-invariant.
Expectations of the two odd Hermitian operators must vanish.
\end{proof}

The clock energy is therefore included from the start without changing
the previously supported mean trajectory. This does not set the quantum
force operator or its fluctuations to zero, and different blanks or
gates need not have the symmetry. The general coupled force equation,
not a zero-force constraint, was used in the numerical calculation.
The thermal reservoir's quantum degrees, both-port material stress and
its local preparation in $\eta$ are not thereby included in that body.

The same finite meter also reads the mechanical-source branch of
Section~\ref{sec:inertial-scalar-support}. Its complete mean energy
$2.5532741$, rather than just the probe offset, is included from the
start by using the common positive rest mass $6.21512398$ for both
equal-energy records. Proposition~\ref{prop:force-neutral-meter} leaves
their mean mechanical and geometric trajectories unchanged.
Their actual acquired record probabilities are $.502991099$ and
$.503655046$, a total variation distance $.000663947$.
Keeping clock, blank and record gives distances $.000376618$ and
$.000458340$ to the closest renewed-probe pure product with its normalized
internal reference. These are not errors against an ideal reference
state. The full vector evolution preserves norm to $1.6\times10^{-11}$
and internal mean energy to $2.6\times10^{-12}$; a second integration
agrees on the record probability to $3.1\times10^{-14}$.
This finite episode obtains the record through physical gates, rather
than reading the numerical solution into a controller. In this force-neutral
example the record does not act on the mechanics. The next section adds
that interaction and includes its reaction during acquisition. The scalar
source has no thermal reservoir whose regeneration could be inferred
from either construction.

For input streams jointly independent of the record and prepared in a
common clock frame, the record can select
either of two contact angles in the next stationary window.
If their transmissions are $\chi_0\ne\chi_1$, its
source-current mean is
\[
 q_{\rm next}=q_{\rm hot}[\chi_0+(\chi_1-\chi_0)\Pr(R=1)].
\]
It depends on the actual phase bit, not a supplied numerical dilaton.
The angle change's pump work in Eq.~\eqref{eq:local-source-pump} is
also charged. This gives the conditional response-to-source map;
it does not yet autonomously prepare the reservoirs or renew the
entire controller. Quantum fluctuations
of the probe force also lie outside the mean-source approximation.
The phase sensor and the earlier fixed-boundary current sensor are
different realizations, not two freely composable laws for one apparatus.

\subsection{An acquired record drives a new material source}
\label{sec:record-actuated-source}

A record can act without first being converted into a supplied mechanical
initial condition. Keep the two classical scalar fields and the material
action of Section~\ref{sec:inertial-scalar-support}, and append
$i\langle\Psi,\dot\Psi\rangle-\langle\Psi,H_c\Psi\rangle$ to its
proper-time Lagrangian. The normalized state $\Psi$ contains a finite
history clock, probe $P$, record $R$, replacement blank $B$ and an old
purification $Q$. Set
\begin{equation}\label{eq:record-spring-core}
 H_c(\Phi,y)=H_{\rm hist}+E_{\rm off}\mathbf1
       +\lambda(\Phi-\phi_0)Z_P
       +\frac{\chi_R}{2}P_R(y-d)^2,
 \qquad P_R=|1\rangle\langle1|_R,\quad\chi_R\ge0.
\end{equation}
$H_{\rm hist}$ is the positive gate-conjugated clock of Q-RAIF~C.
Its first nonidentity edge performs $C_Y$, the next performs a true
SWAP of $P$ and $B$, and the remaining edges are identities.
There is no proper-time switch. Initially $R=0$, so the spring has
zero expectation. The actual acquisition opens its force channel.
$H_c\ge0$ while $E_{\rm off}>|\lambda(\Phi-\phi_0)|$.
The target $d$, couplings and initial clock are prepared hardware.
The scalar fields, mechanical coordinates and geometry use an Ehrenfest
mean-action approximation~\cite[Section~2]{CastroGross2013};
the finite quantum core and its old reference are retained
in full, without sampling or collapsing $R$.

Let $p=\langle P_R\rangle$, $z=\langle Z_P\rangle$ and
$E_c=\langle H_c\rangle$. Variation of this one action gives
\begin{align}\label{eq:record-spring-work}
 i\dot\Psi&=H_c\Psi,& F_y&=\chi_R(y-d)p,& F_\Phi&=\lambda z,\\
 J\ddot y+k(y-\theta)+Ky&=-F_y,&
 \dot e&=-P_+-P_--F_y\dot y,&
 \dot E_c&=F_y\dot y+F_\Phi w.\nonumber
\end{align}
The angle equation and $P_\pm=g_\pm^2\dot\theta^2$ are unchanged.
Here $p$ is an expectation obtained by variation, not a decoded number
fed to a separate control policy. The mass and its tangent balance are
\begin{equation}\label{eq:record-spring-mass}
 M=m_{\rm rest}+U+e+E_c,\qquad
 \dot M=(U'+F_\Phi)w-P_+-P_- .
\end{equation}
Consequently Eq.~\eqref{eq:two-outgoing-wall} holds with $U'$ replaced
by $U'+F_\Phi$ in its acceleration equation; both invariant and null-clock
equations retain their directions. Differentiating the mass junction
with these equations gives zero. The separate work balances are
\begin{align}\label{eq:record-spring-ledgers}
 e(\tau)+Q_{\rm out}(\tau)+\int_0^\tau F_y\dot y\,ds&=e(0),\\
 E_c(\tau)-E_c(0)&=\int_0^\tau(F_y\dot y+F_\Phi w)\,ds.\nonumber
\end{align}
The clock interaction and conditional spring are included in $E_c$.
The spring breaks the force-neutral symmetry used above; $F_\Phi$ cannot
be set to zero in this coupled calculation.

\begin{proposition}[A nonzero record-mediated source return]
\label{prop:record-mediated-source}
Take two initial geometries with the same $\Phi=\phi_0$, $w=0$, mass,
mechanical data and quantum input, but different $\mathcal C_-$;
set $\mathcal C_+$ by the mass junction. In the positive exterior domain,
either $\lambda=0$ or $\chi_R=0$ makes their proper emitted powers
identical. With both interactions present, there are positive finite-clock
data for which the emitted powers differ on a nonempty finite interval.
Their record, mechanical response, source and geometry obey the same action
and work account.
\end{proposition}
\begin{proof}
For $\chi_R=0$, the mechanical subsystem is independent of the reader
and geometry. For $\lambda=0$, the joint mechanical and quantum equations
are independent of geometry. Either condition fixes identical mechanical
sources for the two initial geometries.

For a nonzero example choose $\theta(0)=\dot\theta(0)=0$,
$y_0(y_0-d)\ne0$ and probe state $(I+\upsilon X_P)/2$, $\upsilon\ne0$,
with a blank $R$ and its old purification held. Let $r\ge1$ be the
$C_Y$ edge, $a_j>0$ the clock edge strengths and
$G_r=\prod_{j=1}^r a_j$. The initial acceleration difference
$\delta a$ follows from the junction equations, so that
$\delta\Phi(t)=\delta a\,t^2/2+O(t^3)$.
The shortest-path Dyson expansion gives
\begin{equation}\label{eq:record-spring-leading-record}
 \delta p(t)=K_r t^{2r+3}+O(t^{2r+4}),\qquad
 K_r=\frac{\upsilon G_r^2\lambda\delta a}
             {6r!(r-1)!(r+3)}.
\end{equation}
An insertion of $\lambda\delta a\,t^2Z_P/2$ after $j<r$ forward
edges contributes the time integral
$t^{r+3}(j+1)(j+2)/(r+3)!$.
Their sum is $t^{r+3}r/[3(r+3)r!]$.
Using $C_Y^*P_RC_Y=(I+Y_P)/2$ in the interference with the unperturbed
shortest path gives the coefficient above. The insertion after all
$r$ edges has Hermitian $P_RZ_P$ in that interference and zero real
part. Further graph paths and spring insertions raise the order;
the spring annihilates the blank record before acquisition.

Integrating the mechanical equations then yields
\begin{equation}\label{eq:record-spring-leading-source}
 \delta(P_++P_-)(t)=
 -\frac{2\gamma k^2 y_0\chi_R(y_0-d)K_r}
 {I^2J(2r+4)(2r+5)(2r+6)}t^{2r+7}+O(t^{2r+8}).
\end{equation}
The baseline angle velocity is $ky_0t/I+O(t^2)$.
Different initial $\mathcal C_-$ give a nonzero $\delta a$ for generic
supported data, since initially
$\dot w=-U'\beta_+\beta_-/M-\phi_0/L^2$.
The displayed coefficient is nonzero for the stated nonzero parameters.
The coupled finite-dimensional equations are analytic on their strict
positive exterior domain. Local existence and continuity therefore give
a nonempty admissible family, not just a numerically resolved difference.
\end{proof}

A complete integration uses the earlier $I,J,k,K$, $\gamma=1/(2\pi)$,
$e(0)=.95$, $y_0=.5$ and $d=-.5$, with
$\lambda=.2$, $\chi_R=.05$ and $E_{\rm off}=.04$.
The 33-state clock has edge scale $\pi/160$, $r=3$, mean energy
$.628318531$ and spectral maximum $1.256637061$.
The initial probe has $\upsilon=.8$, its two-dimensional old purification
is retained, and $B$ starts in $|+\rangle$.
The quadratic support has stiffness $100$ and
$U'(3)=-30/(1.95\cdot2.05)$, as in the static mass-10 support.
Both geometries have initial $\Phi=3$, $w=0$, $M=10$ and
$\mathcal C_-=4.7975\pm.1$, with the outer invariant fixed by the junction.
The initial core energy $.668318531$ is already in the mass;
the common rest mass is $8.100079577$. At proper time 100:
\begin{center}
\begin{tabular}{cccc}
\hline
$\mathcal C_-(0)-4.7975$ & $\Pr(R=1)$ & $Q_{\rm out}$
 & $\mathcal C_-(100)-\mathcal C_-(0)$\\
\hline
$-.1$ & $.506528009$ & $.518808111$ & $.010786169$\\
$+.1$ & $.496292018$ & $.517763845$ & $.010480933$\\
\hline
\end{tabular}
\end{center}
The largest proper-power difference on this window is
$1.0350613\times10^{-4}$; with either interaction removed it is below
$3.2\times10^{-16}$. With feedback removed the two record laws still
differ, while their scalar sources agree. Thus acquisition and actuation
are separately tested. The nonzero-source result concerns the emitted
power, not the already different initial geometries.

The full core--reference vector is evolved along with both invariants.
Norm, junction and the three work balances, including $U$'s dilaton work,
have residuals below $3.7\times10^{-14}$ in the primary integration.
A separate tensor-neighbour implementation with another integrator agrees
on the record and source to $6.0\times10^{-13}$ and
$2.0\times10^{-15}$, respectively, and on the mass to $1.2\times10^{-8}$.
The largest probe force is about $.0082$ and is retained.
Clock energy changes during acquisition; it is not an uncharged timer.
Only the $C_Y$ edge changes $P_R$, giving
\[
 |\dot p|\le2a_r\|\Psi_{r-1}\|\|\Psi_r\|.
\]
This bound is checked on the interacting clock, not a free-clock tail.
It supplies finite record retention when that edge is sparsely occupied,
not permanent halt. The construction closes an actual finite
geometry--record--force--source return in the mean material branch.
Full quantum recoil, covariant thermal transport and renewal of the clock,
blank and reference remain required for the sustainable loop below.


\subsection{A spatial return with a physical clock}
\label{app:finite-spatial-return}

The geometry and support below are semiclassical mean fields. The internal
oscillators and both retained spins are quantum. The regulator is part of
the physical action; the statement does not remove its bandwidth cutoff.

Take JT gravity with radius $R$ and coupling $\kappa>0$, a full material
label $a$, and a full synchronization scalar $\theta$, with
\begin{equation}
 S_b=-\int\sqrt{-g}\,M\left[mn+\frac K2(n-n_0)^2\right],\quad
 n^2=(\nabla a)^2,\qquad
 S_\theta=-\frac J2\int\sqrt{-g}(\nabla\theta)^2.                 \label{eq:finite-spatial-1}
\end{equation}
Require $m>Kn_0>0$, $n$ near $n_0$, and a time function
$\theta_t>0$ with timelike gradient. Use the unrestricted inverse embedding
$\theta(X(\tau,\xi))=\tau$, $a(X(\tau,\xi))=\xi$ and
$h=X^*g$. Only the internal expansions are finite:
$\sqrt I\chi=\sum q_{c\alpha}e_\alpha(\xi)$ and
$u=\sum q_{u\alpha}e_\alpha(\xi)$, for smooth supported real $e_\alpha$.
Substitute them into the two scalar kinetic actions and the potential
\begin{equation}
 V=\frac{I\mu^2\chi^2}{2}
   +\frac12\left[\omega u+\frac{A\chi}{\omega}\right]^2
   +E_o\rho_o-D_o,
 \quad A=\lambda+gwY_P+\gamma w_2Y_S,                         \label{eq:finite-spatial-2}
\end{equation}
before Legendre transform and metric/embedding variation. Here
\[
 D_o=\frac{C_\chi}{\omega^2}
 (\lambda gwY_P+\lambda\gamma w_2Y_S+g\gamma ww_2Y_PY_S),
 \qquad C_\chi(\xi)=e^TQ_ce/I .
\]
$Q_c$ is the fixed prepared clock covariance. Choose the paid scalar
offset to dominate $|D_o|$. Quantize the resulting Hamiltonian with Weyl
ordering and a joint-state Berry term on the actual $\theta$ leaves. Its
kinetic matrices are
\[
 C=-\int\sqrt{-h}\,h^{\tau\tau}ee^T,\quad
 D=-\int\sqrt{-h}\,h^{\tau\xi}e(e')^T,\quad
 K_{\rm grad}=\int\sqrt{-h}\,h^{\xi\xi}e'(e')^T,
\]
giving $(p-Dq)^TC^{-1}(p-Dq)/2+q^TKq/2$ plus the paid
offset/calibration terms. A proper-number sensor contributes
\begin{equation}
 H_S=\int n_S(\xi)N[E_S+h_s\Phi(X)Z_S]\,\dd\xi,
 \quad N=\sqrt{-h_{\tau\tau}},\quad \int n_S\dd\xi=1.        \label{eq:finite-spatial-3}
\end{equation}
$n_S$ is smooth, symmetric, and has nonzero width. Include both spins,
their Berry terms, all profile reactions, and any old reference in the
same retained state. Restrict $|h_s|\sup|\Phi|<E_S/2$. No profile is an
external time-dependent signal.

\begin{proposition}[Finite-bandwidth spatial return]
\label{prop:finite-spatial-return}
There are finite positive $I,M,J,R$, smooth finite-width preparations,
and nonzero finite $h_s,\gamma$ for which this one mean-action system has
a positive local solution and two actual histories with the same initial
geometric constraints, differing only in their retained $P$ record,
whose next-program probabilities for $(1+Y_S)/2$ differ at the same
central proper time. The retained $Y_S$ enters the subsequent actual
material Hamiltonian, with a nonzero mixed sensing/return response of
the material mean acceleration. All clock, support, interaction
and probe sources are metric variations of this action, and satisfy the
full Ward identity. This is an existence result, not a uniform gain,
state-independent CPTP map, continuum-controller or thermal-renewal
statement.
\end{proposition}

\begin{proof}
Prepare $P=|\pm Y_P\rangle$, $S=|+X_S\rangle$, a common free product
Gaussian field state with $\langle q\rangle=0$, zero symmetrized $qp$
covariance, and clock momentum $\sqrt I\Omega F$. References are
retained, not reset. At $\gamma=h_s=0$, calibration cancels the initial
expected record-dependent potential as a functional of $X$. Thus the
complete initial metric stresses coincide, not just their integrals.
Take common static AdS data, $\theta=t$ on a finite prepared plateau,
and hydrostatic bar density
$n=n_0+(m/K)[\cos(x/R)-1]$ on a small causal patch. The common initial
JT constraint is a linear spatial equation for $\Phi$; adding a
sufficiently large finite positive homogeneous solution
$B\sec(x/R)$ makes $\Phi>0$. This changes initial geometric data, not
the charged positive source.

We first check the physical time equation. At rest, varying
$k=X_\tau^t$ gives $C=\int x'ee^T/k$,
$K_{\rm grad}=\int ke'(e')^T/x'$ and
$K_{\rm mass}=\int e^{2\rho}kx'ee^T\mathcal M_\pm$.
Hence the local time-momentum current $\delta H/\delta X_\tau^t$
contains only $q^2,p^2$ and commuting spin constants. Its commutator
with the initial Hamiltonian is Weyl $qp$; its expectation, and the
mixed time/spatial current expectation, vanish. The corresponding
initial mixed principal block vanishes too. The full time EL equation
therefore gives $\theta_{tt}(0)=0$, or $\theta=t+O(t^3)$.
It changes no potential trace coefficient of order $t^2$.
Nevertheless, if $\theta=t+\theta_3t^3/6+\cdots$, the actual
$T_{\theta,tt}=T_{\theta,xx}=J/2+J\theta_3t^2/2+\cdots$
is retained; only its trace is identically zero.

Here is a nonempty analytic trace certificate. Choose one normalized
even smooth mode $e$, constant and positive in an open central
neighborhood, $w=1$ on its support, and $F=1$. Put
$k=\int(e')^2$, $\beta=\int\sec^2(x/R)e^2$,
$Q_c=[2\sqrt{k+\mu^2\beta}]^{-1}$,
$Q_u=[2\sqrt{k+\omega^2\beta}]^{-1}$.
The exact static central trace-contrast coefficient is
\begin{align}
 d_2&=-D_0+B_1/I+B_2/I^2,\quad
 D_0=4\lambda g\Omega^2e(0)^2/\omega^2,\notag\\
 B_1/D_0&=\beta[\mu^2Q_c+\omega^2(Q_c+Q_u)]/\Omega^2,
 \\
 B_2/D_0&=2\beta(\lambda^2+g^2)Q_c/(\omega^2\Omega^2).       \label{eq:finite-spatial-4}
\end{align}
These follow by $[t^2]Q=-\{K_\pm-K_0,Q_0\}/2$; the off-diagonal
species covariance is included. The initial spatial momentum current
is locally proportional to $ee'\operatorname{sym}(qp)$, so its force
contrast and spatial derivatives vanish on the central plateau. No
bar recoil trace term of order $1/I$ survives there. For
$(\mu,\omega,\lambda,g,\Omega)=(.3,1.4,.6,.2,.8)$,
$B_1/D_0\le6.432292\beta$ and $B_2/D_0\le1.062926\beta$.
Take finite $R$ with $\beta<1.01$, and $I=100$.
Then $d_2<-0.9349D_0$.
Scale $M=\alpha L,J=\beta_bL$ together, with fixed positive ratios.
Actual common-acceleration corrections to \eqref{eq:finite-spatial-4} are $O_I(L^{-1})$;
hydrostatic preparation cancels the order-$M$ force, not the quantum
reaction. A sufficiently large finite $L$ preserves $d_{2,\rm total}<-D_0/2$.

For completeness the actual coupled equations have a local solution.
Let $\mathcal A=1+\sum(q_j^2+p_j^2)$. A degree-two Hamiltonian gives
$\|\mathcal A^sU_z\psi\|\le e^{C_sT}\|\mathcal A^s\psi\|$ and
\begin{equation}
 \|\mathcal A^2(H_z-H_{\hat z})\psi\|
 \le C\|z-\hat z\|_{H^{r-1}}\|\mathcal A^3\psi\|,
 \qquad r\ge4.                                             \label{eq:finite-spatial-5}
\end{equation}
Here $z$ consists of first-order classical variables. Quadratic
principal/source moments and quartic differentiated-current forces are
Lipschitz in $D(\mathcal A^2)$ on this graph ball. The latter degree,
which arises from noncommuting spin coefficients, is not suppressed.
The full bar and synchronization principal forms have size $L$ and
strict positive kinetic/spatial margins; their finite quantum
corrections, including smooth finite-rank kernels, have bounded norm.
Large finite $L$ preserves those margins.

To express geometry on actual leaves, use ambient-frame derivatives
$P=(\rho_t,\rho_x)$ and $Q=(\Phi_t,\Phi_x)$ pulled back to $X$.
Both first-order wave systems have the symmetric matrices
\[
 A_\tau=\begin{pmatrix}\theta_t&-\theta_x\\-\theta_x&\theta_t\end{pmatrix}>0,
 \qquad A_\xi=\begin{pmatrix}a_t&-a_x\\-a_x&a_t\end{pmatrix}.
\]
Their sources are respectively
$(\kappa e^{2\rho}J_\Phi-e^{2\rho}/R^2,0)$ and
$(-2e^{2\rho}\Phi/R^2-\kappa e^{2\rho}T,0)$.
No second embedding derivative enters this geometric principal block.
The first-order classical energy in an $H^r$ ball contracts in
$H^{r-1}$; \eqref{eq:finite-spatial-5} contracts the state in $D(\mathcal A^2)$ using only a
uniform $D(\mathcal A^3)$ bound. Both difference estimates have a
factor $T$. Time derivatives of principal quadratic moments are at
most quartic, so there is no higher-graph difference cycle. Picard
iteration, followed by closedness and finite smooth-data bootstrap,
gives the full unprojected EL solution. Smooth preparations with the
additional finite set of graph moments give six time derivatives and
the two coupling derivatives used below.
Unrestricted $X$, including its time equation, gives
$\nabla_\mu T^{\mu\nu}=-J_\Phi\nabla^\nu\Phi$; consequently the JT
constraints propagate. The sources include the variation of the
regulator and proper-number probe. At nonzero $\gamma$, retain
$J_\Phi=h_sn_S\langle Z_S\rangle/\ell_\perp$ and the resulting
$R(g)+2/R^2=2\kappa J_\Phi$.

At $\gamma=0$, $Z_S$ is conserved with mean zero; the source evolution
is independent of $h_s$, and $\langle Y_S\rangle=\sin(2h_s\Theta)$,
where $\Theta=\int n_SN\Phi\,\dd\tau\dd\xi$.
At the parity-fixed center, $X=(T(\tau),0)$ and $N=T'$.
Composing with each history's endpoint $\tau_\pm(t)=\theta_\pm(t,0)$
therefore gives exactly $\Theta_\pm(t)=\int_0^t\Phi_\pm(v,0)\dd v$,
with $t$ central proper time. The wave equation with common initial
data yields
\begin{equation}
 \Delta\Phi(t,0)=-\kappa d_{2,\rm total}t^4/12+O(t^5),
 \quad [t^5]\partial_{h_s}\Delta\langle Y_S\rangle\big|_0
       =-\kappa d_{2,\rm total}/30\ne0.                     \label{eq:finite-spatial-6}
\end{equation}
Finite-width proper-time phases converge in this fifth jet to \eqref{eq:finite-spatial-6}.
Scaling $M,J$ together keeps background-dilaton times recoil jets
bounded; the center endpoint contribution cancels exactly. Choose a
small but nonzero symmetric width preserving the coefficient, then
finite $L$ preserving all margins. Lower-order recoil jets are allowed.
Thus some finite short $t>0$ and small nonzero $h_s$ give nonzero
next-program contrast. Smooth parameter dependence preserves it for
sufficiently small finite $\gamma\ne0$, when the retained program
enters the actual material square \eqref{eq:finite-spatial-2}.

One can also check a complete mean-force response, without attributing
every term to geometry. Put $w_2=1$ on $\operatorname{supp}e$ and write
$\bar\Phi_0=\int n_SN\Phi(0)>0$, $b=\beta$, $c=\langle q_c\rangle/\sqrt I$,
$y=\langle q_u\rangle$, $a_\pm=\lambda\pm g$.
At the static-support large-$L$ limit the unperturbed means satisfy
$\ddot c=-\nu_c^2c-a_\pm^2bc/(I\omega^2)-a_\pm by/I$,
$\ddot y=-\nu_u^2y-a_\pm bc$.
For $\delta=\partial_\gamma\partial_{h_s}|_0$, independence at $\gamma=0$
and $\partial_{h_s}\langle Y_S\rangle=2\bar\Phi_0t+O(t^3)$ give
\begin{align*}
 \delta\ddot c={}&-(\nu_c^2+a_\pm^2b/(I\omega^2))\delta c
 -a_\pm b\delta y/I-2\bar\Phi_0t(by/I+2a_\pm bc/(I\omega^2)),\\
 \delta\ddot y={}&-\nu_u^2\delta y-a_\pm b\delta c
                       -2b\bar\Phi_0t c .
\end{align*}
The common $O(t^3)$ sensor term cancels in the branch fourth-jet.
With $\dot c(0)=\Omega$, all other initial means zero, these equations
include the clock/state variation and give
\begin{equation}
 [t^4]\delta(\ddot y_+-\ddot y_-)
 =\bar\Phi_0\left[\frac{8\lambda g b^2\Omega}{3I\omega^2}
                         +O_I(L^{-1})\right]\ne0 .           \label{eq:finite-spatial-7}
\end{equation}
The finite-$L$ correction includes proper-time pullback: mixed
synchronization displacement can start at $t^4$, spatial displacement
at $t^5$, while mixed curvature starts at $t^5$.
Their fourth-jet mean-force contributions, divided by $\bar\Phi_0$,
are $O_I(L^{-1})$ by the scaled positive EL inverse. Initial mixed
source differences vanish. Ordinary phase terms are present: \eqref{eq:finite-spatial-7} is
not a purely geometry-induced coefficient. Geometry's nonzero path
was separately certified by \eqref{eq:finite-spatial-6}.
Choose finite $L$ also preserving the strict sign in \eqref{eq:finite-spatial-7}, then small
finite $h_s,\gamma$ and a short nonzero event. All parameters and the
positive operating window are finite; no external read/reinstall is
performed.
\end{proof}

\subsection{Covariant bulk evolution of a retained quartet}
\label{app:covariant-quartet-bulk}

The full-mode carrier requires a covariant source, not just a conserved
finite-mode Hamiltonian. Use a physical metric $g=e^{2\rho}(-dt^2+dx^2)$
and an independent unit material director
$U=e^{-\rho}(\cosh\eta,\sinh\eta)$, with orthogonal unit $V$.
The work field $\chi$ and director have actions
\begin{equation}
 S_\chi=\frac{I_\chi}{2}\int\sqrt{-g}
       [(U\nabla\chi)^2-C_\chi(V\nabla\chi)^2],\qquad
 S_\eta=\frac{I_\eta}{2}\int\sqrt{-g}
       [(U\Omega)^2-C_\eta(V\Omega)^2],                     \label{eq:bulk-paid-fields}
\end{equation}
where $\Omega=d\eta$ plus the physical spin connection and $0<C_i<1$.
A positive elastic bar with coordinate $a$ anchors this director through
$-\Lambda\int\sqrt{-g}(U\nabla a)^2/2$.
Prepare a strict bar domain $B-\Lambda>K_a>0$ and finite alignment
overcoming the work-clock orientation curvature
$I_\chi(C_\chi-1)$ and the quantum pressure debit. These are paid
material parameters. Deriving $U$ from $\nabla a$ instead would put
second bar derivatives into the optical coefficients and lose the
same-order estimate below.

The quartet comprises two nonchiral Majorana sectors of central charge
$c_s=1/2$, and all three spectator bosons of $c_s=1$. Each sector has
both chiralities and $\gamma_s=c_s/(24\pi)$; the total central charge
is four. Its actual velocity $0<v_s(\chi)<1$ defines
$\bar g_s=-v_s^2U^\flat U^\flat+V^\flat V^\flat$.
Supply the original global representation, density operator and every
old reference. Choose optical-metric covariant point splitting with no
additional finite two-derivative medium counterterm. The anomaly
normalization is the conventional one~\cite[Section~2]{Codello2010};
the actual causal state, rather than a vacuum inverse-curvature
functional, determines its homogeneous stress.

\begin{proposition}[Actual-state covariant quartet bulk]
\label{prop:covariant-quartet-bulk}
Supply a consistent global full-mode representation and a smooth actual
initial state with the required energy moments and old references.
Prepare smooth local clock/bar/director data with strict paid margins,
the causal Hessian conditions \eqref{eq:bulk-cone-domain} and the JT
constraints \eqref{eq:bulk-JT-constraints}. The paid fields, retained
quartet state and renormalized constrained JT equations then have a
smooth common local-time bulk solution. Its mean solution and source
are independent of admissible exterior data continuations inside the
physical domain of dependence. Equality of its quantum restrictions
additionally requires equality of the initial local joint-state
restrictions, not merely stress means. This result does not prescribe
arbitrary local quantum means or a physical moving boundary.
\end{proposition}

\begin{proof}
For one sector, set
\begin{equation}
 d=\cosh^2\eta-v^2\sinh^2\eta,\quad
 N=v/d,\quad \beta=(v^2-1)\sinh\eta\cosh\eta/d,\quad
 \omega=\rho+\tfrac12\log d.                               \label{eq:bulk-optical-ADM}
\end{equation}
Thus $\bar g=e^{2\omega}\hat g$ with
$\hat g=-N^2dt^2+(dx+\beta dt)^2$. Its canonical chiral coefficients
are $h_\sigma=N-\sigma\beta>0$. Write $D_t=\partial_t-\beta\partial_x$,
$k=\beta_x/N$, $E=T_++T_-$ and $P=T_+-T_-$. The actual Virasoro
evolution gives
\begin{align}
 D_tE+NP_x+2N_xP-2\beta_xE&=-2\gamma\beta_{xxx},\\
 D_tP+NE_x+2N_xE-2\beta_xP&= 2\gamma N_{xxx}.               \label{eq:bulk-canonical-Ward}
\end{align}
The exact conversion to the orthonormal $\hat g$ frame is
\begin{equation}
 \hat E=E-\gamma k^2,\quad \hat P=P+2\gamma k_x,\quad
 \hat S=\hat E+\gamma\hat R,\qquad
 \hat R=-2(N_{xx}+D_tk-Nk^2)/N.                           \label{eq:bulk-projective-conversion}
\end{equation}
Substitution gives both covariant conservation equations with zero
residual and trace $\gamma\hat R$. The lapse and shift variations of
$-\gamma\int\beta_x^2/N$ give precisely this projective conversion.
The Weyl cocycle $\gamma\int\sqrt{-\hat g}[(\hat\nabla\omega)^2+
\omega\hat R]$ then yields
\begin{equation}
 \bar T_{ab}=\hat T_{ab}+2\gamma\left[
 \hat\nabla_a\hat\nabla_b\omega-\omega_a\omega_b
 -\hat g_{ab}\hat\Box\omega+
 \tfrac12\hat g_{ab}(\hat\nabla\omega)^2\right].           \label{eq:bulk-Weyl-conversion}
\end{equation}
Its trace is $\gamma\bar R$ and its $\bar g$ divergence vanishes.
In the bulk, the combined conversion is the first-jet density
\begin{equation}
 \mathcal L_{{\rm loc},s}=\gamma_s\left[
 -\frac{(D_t\omega_s-\beta_{s,x})^2}{N_s}
 +N_s\omega_{s,x}^2+2\omega_{s,x}N_{s,x}\right].            \label{eq:bulk-first-jet}
\end{equation}
The integrated-by-parts difference is
$-2\gamma_s[\partial_x(\omega_sN_{s,x})+
\partial_t(\omega_sk_s)-\partial_x(\omega_s\beta_sk_s)]$.
It cannot be dropped at a physical wall.

Define the canonical closed-time-path functional by
$Z_{\rm can}[z_+,z_-]=\tr(U[z_+]\varrho_0U[z_-]^\dagger)$.
Add $\Gamma_{\rm loc}[z_+]-\Gamma_{\rm loc}[z_-]$, with
$\Gamma_{\rm loc}=\sum_s\int\mathcal L_{{\rm loc},s}$.
Variation at coincident branches gives the actual canonical expectation
and the local forces in \eqref{eq:bulk-first-jet}. No new homogeneous
Polyakov field, independent purifier or in-out state is introduced.
Positive $h_\sigma$ makes each canonical Hamiltonian lower bounded,
not nonnegative: diffeomorphism conjugation includes the debit
$-\gamma_s\int h_{\sigma,x}^2/h_\sigma\,/2$ and the chosen Casimir term.
On the prepared smooth coefficient ball these have a finite uniform
bound. One common paid scalar offset covers it; its stress is included
in both histories.

An exact local reduction supplies the same-order estimate. For each
$u_j=\sigma_jh_j$, solve
\begin{align}
 w_{j,t}+u_jw_{j,x}+u_{j,x}w_j&=-u_{j,xx},
 &w_j(0,x)&=0,\notag\\
 A_{j,t}+u_jA_{j,x}+2u_{j,x}A_j&=0,
 &A_j(0,x)&=T_j(0,x).                                    \label{eq:bulk-local-transport}
\end{align}
Differentiation gives
$(\partial_t+u_j\partial_x+2u_{j,x})(w_{j,x}-w_j^2/2)
=-u_{j,xxx}$. Hence, with $Z_j=\gamma_jw_j$,
\begin{equation}
 T_j=A_j-Z_{j,x}+Z_j^2/(2\gamma_j).                       \label{eq:bulk-local-stress}
\end{equation}
This is the exact Ward solution. If
$f_{j,t}+u_jf_{j,x}=0$, $f_j(0,x)=x$, then
$w_j=\partial_x\log f_{j,x}$,
$A_j=f_{j,x}^2T_j(0,f_j)$ and
$T_j=f_{j,x}^2T_j(0,f_j)-\gamma_j\mathcal S(f_j)$.
These are the actual characteristic connection and transported initial
stress, not additional quantum fields or a new source action.

Put $q=(\chi,\eta)$, $P_q=q_t$, $Q_q=q_x$.
Let $K_{tt},K_{tx},K_{xx}$ be the $q$ Hessians of
the complete first-jet classical action plus \eqref{eq:bulk-first-jet},
with $K_{tx}=\mathcal L_{P_qQ_q}+\mathcal L_{Q_qP_q}$, and let
$G$ have columns $g_j=\nabla_qh_j$. The matter principal matrix and
its symmetrizer are
\begin{equation}
 A_m=\begin{pmatrix}
 K_{tt}^{-1}K_{tx}&K_{tt}^{-1}K_{xx}&-K_{tt}^{-1}G\\
 -1&0&0\\0&\operatorname{diag}(\gamma_j\sigma_j)G^T&
                 \operatorname{diag}(\sigma_jh_j)
 \end{pmatrix},\quad
 H_m=\begin{pmatrix}K_{tt}&0&0\\
 0&-K_{xx}&G\\0&G^T&\operatorname{diag}(h_j/\gamma_j)
 \end{pmatrix}.                                          \label{eq:bulk-symmetrizer}
\end{equation}
Multiplication gives $H_mA_m=(H_mA_m)^T$. The Schur conditions for
$H_m$ and $H_m\pm H_mA_m$ are exactly
\begin{align}
 K_{tt}&>0,\qquad
 -K_{xx}-\sum_j\gamma_jg_jg_j^T/h_j>0,\notag\\
 K_{tt}+K_{xx}\pm K_{tx}
       \pm\sum_j\frac{\gamma_j\sigma_jg_jg_j^T}
                              {1\pm\sigma_jh_j}&>0.       \label{eq:bulk-cone-domain}
\end{align}
At rest the director stiffness is corrected by
$\Delta D=\sum_s2\gamma_s(1-v_s^2)^2/v_s$.
Finite large inertias with strict $0<C_i<1$ give a nonempty domain;
\eqref{eq:bulk-cone-domain} uses its actual nonlinear coefficients,
including nonzero director boost. These inequalities imply a uniform
matter cone strictly inside the physical light cone.

The metric and dilaton each have light speeds $\lambda=\pm1$.
There is no resonant trace coupling between them. For the full
$(\chi,\eta,\rho)$ Hessians of a quantum sector, direct substitution
of \eqref{eq:bulk-optical-ADM} gives the exact identity
\begin{equation}
 (K^3_{tt,s}-\lambda K^3_{tx,s}+K^3_{xx,s})
       \begin{pmatrix}0\\\lambda\\1\end{pmatrix}
 =\sum_{\sigma=\pm1}
 \frac{\gamma_s\sigma\lambda h_{\sigma,s,\eta}}
      {\lambda-\sigma h_{\sigma,s}}
       \nabla_{(\chi,\eta,\rho)}h_{\sigma,s}.              \label{eq:bulk-light-residue}
\end{equation}
The right side has zero $\rho$ component. The director connection
also vanishes on this light mode; its clock component is zero.
The corresponding stress is
$Z_{\sigma,s}=\gamma_s\sigma h_{\sigma,s,\eta}\lambda p/
(\lambda-\sigma h_{\sigma,s})$, with
$P_\eta=-p$, $Q_\eta=\lambda p$ and $P_\chi=Q_\chi=0$.
Hence the complete trace residue is zero. The strict matter cone
solves the matter-to-light Sylvester equations; the zero resonant
projection solves the remaining metric-to-dilaton equations.
Pulling back the block positive form produces a smooth full symmetrizer.
The ten $A_j$ variables append the diagonal principal block
$\operatorname{diag}(u_j)$ and positive form $1_{10}$; they introduce
no new light resonance.

For $m>5/2$, use $q\in H^{m+1}$ and $P_q,Q_q,Z,A\in H^m$;
\eqref{eq:bulk-local-stress} lies in $H^{m-1}$. The exact local split is
\begin{align}
 Z_{j,t}+u_jZ_{j,x}+\gamma_j\sigma_jg_j^TQ_{q,x}
 &=-\sigma_j(g_j^TQ_q)Z_j
   -\gamma_j\sigma_jh_{j,qq}[Q_q,Q_q],\notag\\
 A_{j,t}+u_jA_{j,x}&=-2\sigma_j(g_j^TQ_q)A_j.             \label{eq:bulk-local-remainders}
\end{align}
The force is $\sum_jg_jZ_{j,x}
-\sum_jg_j[A_j+Z_j^2/(2\gamma_j)]$. All remainders are smooth
local products in the same $H^m$ algebra. The conversion contributes
first-jet Hessian principal terms and quadratic first-jet remainders.
No spatial inverse, omitted mean or higher derivative of the quantum
state enters the iteration. Symmetric-hyperbolic iteration
gives $d\mathcal E_m/dt\le C_R(1+\mathcal E_m)$ on the declared
operating ball. Tame higher-order estimates keep smooth data smooth
on that same local time. The actual full-mode propagator therefore
exists for the solved coefficients; uniqueness of the distributional
Ward transport identifies its stress with the solved means on the
interior patch, with the same initial state and all references retained.
A norm-controlled exterior extension of the reduced initial variables
can be chosen for the local iteration. It does not assign artificial
means to that quantum state outside the initial interval, and no
common-action solution is inferred there. Its bound fixes the local
existence time; arbitrary exterior continuations are compared only on
their common interval of existence.

This local reduction also proves physical bulk causality. Let two
admissible continuations agree initially on $(a,b)$, including scalar
values and actual stress means. Their local difference energy on
$(a+t,b-t)$ has endpoint fluxes
$-V^T(H-HA)V$ on the left and $-V^T(H+HA)V$ on the right.
They are nonpositive; the interior is bounded by its energy. Gronwall
therefore gives identical mean solutions until the interval closes or
the operating domain is left. The interval boundaries are the null
curves of $g$. For quantum restrictions require the same initial joint
local state on the causal interval and the compared old references.
The actual sector diffeomorphisms pull local observables into that
interval, giving the corresponding restriction equality. Stress means
alone do not give quantum-state equality, nor is a Type-III local
density matrix assumed.

Finally write the minimally coupled JT action as
$(2\kappa)^{-1}\int\sqrt{-g}\Phi(R-2K_{\rm geo})$.
Its metric curvature equation is
$\rho_{tt}-\rho_{xx}=K_{\rm geo}e^{2\rho}$.
Set $w=\Phi_x$, $\pi=\Phi_t$. The initial constraints are
\begin{equation}
 \Phi_x=w,\qquad
 w_x=\rho_xw+\rho_t\pi-K_{\rm geo}e^{2\rho}\Phi-\kappa T_{00},
 \qquad \pi_x=\rho_x\pi+\rho_tw-\kappa T_{0x}.             \label{eq:bulk-JT-constraints}
\end{equation}
Here $T$ is the complete source just constructed, including initial
conversion jets; needed accelerations come from its invertible time
matrix. This spatial ODE has a local solution for three integration
constants, with positive $\Phi$ on a sufficiently small interval.
Conservation propagates its two traceless constraints locally along
physical light directions. Arbitrary periodic positive-energy JT data
are not assumed. A consistent global representation and state realize
the quantum solution; the causal interior restriction no longer depends
on an arbitrary exterior coefficient continuation.
\end{proof}

This closes the quartet's nonlinear bulk source--clock--director--JT
interface. It supplies neither moving-wall boundary terms nor the
other gate windows' errors. The original mixed target and its purifier
remain in the actual state, but clock and geometry changes alter gate
areas. Their same-history finite-horizon comparison and the physical
exhaust still have to be incorporated before the complete thermal
apparatus is obtained.

\subsection{A physical wall retains its thermal friction}
\label{app:Dirac-wall-response}

A finite rest mass cannot repair a perfect mirror by ignoring its
high-frequency reaction. Partial transmission has a causal passive
linear-response formulation~\cite{JaekelReynaud1992}. For the four
actual Dirac flavours, the corresponding force can be derived from a
smooth material mass layer rather than imported from a scalar mirror.
In a static material frame let
$\mathsf h_z=-iv\sigma_z\partial_x+M(x-z)\sigma_x$, with
$M\in C_c^\infty$, $M\ge0$, and positive bare mechanical mass $m_b$.
The same mass action gives the translation force
$F=\int M'(x-z)\psi^\dagger\sigma_x\psi\,dx+F_{\rm ct}$,
distributed bar reaction and metric mass stress. Use one fixed local
counterterm action, an actual static interacting Hadamard ground or
KMS preparation, and smooth flat-start displacement tests.

\begin{proposition}[Passive Dirac-layer response]
\label{prop:Dirac-wall-passivity}
For the static ground-state linear response, the full-continuum induced
vacuum inertia is finite. With positive bare mass and total positive
elastic stiffness, the unreduced causal mechanical impedance is
positive-real. In the static zero-charge KMS continuum the layer instead
has positive ohmic friction. Its vacuum-style low-frequency inertia
integral diverges, while the unreduced thermal impedance remains
convergent and positive-real. These are linear statements about the
specified passive preparation, not arbitrary driven quartet states.
\end{proposition}

\begin{proof}
Let $u_{E,a}$, $a=L,R$, be energy-normalized scattering modes with
incoming amplitude $(2\pi v)^{-1/2}$. Define
$A=\int M/v$, $Q=\|M'\|_1$ and $B=e^A/\sqrt{\pi v}$.
The exact transfer equation
$u_x=(iE\sigma_z+M\sigma_y)u/v$ bounds $\|u\|\le B$.
Integration by parts gives
\begin{align}
 f_{ab}(E,E'):=\int M'u_{E,a}^\dagger\sigma_xu_{E',b}
 &=-(E+E')\int(M/v)u_{E,a}^\dagger\sigma_yu_{E',b},\notag\\
 |f_{ab}(E,E')|&\le B^2\min\{Q,A|E+E'|\}.                \label{eq:Dirac-force-identity}
\end{align}
The mass anticommutator vanishes and $M$ removes the endpoint term.
For $N=4$ flavours the filled-negative-state force spectrum obeys
\begin{align}
 \chi_I(\nu)=\pi N\sum_{a,b}\int_0^\nu
                 |f_{ab}(E,E-\nu)|^2\,dE
 &\le4\pi NB^4\min\{Q^2\nu,A^2\nu^3/3\},\notag\\
 \mu=\frac2\pi\int_0^\infty\frac{\chi_I(\nu)}{\nu^3}\,d\nu
 &\le\frac{16NB^4AQ}{\sqrt3}.                             \label{eq:Dirac-vacuum-mass}
\end{align}
At zero energy the exact transfer gives reflection $\tanh^2A$,
and hence $\chi_I(\nu)=N\tanh^2A\,\nu^3/(6\pi v^2)+O(\nu^4)$.
Retain the complete causal memory. For $\Re s>0$ its impedance is
\begin{equation}
 Z(s)=m_bs+k/s+\frac2\pi s\int_0^\infty
              \frac{\chi_I(\nu)}{\nu(\nu^2+s^2)}\,d\nu.  \label{eq:Dirac-passive-impedance}
\end{equation}
The measure is positive and
$s/(s^2+\nu^2)=\tfrac12[(s-i\nu)^{-1}+(s+i\nu)^{-1}]$
has positive real part. Thus $\Re Z(s)>0$; the response has no
right-half-plane pole. The observed vacuum inertial mass is
$m_{\rm phys}=m_b+\mu$, not the action's bare $m_b$.
For two separated layers the same supportwise identity gives a
positive matrix spectral measure. Adding their actual finite static
polarization/Casimir Hessian to the paid elastic stiffness extends
this argument when that total stiffness and the bare mass matrix are
positive.

In the KMS continuum, with Fermi function $n_T(E)$,
\begin{equation}
 \chi_{I,T}(\nu)=\pi N\sum_{a,b}\int_{\mathbb R}
 [n_T(E)-n_T(E+\nu)]|f_{ab}(E+\nu,E)|^2\,dE.             \label{eq:Dirac-KMS-force}
\end{equation}
Its weights are positive. On shell,
$\sum_{a,b}|f_{ab}(E,E)|^2=2E^2R(E)/(\pi^2v^2)$,
where $R(E)$ is the actual reflection probability. Therefore
\begin{equation}
 \chi_{I,T}(\nu)=\gamma_T\nu+o(\nu),\qquad
 \gamma_T=\frac{2N}{\pi v^2}\int E^2R(E)[-n_T'(E)]\,dE>0.
                                                               \label{eq:Dirac-thermal-friction}
\end{equation}
The vacuum $\mu$ integral now diverges in the infrared.
Equation~\eqref{eq:Dirac-passive-impedance} itself stays finite:
its infrared integrand uses this ohmic law, while
$\int[n_T(E)-n_T(E+\nu)]dE=\nu$ and
\eqref{eq:Dirac-force-identity} give an ultraviolet linear bound.
Its positive-real proof is unchanged. A finite prepared exterior has
the corresponding discrete positive Kubo measure; this continuum
friction limit is not an uncharged infinite thermal reservoir.
\end{proof}

All transmitted exterior channels, old purifiers, static polarization
and motion work remain in this construction. Smooth spatial profiles
alone do not justify applying an instantaneous force to the old free
sea as a finite-energy vector; the valid response tests are also smooth
in time. More importantly, the actual driven body need not be passive.
Its opposite-chirality mass mixing and non-Gaussian quartet require a
joint local-source and task-error estimate. The common-velocity Ward
closure inside a massless gate does not close that boundary hierarchy.
The passive wall result establishes a physical linear margin, without
substituting that margin for the retained mixed-target reset proof.

\begin{thebibliography}{99}

\bibitem{JaekelReynaud1992}
M.-T.~Jaekel and S.~Reynaud, \emph{Causality, stability and passivity
for a mirror in vacuum}, Physics Letters A \textbf{167}, 227--232
(1992), doi:10.1016/0375-9601(92)90195-R,
arXiv:quant-ph/0101080.

\bibitem{Codello2010}
A.~Codello, \emph{Polyakov Effective Action from Functional
Renormalization Group Equation}, Annals of Physics \textbf{325},
1727--1738 (2010), doi:10.1016/j.aop.2010.04.013,
arXiv:1004.2171.

\bibitem{Vassilevich2003}
D.~V.~Vassilevich, \emph{Heat kernel expansion: user's manual},
Physics Reports \textbf{388}, 279--360 (2003),
doi:10.1016/j.physrep.2003.09.002, arXiv:hep-th/0306138.

\bibitem{Friedan1985}
D.~Friedan, \emph{Nonlinear Models in $2+\epsilon$ Dimensions},
Annals of Physics \textbf{163}, 318--419 (1985),
doi:10.1016/0003-4916(85)90384-7.

\bibitem{BernardDoyonViti2015}
D.~Bernard, B.~Doyon and J.~Viti,
\emph{Non-Equilibrium Conformal Field Theories with Impurities},
J.~Phys.~A \textbf{48}, 05FT01 (2015), arXiv:1411.0470,
doi:10.1088/1751-8113/48/5/05FT01.

\bibitem{Shen2024}
J.~Shen, C.~Peng and L.-X.~Li, \emph{Multiway Junction Conditions:
Booklets and Webs}, Phys.\ Rev.\ D \textbf{110}, 066021 (2024),
arXiv:2402.00694.

\bibitem{Hollands2019}
S.~Hollands, \emph{Relative entropy for coherent states in chiral CFT},
arXiv:1903.07508, especially Section~2.

\bibitem{Landulfo2016}
A.~G.~S.~Landulfo, \emph{Nonperturbative approach to relativistic quantum communication channels},
Phys.\ Rev.\ D \textbf{93}, 104019 (2016), arXiv:1603.06641.

\bibitem{Camassa2012KMS}
P.~Camassa, R.~Longo, Y.~Tanimoto and M.~Weiner,
\emph{Thermal States in Conformal QFT. I}, Commun.\ Math.\ Phys.
\textbf{309}, 703--735 (2012), arXiv:1101.2865.

\bibitem{Pendry1983}
J.~B.~Pendry, \emph{Quantum limits to the flow of information and entropy},
J.\ Phys.\ A: Math.\ Gen.\ \textbf{16}, 2161--2171 (1983),
doi:10.1088/0305-4470/16/10/012.

\bibitem{BernardDoyon2012}
D.~Bernard and B.~Doyon, \emph{Energy flow in non-equilibrium conformal field theory},
J.\ Phys.\ A: Math.\ Theor.\ \textbf{45}, 362001 (2012), arXiv:1202.0239.

\bibitem{BernardDoyon2015}
D.~Bernard and B.~Doyon, \emph{Non-equilibrium steady-states in conformal field theory},
Ann.\ Henri Poincar\'e \textbf{16}, 113--161 (2015), arXiv:1302.3125,
doi:10.1007/s00023-014-0314-8.

\bibitem{ReebWolf2014}
D.~Reeb and M.~M.~Wolf, \emph{An improved Landauer principle with finite-size corrections}, New J.\ Phys.\ \textbf{16}, 103011 (2014), arXiv:1306.4352.

\bibitem{QRC}
S.~Liu, \emph{Retained Quantum History, Physical Sources and Receiver Renewal in Q-RAIF}, (2026), series record DOI: 10.5281/zenodo.18528934.

\bibitem{ArakiZsido2005}
H.~Araki and L.~Zsid\'o, \emph{Extension of the structure theorem of Borchers and its application to half-sided modular inclusions}, Rev.\ Math.\ Phys.\ \textbf{17}, 491--543 (2005), arXiv:math/0412061.

\bibitem{Kiefer2012}
C. Kiefer, \emph{Quantum Gravity}, 3rd ed., Oxford University Press (2012).

\bibitem{Zanardi2001}
P. Zanardi, \emph{Virtual Quantum Subsystems}, Phys. Rev. Lett. \textbf{87}, 077901 (2001).

\bibitem{Zanardi2004}
P. Zanardi, D. A. Lidar, and S. Lloyd, \emph{Quantum Tensor Product Structures are Observable Induced}, Phys. Rev. Lett. \textbf{92}, 060402 (2004).

\bibitem{Liu2026PaperA}
S.~Liu, \emph{Physical Access and Correlation Geometry: The Holographic Alaya--Field Framework}, (2026), series record DOI: 10.5281/zenodo.18361706.

\bibitem{Liu2026PaperB}
S.~Liu, \emph{Coupling, Records, and Stable Relations: Formation, Protected Content and the Lifetime of a Geometry}, (2026), series record DOI: 10.5281/zenodo.18367060.

\bibitem{Liu2026QRAIFB}
S.~Liu, \emph{Thermodynamic Stability Constraints on Persistent Open Quantum Subsystems: Control Algebra, Physical Actuation and Energy Cost}, (2026), series record DOI: 10.5281/zenodo.18525890.

\bibitem{Liu2026Chain}
S.~Liu, \emph{The Imprint Chain: Records, Self-Calibration, and Physical Feedback}, (2026), series record DOI: 10.5281/zenodo.20151594.

\bibitem{Liu2026QRAIF_A}
S.~Liu, \emph{Lorentzian Read-out from Local Unitary Propagation: Clifford Relations, Causal Cones and Spectral Mass}, (2026), series record DOI: 10.5281/zenodo.18525876.

\bibitem{PaperCII}
S.~Liu, \emph{Higher-Dimensional Consistency of Geometry and Modular Flow: Physical Source Identification and a Second-Order Extension} (Consciousness~II), (2026), series record DOI: 10.5281/zenodo.19607543.

\bibitem{PaperCI}
S.~Liu, \emph{Geometry Relative to a Cut: Selection, Generation and Modular Inseparability}, (2026), series record DOI: 10.5281/zenodo.19607387.

\bibitem{QRNV}
S.~Liu, \emph{No View from Nowhere: Read-out Closure and the Vacuum Curvature Test in Q-RAIF A}, Zenodo (2026), DOI: 10.5281/zenodo.20284454.

\bibitem{RyuTakayanagi2006}
S. Ryu and T. Takayanagi, \emph{Holographic Derivation of Entanglement Entropy from AdS/CFT}, Phys. Rev. Lett. \textbf{96}, 181602 (2006).

\bibitem{VanRaamsdonk2010}
M. Van Raamsdonk, \emph{Building up spacetime with quantum entanglement}, Gen. Relativ. Gravit. \textbf{42}, 2323 (2010).

\bibitem{Rovelli2004}
C. Rovelli, \emph{Quantum Gravity}, Cambridge University Press (2004).

\bibitem{Smolin2006}
L. Smolin, \emph{The case for background independence}, in \emph{The Structural Foundations of Quantum Gravity}, eds. D. Rickles, S. French, J. Saatsi, Oxford University Press (2006).

\bibitem{Maldacena1999}
J. M. Maldacena, \emph{The Large N limit of superconformal field theories and supergravity}, Int. J. Theor. Phys. \textbf{38}, 1113 (1999).

\bibitem{Vidal2008}
G. Vidal, \emph{Class of quantum many-body states that can be efficiently simulated}, Phys. Rev. Lett. \textbf{101}, 110501 (2008).

\bibitem{Swingle2012}
B. Swingle, \emph{Entanglement Renormalization and Holography}, Phys. Rev. D \textbf{86}, 065007 (2012).

\bibitem{Jacobson1995}
T. Jacobson, \emph{Thermodynamics of Spacetime: The Einstein Equation of State}, Phys. Rev. Lett. \textbf{75}, 1260 (1995).

\bibitem{Thiemann2007}
T. Thiemann, \emph{Modern Canonical Quantum General Relativity}, Cambridge University Press (2007).

\bibitem{Petz1996}
D. Petz, \emph{Monotone metrics on matrix spaces}, Linear Algebra Appl. \textbf{244}, 81 (1996).

\bibitem{Bengtsson2006}
I. Bengtsson and K. \.Zyczkowski, \emph{Geometry of Quantum States: An Introduction to Quantum Entanglement}, Cambridge University Press (2006).

\bibitem{Lashkari2014}
N. Lashkari, M. B. McDermott, and M. Van Raamsdonk, \emph{Gravitational dynamics from entanglement ``thermodynamics''}, JHEP \textbf{04}, 195 (2014).

\bibitem{Faulkner2014}
T. Faulkner, M. Guica, T. Hartman, R. C. Myers, and M. Van Raamsdonk, \emph{Gravitation from Entanglement in Holographic CFTs}, JHEP \textbf{03}, 051 (2014).

\bibitem{Takesaki2003}
M.~Takesaki, \emph{Theory of Operator Algebras II}, Encyclopaedia of Mathematical Sciences \textbf{125}, Springer (2003).

\bibitem{BratteliRobinson1997}
O.~Bratteli and D.~W.~Robinson, \emph{Operator Algebras and Quantum Statistical Mechanics}, Vols.\ 1--2, 2nd ed., Springer (1997).

\bibitem{Haag1996}
R. Haag, \emph{Local Quantum Physics: Fields, Particles, Algebras}, Springer-Verlag (1996).

\bibitem{Yngvason2005}
J.~Yngvason, \emph{The role of type III factors in quantum field theory}, Rep.\ Math.\ Phys.\ \textbf{55}, 135 (2005), arXiv:math-ph/0411058.

\bibitem{Wiesbrock1993}
H.-W.~Wiesbrock, \emph{Half-sided modular inclusions of von Neumann algebras}, Commun.\ Math.\ Phys.\ \textbf{157}, 83 (1993).

\bibitem{BisognanoWichmann1975}
J.~J.~Bisognano and E.~H.~Wichmann, \emph{On the duality condition for a Hermitian scalar field}, J.\ Math.\ Phys.\ \textbf{16}, 985 (1975).

\bibitem{BisognanoWichmann1976}
J.~J.~Bisognano and E.~H.~Wichmann, \emph{On the duality condition for quantum fields}, J.\ Math.\ Phys.\ \textbf{17}, 303 (1976).

\bibitem{CasiniHuertaMyers2011}
H.~Casini, M.~Huerta, and R.~C.~Myers, \emph{Towards a derivation of holographic entanglement entropy}, JHEP \textbf{1105}, 036 (2011), arXiv:1102.0440.

\bibitem{JLMS2016}
D.~L.~Jafferis, A.~Lewkowycz, J.~Maldacena, and S.~J.~Suh, \emph{Relative entropy equals bulk relative entropy}, JHEP \textbf{1606}, 004 (2016), arXiv:1512.06431.

\bibitem{Jacobson2016}
T.~Jacobson, \emph{Entanglement equilibrium and the Einstein equation}, Phys.\ Rev.\ Lett.\ \textbf{116}, 201101 (2016), arXiv:1505.04753.

\bibitem{Connes1994}
A.~Connes, \emph{Noncommutative Geometry}, Academic Press (1994).

\bibitem{OhyaPetz1993}
M.~Ohya and D.~Petz, \emph{Quantum Entropy and Its Use}, Springer (1993).

\bibitem{AlmheiriPolchinski2015}
A.~Almheiri and J.~Polchinski, \emph{Models of AdS$_2$ backreaction and holography}, JHEP \textbf{1511}, 014 (2015), arXiv:1402.6334.

\bibitem{Moitra2020}
U.~Moitra, S.~K.~Sake, S.~P.~Trivedi and V.~Vishal, \emph{Jackiw--Teitelboim model coupled to conformal matter in the semi-classical limit}, JHEP \textbf{2020}(4), 199 (2020), arXiv:1908.08523.

\bibitem{FoundationI}
S.~Liu, \emph{Foundation I: From Source to Manifestation. Accessible Algebras from a Source, Its Noise and a Readable Seed}, (2026), series record DOI: 10.5281/zenodo.19627786.

\bibitem{Liu2026PaperG}
S.~Liu, \emph{Internal Selection and Revisable Organization: How History
Changes Permissions and Physical Sources}, (2026),
series record DOI: 10.5281/zenodo.18402907.

\bibitem{FoundationIII}
S.~Liu, \emph{Foundation III: The Perfuming Loop. Held Content,
Learned Frames and Physical Return}, (2026),
series record DOI: 10.5281/zenodo.19642189.

\bibitem{BrainAnchor}
S.~Liu, \emph{The Brain as a Plastic Observable-Algebra Network:
Learning Shared Access from Sensorimotor History}, 2026.

\bibitem{Callebaut2019}
N.~Callebaut, \emph{The gravitational dynamics of kinematic space}, JHEP \textbf{2019}(2), 153 (2019), arXiv:1808.10431.

\bibitem{Perche2025}
T.~R.~Perche, J.~P.~M.~Pitelli, and D.~A.~T.~Vanzella, \emph{The stress-energy tensor of an Unruh--DeWitt detector}, Phys.\ Rev.\ D \textbf{111}, 045004 (2025), arXiv:2411.09732.

\bibitem{Griguolo2026Cutoff}
L.~Griguolo, J.~Papalini, L.~Russo and D.~Seminara,
\emph{A new perspective on dilaton gravity at finite cutoff},
arXiv:2512.21774v2 (2026).

\bibitem{CastroGross2013}
A.~Castro and E.~K.~U.~Gross,
\emph{Optimal control theory for quantum-classical systems: Ehrenfest
Molecular Dynamics based on time-dependent density-functional theory},
arXiv:1308.4162 (2013), Section~2, Eqs.~(2)--(6).

\end{thebibliography}

\end{document}
